% Chapter generated from the companion chapter on the rank-two wild place, Draft~8 by booktools/assemble.py
\begingroup
% macros of the companion chapter on the rank-two wild place, Draft~8
\providecommand{\F}{}\renewcommand{\F}{F}
\providecommand{\Z}{}\renewcommand{\Z}{\mathbb Z}
\providecommand{\C}{}\renewcommand{\C}{\mathbb C}
\providecommand{\Q}{}\renewcommand{\Q}{\mathbb Q}
\providecommand{\Fq}{}\renewcommand{\Fq}{\mathbb F}
\providecommand{\Nm}{}\renewcommand{\Nm}{\operatorname N}
\providecommand{\Nrd}{}\renewcommand{\Nrd}{\operatorname{Nrd}}
\providecommand{\Trd}{}\renewcommand{\Trd}{\operatorname{Trd}}
\providecommand{\Tr}{}\renewcommand{\Tr}{\operatorname{Tr}}
\providecommand{\Irr}{}\renewcommand{\Irr}{\operatorname{Irr}}
\providecommand{\Ind}{}\renewcommand{\Ind}{\operatorname{Ind}}
\providecommand{\Ad}{}\renewcommand{\Ad}{\operatorname{Ad}}
\providecommand{\tH}{}\renewcommand{\tH}{\widetilde H}
\providecommand{\tZ}{}\renewcommand{\tZ}{\widetilde Z}
\providecommand{\tU}{}\renewcommand{\tU}{\widetilde U}
\providecommand{\tV}{}\renewcommand{\tV}{\widetilde V}
\providecommand{\tth}{}\renewcommand{\tth}{\widetilde\theta}
\providecommand{\Ocal}{}\renewcommand{\Ocal}{\mathcal O}
\providecommand{\Pfr}{}\renewcommand{\Pfr}{\mathfrak P}
\providecommand{\pfr}{}\renewcommand{\pfr}{\mathfrak p}
\providecommand{\Reg}{}\renewcommand{\Reg}{\operatorname{Reg}}
\providecommand{\Heis}{}\renewcommand{\Heis}{\mathrm{Heis}}
\providecommand{\Cubes}{}\renewcommand{\Cubes}{\mathrm{Cubes}}
\providecommand{\Good}{}\renewcommand{\Good}{\mathrm{Good}}
\providecommand{\Ptag}{}\renewcommand{\Ptag}{\textnormal{\textbf{(P)}}}
\providecommand{\Pctag}{}\renewcommand{\Pctag}{\textnormal{\textbf{(P$^{\mathrm c}$)}}}
\providecommand{\Etag}{}\renewcommand{\Etag}{\textnormal{\textbf{(E)}}}
\providecommand{\Otag}{}\renewcommand{\Otag}{\textnormal{\textbf{(O)}}}
\providecommand{\Ctag}{}\renewcommand{\Ctag}{\textnormal{\textbf{(C)}}}
\chapter[The wild place in degree two]{The wild place in degree two: toral positivity and the nonabelian gluing problem}\label{ch:W}
\paragraph{Overview.} We study irreducible-character integrality for the scalar-central
odd-degree metaplectic cover of a quaternion division algebra at a
wild place.  The general result is a uniform quadratic toral
permutation-character theorem: for every quadratic $E/F$ in odd
residue characteristic and every odd $n$ with $\mu_n\subset F$, the
discriminant ratio $\Delta(t)/\Delta(t^n)$ is a finite permutation
character of $E^\times/F^\times$.  The main new result concerns the
running example $F=\Q_3(\zeta_3)$, $n=3$, $D$ quaternion.  We prove
that the division-side transform $L_D$ of Chapter~\ref{ch:III} carries the
trivial character of $D^\times$ to the character of a single genuine
irreducible representation $\tau$ of the cover, of dimension $27$
(Target~A), and we describe $\tau$: its restriction to the
pro-$3$ part of the cover is induced from a canonical genuine
character of the covered third congruence subgroup.  The proof has
three layers: a reduction, via Brauer's characterization and the
toral theorem, to the pro-$3$ Sylow subgroup; a structural analysis
of the good elements of the congruence filtration which turns the
question into a finite combinatorial problem about $72$ vectors in
$\Fq_3^5$ that has exactly two mirror-image solutions; and an
explicit cocycle model of the cover, in which the one remaining
identity (multiplicativity of the canonical star section on the sixth
congruence subgroup) is proved by a short computation with cubic
Hilbert symbols.  The same holds for every irreducible representation of $D^\times$
trivial on the third congruence subgroup and compatible with the cover
(central character trivial on $\mu_3$): the transform of its
character is the character of a genuine irreducible representation
whose restriction to the pro-$3$ part is a sum of irreducible
representations induced from the covered third congruence subgroup,
one for each orbit of genuine constituents of a ``cube-root transform''
of $\rho|_{K^1}$ (part of Target~B).  For an arbitrary row
we prove the reduction to the pro-$3$ Sylow subgroup and a sufficient
criterion in terms of a cube-root transform of $\Theta_\rho$.  Finally,
the $p$-th power identity and the multiplicativity of the canonical star
section on the $2d$-th congruence subgroup are proved for every odd prime
$p$ and every $F\supset\mu_p$ (part of Target~C); what is not uniform is
identified.

\section{Purpose, scope, and status}\label{W:sec:intro}
This chapter separates the wild division-side character problem from the
conditional global trace comparison in Chapter~\ref{ch:IV}.
The chapter proves the general toral theorem (Section~\ref{W:sec:general-toral}),
works out the three quadratic tori of the running example
(Section~\ref{W:sec:examplewild}), proves the positivity--integrality
criterion (Section~\ref{W:sec:positive}), and formulates three targets:
\begin{description}[style=nextline]
\item[Target A] For $F=\Q_3(\zeta_3)$, $D$ quaternion and $n=3$, decide
whether $L_D\mathbf1$ is the character of a single $27$-dimensional
genuine irreducible representation of $\tH$ with the prescribed central
character.
\item[Target B] For every ordinary irreducible $\rho$ with a fixed
compatible central character, prove that all coefficients $a_{\rho\tau}$
are integral, or exhibit a nonintegral one.
\item[Target C] Extend Target~B to all odd $n$, all $F$ of odd residual
characteristic with $\mu_n\subset F$, and quaternion $D/F$.
\end{description}
The later sections settle Target~A affirmatively (Theorem~\ref{W:thm:main-targetA}),
settle Target~B for all representations trivial on $K^3$ and their norm
twists (Theorems~\ref{W:thm:levelone}, \ref{W:thm:leveltwo}), prove the reduction theorem for
every row (Theorem~\ref{W:thm:reduction-all}) and a sufficient criterion
for a row (Proposition~\ref{W:prop:criterion}), and prove the two
cocycle-level ingredients of Target~C uniformly in $p$
(Theorems~\ref{W:thm:pthpower}, \ref{W:thm:H1-gen}).  The dependencies of
the main theorem are listed in Remark~\ref{W:rem:dependencies}.

\paragraph{Labels.}
Proved here is the default and carries no visible tag; \Etag\ marks a
statement established elsewhere and used as a black box; \Otag\
marks one outlined but not fully written out or checked; \Ctag\ marks
one that is conjectural. A remark immediately follows every
\Etag, \Otag\ or \Ctag\ statement, saying where it is proved or, if
conjectural or outlined, the motivation or what remains. The label
\Pctag\ marks a statement whose proof uses, besides hand arguments, an
exact finite computation (linear algebra over $\Fq_3$, or exact
arithmetic in a number field) that is \emph{exhaustive} over an
explicitly described finite set, and is proved here in that sense, so
it also carries no further tag beyond \Pctag\ itself. Every such
computation is described in the text, and the scripts are listed in
Appendix~\ref{W:app:scripts}. No statement labeled \Pctag\ depends on
floating-point arithmetic or on sampling. Sampled computations occur
in the chapter only as corroboration of statements that are proved by
hand, and are identified as such where they are mentioned.

\begin{center}\small
\begin{tabular}{@{}p{0.60\linewidth}p{0.26\linewidth}p{0.08\linewidth}@{}}
\toprule
Statement & Source & Status\\\midrule
Unitary local $L_D$ and identity factor & Chapter~\ref{ch:III} & \Etag\\
Heisenberg commutator and degree on maximal tori & Chapter~\ref{ch:I} & \Etag\\
Quadratic toral permutation-character theorem & Theorem~\ref{W:thm:main} & \Ptag\\
Three exact cubic wild field-torus calculations & Section~\ref{W:sec:examplewild} & \Ptag\\
All-row positivity--integrality equivalence & Proposition~\ref{W:prop:positive} & \Ptag\\
Structure of good elements in the congruence filtration & Section~\ref{W:sec:depth} & \Ptag\\
Reduction of Target A to the pro-$3$ Sylow subgroup & Theorem~\ref{W:thm:reduction} & \Ptag\\
Finite model; the defect function; uniqueness of the design & Section~\ref{W:sec:finite} & \Pctag\\
Explicit cocycle model (Kubota--Kazhdan--Patterson) & Lemma~\ref{W:lem:model} & \Etag/\Ptag\\
Cube identity; star section on $K^6$; proof of (H1) & Section~\ref{W:sec:model} & \Ptag\\
$L_D\mathbf1$ is a genuine irreducible character (Target A) & Theorem~\ref{W:thm:main-targetA} & \Pctag\\
Reduction theorem for every row & Theorem~\ref{W:thm:reduction-all} & \Ptag\\
Cube-root transform; sufficient criterion for a row & Proposition~\ref{W:prop:criterion} & \Ptag\\
All rows trivial on $K^2$, and norm twists (part of Target B) & Theorem~\ref{W:thm:levelone} & \Pctag\\
All rows trivial on $K^3$ (part of Target B) & Theorem~\ref{W:thm:leveltwo} & \Pctag\\
Target B for level $\ge3$ & Sections~\ref{W:subsec:cuberoot}, \ref{W:sec:future} & \Otag\\
$p$-th power identity and (H1) for every odd $p$ (part of Target C) & Section~\ref{W:sec:targetC} & \Ptag\\
Target C: Brauer reduction; finite design & Remark~\ref{W:rem:uniform} & \Otag/\Ctag\\
\bottomrule
\end{tabular}
\end{center}

\section{Background}\label{W:sec:background}
This section collects standard facts used throughout, with references;
nothing in it is new.  Readers familiar with division algebras and with
the Heisenberg--Weil machinery can skip to Section~\ref{W:sec:transform}
and return for the heuristic of Section~\ref{W:subsec:howe}, which explains
how the results of Sections~\ref{W:sec:finite}--\ref{W:sec:plan} should be
read.

\subsection{Division algebras over local and global fields}\label{W:subsec:divalg}
The general theory recalled in this subsection is developed with
more examples, and tied to the running example throughout, in
\S\ref{gl:divalg}; what follows specializes it to this chapter's
running example and extends it with the congruence filtration and
the notion of level on which the rest of the chapter depends.

Let $F$ be a nonarchimedean local field.  A central simple algebra over
$F$ has dimension $r^2$ for an integer $r$ (its degree), and its class in
the Brauer group is determined by the Hasse invariant
$\mathrm{inv}:\mathrm{Br}(F)\xrightarrow{\sim}\Q/\Z$ \cite[Ch.~XIII]{Ser79}; \cite[\S31]{Rei75}.  The algebras of degree $r$ are the $M_m(D)$ with $D$ a division
algebra of degree $d\mid r$, $md=r$, and $D$ is determined by its
invariant $a/d\bmod1$ with $(a,d)=1$; in particular there is exactly one
quaternion division algebra over $F$, with invariant $\tfrac12$.  A
division algebra $D$ of degree $r$ has the following structure
\cite[\S\S12--14]{Rei75}; \cite[Ch.~I]{Wei73}.  It carries a discrete valuation $v_D$
extending $v_F$, with $v_D(F^\times)=r\Z$ when normalized by
$v_D(\Pi)=1$ for a uniformizer $\Pi$; its valuation ring $\Ocal_D$ is the
unique maximal order, with unique maximal two-sided ideal
$\Pfr=\Pi\Ocal_D$, and residue field $k_D=\Ocal_D/\Pfr$ the extension of
$k_F$ of degree $r$.  The unramified extension $E_0/F$ of degree $r$
embeds in $D$, $\Ocal_D=\Ocal_{E_0}[\Pi]$, and $D$ is the cyclic algebra
$\bigoplus_{k<r}E_0\Pi^k$ with $\Pi x\Pi^{-1}=\sigma^a(x)$ for a
generator $\sigma$ of $\mathrm{Gal}(E_0/F)$ and $\Pi^r=\varpi$ up to a
unit; for the quaternion algebra one may take $\Pi^2=\varpi$ and
$\Pi x\Pi^{-1}=\sigma(x)$, which is the description
$D=E_0\oplus E_0j$ used in Example~\ref{W:ex:ranktwo-data}.  Every field
extension of $F$ of degree $r$ embeds in $D$ as a maximal subfield, and
when $r$ is prime (the case of this book), every element of
$D\setminus F$ generates a maximal subfield; the
maximal subfields of the quaternion algebra are the three quadratic
extensions of $F$ ($p$ odd), one unramified and two ramified, and the
maximal tori of $D^\times$ are their multiplicative groups.  The reduced
norm $\Nrd:D^\times\to F^\times$ is surjective, $D^\times/F^\times$ is
compact, and $D^\times$ has the filtration by the normal subgroups
$K^m=1+\Pfr^m$ ($m\ge1$) with $K^m/K^{m+1}\simeq k_D$ (additive) and
$\Ocal_D^\times/K^1\simeq k_D^\times$; the group $D^\times/F^\times K^1$
is finite, of order $r(q^r-1)/(q-1)$, and $K^1$ is a pro-$p$ group.  For
the running example ($F=\Q_3(\zeta_3)$, $q=3$, $r=2$) these are the
groups $Q=K^1F^\times/F^\times$, $Q_m$ and $D^\times/F^\times K^1\simeq D_8$
of Section~\ref{W:sec:depth}.  Every irreducible smooth representation of
$D^\times$ is finite dimensional and trivial on some $K^m$; its
\emph{level} is the least $m$ with $\rho|_{K^{m+1}}$ trivial. (The
general definition of level and the congruence filtration, and the
running example's own level-$\le2$ representation count, are given
once, for reuse book-wide, in the Congruence filtration and level
entry of \S\ref{gl:divalg}; the rest of this paragraph specializes
that entry to the combinatorics this chapter actually needs.)

\paragraph{Levels.}  The level organizes everything below, so we
record what it means.  Since $K^{m+1}$ is normal and
$D^\times/F^\times K^{m+1}$ is finite, the representations of level
$\le m$ with a given central character are the representations of a
finite group; for the running example $D^\times/F^\times K^{m+1}$ has
order $8\cdot|Q/Q_{m+1}|=8\cdot3^{m+\lceil m/2\rceil}$ and the
representations of level $\le2$ with trivial central character are the
irreducible representations of the group $\mathrm{Heis}\rtimes D_8$ of
order $216$: five of level zero (dimensions $1,1,1,1,2$, the
representations of $D_8$), four of level one (dimension $4$, those
nontrivial on $K^1/K^2\simeq\Fq_9$ but trivial on $K^2$) and four of
level two (dimension $6$, those nontrivial on $K^2/K^3$).  The level is
the integer form of the Moy--Prasad depth: for $D^\times$ with its
standard filtration a representation of level $m$ has depth $m/r$, and
depth is preserved by the local Jacquet--Langlands correspondence
\cite[Ch.~13]{BH06}; \cite{DKV84}.  For $p\nmid r$, which includes the running
example ($r=2$, $p=3$), Howe and Corwin \cite{How77b}; \cite{Cor74} parametrize the
irreducible representations of $D^\times$ by admissible pairs $(E,\chi)$,
$E\subset D$ a subfield and $\chi$ a character of $E^\times$, with
$\rho_{E,\chi}$ induced from a character of $E^\times K^{m'}$ extending
$\chi$; the level of $\rho_{E,\chi}$ is governed by the conductor of
$\chi$, and this induction through a congruence subgroup is the model
for the inducing data in Sections~\ref{W:sec:maintheorem}
and~\ref{W:sec:plan}.  Levels are unbounded: every irreducible
representation has some finite level, but no finite set of levels
exhausts $\Irr(D^\times)$, which is why the results of
Section~\ref{W:sec:plan}, which cover levels $\le2$, do not by themselves
give Target~B and why Section~\ref{W:sec:future} asks for an inductive
mechanism.

Over a number field $F$ the Brauer group is described by the
Albert--Brauer--Hasse--Noether theorem: the sequence
$0\to\mathrm{Br}(F)\to\bigoplus_v\mathrm{Br}(F_v)\xrightarrow{\ \Sigma\ }\Q/\Z\to0$
is exact, so a central simple algebra is determined by its local
invariants, which vanish at all but finitely many places and sum to
zero; the period of a class equals its index, so a class of order $r$ is
the class of a division algebra of degree $r$ \cite[\S32]{Rei75}; \cite[Ch.~XIII]{Wei73}.
A quaternion algebra over a number field is ramified at an even number
of places (where its invariant is $\tfrac12$) and split elsewhere.  The
local components $D_v=D\otimes F_v$ at ramified $v$ are the local
division algebras above; the group $D^\times(F)\backslash D^\times(\mathbb A_F)/F^\times(\mathbb A_F)$
is compact, so the automorphic spectrum of $D^\times$ is discrete, and
the Jacquet--Langlands correspondence \cite[Ch.~III]{JL70}; \cite{DKV84}; \cite{Rog83}; \cite{Bad08}
identifies the automorphic representations of $D^\times$ with those of
$\mathrm{GL}_r$ that are discrete series at the ramified places; its
local form is a bijection between the irreducible representations of
$D_v^\times$ and the discrete series of $\mathrm{GL}_r(F_v)$ characterized
by the character relation $\Theta_\pi(g)=(-1)^{r-1}\Theta_\rho(g')$ on
matching regular elliptic elements.  Chapters~\ref{ch:I}--\ref{ch:IV} concern
the metaplectic analogue of this picture: an $n$-fold cover of
$D^\times$ and the Kazhdan--Patterson cover of $\mathrm{GL}_r$, related
by a transfer of characters.  The present chapter is purely local and
concerns one ramified place of the running example.

\subsection{The cover, the star map and good elements}\label{W:subsec:coverbackground}
We use the following from Chapter~\ref{ch:I}.  Let $\mu_n\subset F$.  The
$n$-fold scalar-central cover $\tH\to H=D^\times$ is the central
extension by $\mu_n$ obtained by restricting the Kazhdan--Patterson cover
of $\mathrm{GL}_r(E_0)$ \cite{KP84} along $D\hookrightarrow M_r(E_0)$; its
preimage $\tZ$ of the center $F^\times$ is central, and
$\tZ\simeq\mu_n\times F^\times$ noncanonically (\cref{I:prop:restriction}).  The
restriction of the cover to a maximal torus $E^\times$ is a Heisenberg
group: the commutator of lifts is a bimultiplicative pairing
$E^\times\times E^\times\to\mu_n$ with kernel $F^\times E^{\times n}$ on
either side (\cref{I:prop:comm} and \eqref{W:eq:ranktwo-comm}), so the
covered torus has center $C_E$ of index $|n|_F^{-(r-1)}$ in
$\widetilde{E^\times}$, and every genuine irreducible representation of
$\widetilde{E^\times}$ has degree $b_D=|n|_F^{-(r-1)/2}$ and is determined
by its central character (Stone--von Neumann; see
Section~\ref{W:subsec:howe}).  The \emph{star map} $x\mapsto x^*=\tilde x^{\,n}$
is well defined on $H$ (independent of the lift, since $\mu_n$ is
central) and is a homomorphism on each covered torus, because
$(\tilde x\tilde y)^n=\tilde x^n\tilde y^n[\tilde y,\tilde x]^{n(n-1)/2}$
with $[\tilde y,\tilde x]\in\mu_n$ and $n$ odd; an element of $H$ is \emph{good} if it
lies in $F^\times\cdot\{x^n\}$, and the division-side transform
\eqref{W:eq:LD} is supported on the star lifts $\tilde z\,t^*$ of good
elements.  In the running example $n=3$, $b_D=3$, $c_D=27$.

\subsection{Heisenberg groups, the Weil representation and the Howe correspondence}\label{W:subsec:howe}
This subsection recalls the Heisenberg group, the Stone--von Neumann theorem, the Weil representation and the Howe correspondence. They serve as a heuristic for the later sections, and Remark~\ref{W:rem:howe-duality} uses them to restate Target~B as a duality.

Let $(W,\langle\ ,\ \rangle)$ be a finite-dimensional symplectic space
over a local field or a finite field $k$, and $\mathcal H(W)=W\times k$
the Heisenberg group, $(w,s)(w',s')=(w+w',s+s'+\tfrac12\langle w,w'\rangle)$.
The Stone--von Neumann theorem says that for each nontrivial character
$\psi$ of the center $k$ there is a unique irreducible (smooth, for local
$k$) representation $\rho_\psi$ of $\mathcal H(W)$ with central character
$\psi$; its character is supported on the center.  It is realized in the
\emph{Schr\"odinger model}: for a Lagrangian decomposition $W=X\oplus Y$,
$\rho_\psi=\Ind_{Y\times k}^{\mathcal H(W)}(\mathbf1\otimes\psi)$, acting on
functions on $X$ \cite{Wei64}; \cite[Ch.~2]{MVW87}.  The symplectic group $\mathrm{Sp}(W)$
acts on $\mathcal H(W)$ fixing the center, so by uniqueness it acts
projectively on $\rho_\psi$; the resulting genuine representation of the
metaplectic double cover $\mathrm{Mp}(W)$ is the Weil representation
$\omega_\psi$.  For a \emph{dual pair} $(G,G')$ of mutually
centralizing reductive subgroups of $\mathrm{Sp}(W)$, Howe's duality
conjecture \cite{How79}, proved by Waldspurger for $p\ne2$ \cite{Wal90} and in
general by Gan and Takeda \cite{GT16}, says that the restriction of
$\omega_\psi$ to $\tilde G\times\tilde G'$ decomposes as
$\bigoplus_\pi\pi\otimes\theta(\pi)$ with $\pi\mapsto\theta(\pi)$
injective on the set of $\pi$ that occur: a bijection between
representations of the two members of the pair, with multiplicity one,
whose ``kernel'' is the character of $\omega_\psi$.  The first
nontrivial instance is Shimura's correspondence: Waldspurger realized
the correspondence between genuine representations of
$\widetilde{\mathrm{SL}}_2$ and representations of $\mathrm{PGL}_2$,
or of $PD^\times$ for the quaternion algebra $D$, through the dual pairs
$(\widetilde{\mathrm{SL}}_2,O(V_3))$ with $V_3$ a three-dimensional
quadratic space, split or anisotropic \cite{Wal80}; \cite{Wal91}.  So for $n=2$ the
link between a metaplectic cover and a quaternion algebra is literally
a Howe correspondence, with the cover on the split side.  For odd $n$
no classical dual pair produces an $n$-fold cover, and the metaplectic
correspondence of Flicker and Kazhdan \cite{FK86} is established by the trace
formula instead; the transform $L_D$ of Chapter~\ref{ch:III} is likewise defined
analytically, by \eqref{W:eq:LD}, not by a kernel.

Three points of contact with the present chapter.  First, the covered
tori are Heisenberg groups and the degree-$b_D$ representations of
Section~\ref{W:sec:fourier} are their Stone--von Neumann representations.
Second, the finite groups of Sections~\ref{W:sec:depth}--\ref{W:sec:finite}
have a Heisenberg flavor: $Q_1/Q_3$ is the Heisenberg group of order
$27$, and the character $\lambda_0$ of $\tU_3$ with
$\tau|_{\tU_1}=\Ind_{\tU_3}^{\tU_1}\lambda_0$ is reminiscent of a
Schr\"odinger model, a character of an abelian normal subgroup of index
$27$ induced up, with the freeness statement of
Proposition~\ref{W:prop:free} playing the role of nondegeneracy (every
genuine character of $\tU_3$ of the relevant level is regular, so every
such induced representation is irreducible).  This is an analogy, not an
identification: $\tU_3/N_6\simeq\mu_3\times V$ is abelian by (H3), $V$
has dimension five, and $\tU_1/N_6$ is not a Heisenberg group with a
nondegenerate symplectic quotient; what is used is ordinary Clifford
theory for an abelian normal subgroup.  Third, and most
to the point, Target~B is itself a Howe-duality statement about the pair
$(\tH,H)$: if $L_D\Theta_\rho=\Theta_{\tau_\rho}$ with $\tau_\rho$
irreducible for every $\rho$, then $\Omega:=\bigoplus_\rho\tau_\rho\otimes\rho^\vee$
is a representation of $\tH\times H$ in which each $\rho$ is paired
with exactly one $\tau_\rho$, with multiplicity one, and its
level-by-level truncations have the explicit characters computed in
Remark~\ref{W:rem:howe-duality} (the full sum is a formal object; see
there).
Whether $\Omega$ has an intrinsic construction, as $\omega_\psi$ has,
is not known to us; the induced models of
Theorems~\ref{W:thm:levelone} and~\ref{W:thm:leveltwo} are what such a
construction would have to produce level by level.  Finally, the
obstruction that distinguishes the wild place is the one that limits
Howe's Kirillov theory for compact $p$-adic groups \cite{How77a}: at a tame
place the $n$-th root transform of characters is a dilation of coadjoint
orbits by $1/n$, hence a Galois conjugation of characters (\cref{III:lem:tame-adams}); at $p\mid n$ the dilation is not available on the Lie
lattice, and Sections~\ref{W:sec:finite}--\ref{W:sec:plan} show how the
corresponding prediction nevertheless survives, with the character
$\lambda_0$ absorbing the nonlinearity of the cube root.

\section{Division-side character transform and its normalization}\label{W:sec:transform}
This section recalls the transform $L_D$ of Chapter~\ref{ch:III} and its normalizing constants, in the notation of this chapter. It then recalls the criterion that the rest of the chapter tests. The transform permutes irreducible characters exactly when its matrix is integral.

Let $F$ be nonarchimedean, $D/F$ a division algebra of prime
degree $r$, $H=D^\times$, and $\tH$ its specified cover.
Set $Z=F^\times$ and $\tZ=p^{-1}(Z)$, equipped with compatible
Haar measures on $H/Z\simeq\tH/\tZ$.  The Chapter~\ref{ch:III} transform is
\begin{equation}\label{W:eq:LD}
 (L_D\Phi)(z\,t^*)
 =b_D\,\tth(z)
  \frac{\Delta(t)}{\Delta(t^n)}\Phi(t),\qquad
 b_D=|n|_F^{-(r-1)/2},
\end{equation}
for regular $t$, $z\in\tZ$; it vanishes at nongood regular lifts.
Here $t^*=\tilde t^{\,n}$ is the star of any lift $\tilde t$ of $t$, and
\begin{equation}\label{W:eq:Delta}
 \Delta(t)=\left|\det\bigl(1-\operatorname{Ad}(t)
 \mid\mathfrak h/\mathfrak t\bigr)\right|_F^{1/2}
 =\prod_{i<j}\left|
 \frac{(\lambda_i-\lambda_j)^2}{\lambda_i\lambda_j}
 \right|_F^{1/2}.
\end{equation}
This equals Chapter~\ref{ch:I}'s $|D(t)|^{1/2}|\Nrd(t)|_F^{-(r-1)/2}$.
Since the power map on a field torus modulo scalars has index
$|n|_F^{-(r-1)}$, Weyl integration gives $b_D^2$
equal to this index.  The factor $b_D$ is independent of
\emph{common} rescalings of the two identified quotient Haar
measures.  The ratio of discriminants extends at central elements
to $|n|_F^{-r(r-1)/2}$ (\cref{III:lem:ratio}), hence
\begin{equation}\label{W:eq:cD}
 (L_D\Theta_\rho)(1)=c_D\dim\rho,\qquad
 c_D=b_D|n|_F^{-r(r-1)/2}
 =|n|_F^{-(r^2-1)/2}.
\end{equation}
The identity factor is pointwise and does not depend on Haar
measures.  \cref{III:thm:test-transfer}, separately gives the adjoint
orbital-integral/test-function transfer.  Test functions transform
by $\tth^{-1}$, not by $\tth$.

\subsection{Integral-matrix criterion}\label{W:subsec:integral}
This subsection writes $L_D$ as a matrix in the bases of irreducible characters and recalls the integrality criterion of Chapter~\ref{ch:III}. The criterion turns a question about characters into a question about integers.

Choose $\tth$ first and put
$\theta(a)=\tth(a^*)$.
\cref{III:thm:unitary} and \cref{III:cor:matrix}, show that $L_D$ is a
unitary \emph{isomorphism} of invariant $L^2$ spaces and that
\begin{equation}\label{W:eq:matrixchars}
 L_D\Theta_\rho=\sum_\tau a_{\rho\tau}\Theta_\tau,
 \qquad a_{\rho\tau}=\langle L_D\Theta_\rho,
 \Theta_\tau\rangle,
\end{equation}
with each row and column finite.  The characters on either side
are orthonormal bases in their fixed-central-character spaces.

\begin{proposition}[\Etag\ Integral-lattice criterion, Chapter~\ref{ch:III}]\label{W:prop:int}
The transform $L_D$ permutes the irreducible characters with
positive signs if and only if every $a_{\rho\tau}$ in
\eqref{W:eq:matrixchars} is an integer.
\end{proposition}
\begin{proof}
An integral unit vector in $\ell^2$ is a signed basis vector;
unitarity gives injectivity and surjectivity because the inverse
is the adjoint of a unitary map.  The value at identity in
\eqref{W:eq:cD} is positive and fixes the sign.
This is \cref{III:cor:integrality}, with its converse.
\end{proof}

\begin{remark}[The actual obstruction]
For wild $v$, Chapter~\ref{ch:III} does not prove
$a_{\rho\tau}\in\Z$.  The tame Adams argument works with a
finite quotient $Q$ and a central subgroup $C$, under the
hypothesis $\gcd(n,|Q/C|)=1$ (\cref{III:lem:tame-adams}).
At a wild place that hypothesis need not hold.  For the
simplest illustration take $Q=C_p$, $C=1$, $n=p$: then the
Adams operation sends every linear character to the trivial
character.  This does not prove a nonintegral coefficient for
$L_D$; it only shows why that particular proof method fails.
The present chapter replaces the Adams argument, in the running
example, by Brauer's characterization of characters and an
explicit analysis of the pro-$3$ part of the cover.
\end{remark}

\section{A useful global criterion: positivity of every row}\label{W:sec:positive}
This section proves that integrality of the whole character matrix is equivalent to positive definiteness of every row. It also explains why the positivity of a single row decides nothing about that row.

\begin{proposition}[Positivity versus integrality]\label{W:prop:positive}
With the unitary character matrix $(a_{\rho\tau})$ of
\eqref{W:eq:matrixchars}, the following are equivalent:
\begin{enumerate}[label=(\alph*)]
\item $L_D$ permutes irreducible characters with positive signs;
\item every $a_{\rho\tau}$ is an integer;
\item every $a_{\rho\tau}$ is a nonnegative real number;
\item every $L_D\Theta_\rho$ is positive definite on $\tH$.
\end{enumerate}
The implication from positive definiteness to integrality
requires \emph{every} row; positivity of a single row does not
force integrality or irreducibility for that row.
\end{proposition}
\begin{proof}
By Chapter~\ref{ch:III}, each row and column of $A$ has $\ell^2$ norm one.
An integer row is a signed basis vector, and its sign is positive
by the identity value \eqref{W:eq:cD}.  If all entries are
nonnegative, orthogonal rows have disjoint supports.  Every
column meets exactly one row, and since its norm is one its
nonzero entry is one; row norms then imply the permutation
property.  The equivalence of nonnegative character coefficients
and positive definiteness follows from Peter--Weyl/GNS for a
compact-mod-center group with fixed unitary central character.
One row could instead be $(2^{-1/2},2^{-1/2},0,\ldots)$.
\end{proof}
\begin{remark}\label{W:rem:onerow}
For a \emph{single} row the usable criterion is different and is
the one we shall apply: a class function of $\ell^2$ norm one with
positive value at the identity is an irreducible character as soon
as it is a \emph{virtual} character (an integral combination of
irreducible characters).  Brauer's theorem reduces the virtual-character
property to restrictions to elementary subgroups; Section~\ref{W:sec:reduction}
carries this out.
\end{remark}

\section{A general quadratic theorem}\label{W:sec:general-toral}
This section proves that the discriminant ratio $R_E$ is the character of a permutation representation, for every quadratic extension $E/F$ and every odd $n$. This is the toral positivity statement of the chapter. The proof computes $R_E$ on the filtration of the norm-one group, first in general and then separately for unramified and ramified extensions.

Let $F$ be a finite extension of $\Q_p$ with $p$ odd, residue field of cardinality $q$, and let $n\geq1$ be odd with $\mu_n\subset F$. Let $E/F$ be any quadratic field extension, with nontrivial Galois involution $\sigma$. Put
\[
 A=E^\times/F^\times,\qquad G=E^1=\ker(N_{E/F}:E^\times\to F^\times).
\]
Hilbert's Theorem 90 identifies $A\xrightarrow{\sim}G$ via $t\mapsto u=t/\sigma(t)$. The normalized absolute value on $E$ is $|x|_E=q_E^{-v_E(x)}$, so $|x|_E=|N_{E/F}x|_F$.

For noncentral $t\in E^\times$ let
\[
 R_E(t)=\frac{\Delta(t)}{\Delta(t^n)},\qquad
 \Delta(t)=\bigl|\det(1-\operatorname{Ad}(t)\mid D/E)\bigr|_F^{1/2}.
\]
The ratio extends continuously at scalar elements, with value $|n|_F^{-1}$.

\begin{theorem}[Quadratic toral permutation positivity]\label{W:thm:main}
For any $E/F,n$ as above, $R_E$ is the character of a finite-dimensional permutation representation of a finite quotient of $A$. In particular,
\[
 R_E(t)=\sum_{\chi\in\widehat A}m_\chi\chi(t),\qquad m_\chi\in\mathbb Z_{\geq0},
\]
where only finitely many $m_\chi$ are nonzero. This includes every wild case $p\mid n$. At the identity $R_E(1)=|n|_F^{-1}$.
\end{theorem}

\subsection{A radial valuation formula}
This subsection shows that $R_E(u)$ depends only on the depth $v_E(1-u)$ of $u$ in the norm-one group, and gives the formula. The two subsections that follow turn the formula into a sum of permutation characters.

Write $n=p^a m$ with $(m,p)=1$, and put
\[
 t_\zeta=v_E(1-\zeta),\qquad
 G_k=\{u\in G:v_E(1-u)\geq k\}\quad(k\geq1),\qquad G_0=G.
\]
If $a=0$ then $R_E\equiv1$, so assume $a\geq1$.

\begin{lemma}\label{W:lem:valuations}
For $u\in G$ and any $\zeta\in\mu_n$ whose prime-to-$p$ component is nontrivial, $|1-\zeta u|_E=1$. For $\zeta\in\mu_{p^a}\setminus\{1\}$,
\[
 v_E(1-\zeta u)=\min\{v_E(1-u),t_\zeta\}.
\]
Consequently, with $k=v_E(1-u)\in\mathbb Z_{\geq0}\cup\{\infty\}$,
\begin{equation}\label{W:eq:radial}
 R_E(u)=q_E^{\frac12\sum_{\zeta\in\mu_{p^a}\setminus\{1\}}\min(k,t_\zeta)}.
\end{equation}
\end{lemma}
\begin{proof}
The reduction of $u\in E^1$ lies in a norm-one subgroup. In the unramified case its intersection with $k_F^\times$ is $\{\pm1\}$, and in the ramified case the reduction of $u$ is itself in $\{\pm1\}$. A root of unity with nontrivial prime-to-$p$ component has reduction in $k_F^\times\setminus\{\pm1\}$, since $n$ is odd, proving the first assertion.

For a nontrivial $p$-power root $\zeta$, the ultrametric inequality gives the asserted minimum when the two valuations differ. If both equal $t_\zeta$ and $v_E(1-\zeta u)>t_\zeta$, then $u\equiv\zeta^{-1}\pmod{\mathfrak p_E^{t_\zeta+1}}$. Applying $\sigma$ and using $\sigma(u)=u^{-1}$ gives $u^{-1}\equiv\zeta^{-1}$ to the same precision, hence $\zeta\equiv\zeta^{-1}$, a contradiction to $v_E(\zeta-\zeta^{-1})=t_\zeta$ (as $p$ is odd). Thus equality holds.

For $u=t/\sigma(t)$ we have
$\Delta(t)=|N_{E/F}(1-u)|_F^{1/2}=|1-u|_E^{1/2}$.
Factor $1-u^n=\prod_{\zeta\in\mu_n}(1-\zeta u)$ and use the two valuation statements. The identity follows for noncentral $t$ and extends continuously to $u=1$. At $u=1$, the cyclotomic product $\prod_{\zeta\in\mu_{p^a}\setminus\{1\}}(1-\zeta)=p^a$ gives $R_E(1)=|p^a|_E^{-1/2}=|n|_F^{-1}$.
\end{proof}

Define for $k\geq0$
\[
 M_k=q_E^{\frac12\sum_{\zeta\in\mu_{p^a}\setminus\{1\}}\min(k,t_\zeta)},\qquad M_0=1,
\]
and put $T=\max_{\zeta\ne1}t_\zeta$. All values $M_k$ are positive integers because roots $\zeta,\zeta^{-1}$ occur in pairs of equal valuation. We have $M_k=M_T$ for $k\geq T$.

\subsection{The unramified filtration}
This subsection treats an unramified extension $E/F$. It computes the indices of the filtration subgroups and writes $R_E$ as a nonnegative integral combination of the permutation characters of $G$ on the quotients $G/G_k$.

Suppose $E/F$ is unramified. Then $q_E=q^2$ and reduction of the norm-one group has order $q+1$; at every deeper level the trace-zero subgroup of the residue-field additive group has order $q$. Thus the filtration satisfies
\[
 [G:G_1]=q+1,\qquad [G:G_k]=(q+1)q^{k-1}\quad(k\geq1).
\]
Write $s_k=\#\{\zeta\in\mu_{p^a}\setminus\{1\}:t_\zeta\geq k\}$. Each $s_k$ is even, and $M_k/M_{k-1}=q^{s_k}$.
Set
\[
 c_k=\frac{M_k-M_{k-1}}{(q+1)q^{k-1}}\quad(1\leq k\leq T).
\]
For $k=1$, $M_0=1$ and $q+1$ divides $q^{s_1}-1$ because $s_1$ is even. For $k>1$ with $s_k>0$, $M_{k-1}$ is divisible by $q^{k-1}$, since at least two roots contribute at least $k-1$ to its exponent. Again $q+1$ divides $q^{s_k}-1$. Hence $c_k\in\mathbb Z_{\geq0}$ (and $c_k=0$ if $s_k=0$).

By telescoping \eqref{W:eq:radial},
\begin{equation}\label{W:eq:unram}
 \boxed{R_E=\mathbf1_G+\sum_{k=1}^T c_k\,[G:G_k]\,\mathbf1_{G_k}.}
\end{equation}
Here $[G:G_k]\,\mathbf1_{G_k}$ is precisely the permutation character of $G$ on $G/G_k$.

\subsection{The ramified filtration}
This subsection treats a ramified extension $E/F$. The positive depths are now odd, so the indices change, but the conclusion has the same form as in the unramified case.

Suppose $E/F$ is ramified. Then $q_E=q$ and $t_\zeta=2v_F(1-\zeta)$ is even. The norm-one group has filtration
\[
 [G:G_1]=2,\qquad G_{2j}=G_{2j+1},\qquad [G:G_{2j}]=2q^j\quad(j\geq1).
\]
For completeness, write $E=F(\sqrt d)$ with $v_F(d)=1$. Hilbert 90 gives $u=(a+b\sqrt d)/(a-b\sqrt d)$; normalize $a,b$ to be integral and not both divisible by the uniformizer of $F$. If $a$ is a unit, scale it to $1$ and write $c=b/a\in\mathcal O_F$: then the depth of $u-1$ is $1+2v_F(c)$. Moreover, two such parameters $c,c'$ define elements congruent modulo $\mathfrak p_E^{2j}$ precisely when $c\equiv c'\pmod{\mathfrak p_F^j}$, so $[G_1:G_{2j}]=q^j$. If $a$ is not a unit and $b$ is, then $u\equiv-1$ modulo $\mathfrak p_E$. These observations give the displayed indices and the fact that positive depths are odd.

Define
\[
 c_1=\frac{M_1-1}{2},\qquad
 c_{2j}=\frac{M_{2j+1}-M_{2j-1}}{2q^j}\quad(1\leq j\leq T/2).
\]
The first coefficient is an integer because $M_1=q^{(p^a-1)/2}$ is odd. For $j\geq1$ write
\[
 \frac{M_{2j+1}}{M_{2j-1}}=q^{e_j},\qquad
 e_j=\frac12\sum_{\zeta\ne1}\bigl(\min(2j+1,t_\zeta)-\min(2j-1,t_\zeta)\bigr)\in\mathbb Z_{\geq0}.
\]
If $e_j=0$, then $c_{2j}=0$. Otherwise some pair $\{\zeta,\zeta^{-1}\}$ has $t_\zeta\geq2j$, so $M_{2j-1}$ is divisible by $q^{2j-1}$, hence by $q^j$. Moreover $q^{e_j}-1$ is even. Therefore $c_{2j}\in\mathbb Z_{\geq0}$.

As all positive depths are odd, telescoping now gives
\begin{equation}\label{W:eq:ram}
\boxed{R_E=\mathbf1_G+c_1\,[G:G_1] \mathbf1_{G_1}
 +\sum_{j=1}^{T/2}c_{2j}\,[G:G_{2j}]\,\mathbf1_{G_{2j}}.}
\end{equation}
Again each summand is the character of a permutation representation. Equations \eqref{W:eq:unram} and \eqref{W:eq:ram} prove Theorem~\ref{W:thm:main}.

\section{From Fourier coefficients to covered-torus restrictions}\label{W:sec:fourier}
This section explains what the Fourier coefficients of $R_E$ mean on the cover. They are the multiplicities with which genuine representations of the covered torus occur in the restriction of a transferred character. The section also states the limit of this test. It concerns restrictions to a torus and does not decide integrality on the whole cover.

\begin{lemma}[Fourier test for restrictions to covered field tori]\label{W:lem:toral-Fourier}
Let $E/F$ be any maximal field in $D$, with $r$ prime and
$(r,n)=1$, and put $A_E=E^\times/F^\times$.  The function
\[
 R_E(t)=\frac{\Delta(t)}{\Delta(t^n)}
\]
extends continuously at scalars, is a locally constant function
on the compact abelian group $A_E$, and has a finite Fourier expansion
\begin{equation}\label{W:eq:R-Fourier}
 R_E(t)=\sum_{\chi\in\widehat{A_E}}m_\chi\chi(t),
 \qquad
 m_\chi=\int_{A_E}R_E(t)\overline{\chi(t)}\,dt.
\end{equation}
Here the Haar measure of $A_E$ has mass $1$.
Choose a compatible genuine central character
$\tth$ of $\widetilde Z$ and set
$\theta(a)=\tth(a^*)$.  For any smooth character
$\alpha:E^\times\to\mathbb C^\times$ with
$\alpha|_{F^\times}=\theta$, and for each $\chi\in\widehat{A_E}$,
let $\eta_{\alpha\chi}$ be the unique genuine irreducible
representation of $\widetilde{E^\times}$ whose central character
satisfies
\[
 \zeta_{\alpha\chi}(z\,t^*)=
 \tth(z)\alpha(t)\chi(t)
 \qquad(z\in\widetilde Z,\ t\in E^\times).
\]
Define a torus-level class function $\mathcal L_E\alpha$ to equal
$b_D\tth(z)R_E(t)\alpha(t)$ on $z t^*\in C_E$
and to vanish outside $C_E$.  Then, as a genuine class function
on the \emph{covered torus},
\begin{equation}\label{W:eq:toral-branch}
 \mathcal L_E\alpha
 =\sum_\chi m_\chi\,\Theta_{\eta_{\alpha\chi}}.
\end{equation}
For an ordinary character $\Theta_\rho$ of $D^\times$,
its actual $L_D$ transform restricts to the sum of these
$\mathcal L_E\alpha$ with the multiplicities in
$\Theta_\rho|_{E^\times}$.
Consequently $m_\chi\in\mathbb Z_{\geq0}$ for every $\chi$
is sufficient for the restriction of every transferred ordinary
irreducible character to this covered torus to be a genuine
representation character.  It is not sufficient for the full
$\widetilde{D^\times}$ character to be a representation character.
\end{lemma}
\begin{proof}
The extension of $R_E$ at scalars follows from \cref{III:lem:ratio}.  Since $E/F$ is a degree-$r$ field extension with
$(r,n)=1$, the $n$th power of a noncentral $t$ stays noncentral.
The compactness of $A_E$ and local constancy imply that $R_E$
factors through a finite quotient, giving \eqref{W:eq:R-Fourier}.

By Chapter~\ref{ch:I} the center of the covered torus is
$C_E=\widetilde Z\cdot\{t^*:t\in E^\times\}$.  The displayed
formula for $\zeta_{\alpha\chi}$ is well-defined: the
power map on $A_E$ is injective, and on scalar overlaps the
compatibility $\alpha(a)=\tth(a^*)$ is precisely
what is required.  Such genuine central characters extend over
$C_E$, and Stone--von Neumann gives a unique $\eta_{\alpha\chi}$.
Its character vanishes off $C_E$, while on $C_E$ it equals
$b_D\zeta_{\alpha\chi}$ because its degree is $b_D$.
The defining formula \eqref{W:eq:LD} for $L_D$ has the identical
vanishing and equals $b_D\tth(z)R_E(t)\alpha(t)$
on $z t^*$.  Comparing with \eqref{W:eq:R-Fourier} proves
\eqref{W:eq:toral-branch}.  Finally, restricting a finite-dimensional
ordinary representation to $E^\times$ decomposes it into
one-dimensional characters, all restricting to $\theta$ on
$F^\times$; summing the displayed formula proves the last claim.
\end{proof}

\section{The explicit cubic wild example}\label{W:sec:examplewild}
The following calculation verifies a \emph{necessary representation-ring
integrality test on every maximal torus}, including the two ramified
quadratic field types.  Sections~\ref{W:sec:depth}--\ref{W:sec:maintheorem}
then settle the full-group statement.

\begin{example}[The field, algebra, and cover]\label{W:ex:ranktwo-data}
Take
\[
 F=\Q_3(\zeta_3),\qquad \varpi=\zeta_3-1,\qquad
 3=-\zeta_3^{-1}\varpi^2,\qquad q_F=3.
\]
Thus $v_F(3)=2$ and $|3|_F=3^{-2}=1/9$.  The compatible
covering parameters in \cite{JC26a} and Chapter~\ref{ch:I} are respectively
$c=1$ and $c_v=2\pmod3$.  Let $D/F$ be the unique quaternion
division algebra.  With $E_0=F(i)$, $i^2=-1$, its unramified
quadratic maximal subfield, we may write
\[
 D=E_0\oplus E_0j,\qquad j^2=\varpi,\qquad
 jx=\sigma(x)j\quad(x\in E_0),
\]
where $\sigma$ is the nontrivial automorphism of $E_0/F$.
\cref{I:prop:comm}, gives, for every quadratic maximal
subfield $E/F$, the covered-torus commutator
\begin{equation}\label{W:eq:ranktwo-comm}
 [x,y]_{\widetilde{E^\times}}=
 (N_{E/F}x,N_{E/F}y)_{F,3}^{\,2}
 (x,y)_{E,3}^{-1}.
\end{equation}
The center $C_E$ of $\widetilde{E^\times}$ is the preimage of
$F^\times E^{\times3}$ (\cref{I:prop:good}), and
\[
 [E^\times:F^\times E^{\times3}]=9,
 \qquad b_D=3,\qquad c_D=27.
\]
By the Heisenberg-type Stone--von Neumann theorem
(\cite[Proposition~2.4]{JC26a}; \cref{I:prop:good}), every genuine
irreducible representation of $\widetilde{E^\times}$ has degree $3$.
This is \emph{not} a degree calculation for representations of
$\widetilde{D^\times}$.
\end{example}

\subsection{The unramified quadratic torus}
This subsection computes $R_{E_0}$ on the unramified quadratic torus of the example. It then decomposes $R_{E_0}$ into characters of a finite quotient.

Put $E_0=F(i)$ as above.  Since $X^2+1$ is irreducible over
$\mathbb F_3$, $E_0/F$ is unramified and its residue field is
$\mathbb F_9$.  For $t=a+bi$ with $a,b\in F$ the reduced norm
and trace give the elementary identity
\begin{equation}\label{W:eq:quadratic-R}
 R_{E_0}(a+bi)=
 \frac{|a^2+b^2|_F}{|3a^2-b^2|_F}.
\end{equation}
Multiplying $t$ by an element of $F^\times$, suppose
$a,b\in\mathcal O_F$ and at least one is a unit.
Then $a^2+b^2$ is always a unit.  If $b$ is a unit, so is
$3a^2-b^2$.  If $b=\varpi u$, then
\[
 3a^2-b^2=3(a^2+\zeta_3u^2),
\]
and $a$ is a unit; the reduction of $a^2+\zeta_3u^2$ modulo
$\varpi$ is either $1$ or $2$, never $0$.  Hence
\begin{equation}\label{W:eq:unram-values}
 R_{E_0}(a+bi)=
 \begin{cases}
 9,&b\in\varpi\mathcal O_F,\\
 1,&b\in\mathcal O_F^\times.
 \end{cases}
\end{equation}
The finite residue quotient is
\[
 B_0=k_{E_0}^\times/k_F^\times
 \simeq C_4.
\]
The condition $b\in\varpi\mathcal O_F$ says exactly that the
image of $t$ in $B_0$ is the identity.  Denote its four characters
by $\chi_0=1,\chi_1,\chi_2,\chi_3$.  The regular character of
$C_4$ equals $4$ at the identity and $0$ elsewhere. Thus
\begin{equation}\label{W:eq:unram-fourier}
 \boxed{\quad
 R_{E_0}=1+2\,\mathrm{reg}_{C_4}
 =3\chi_0+2\chi_1+2\chi_2+2\chi_3.
 \quad}
\end{equation}
This is a positive integral Fourier decomposition of degree
$3+2+2+2=9$; it is stronger than merely checking that
$R_{E_0}$ takes integer values.

\subsection{Both ramified quadratic tori}
This subsection does the same computation for the two ramified quadratic extensions of $F$. Proposition~\ref{W:prop:all-three-tori} then collects the three cases.

Every ramified quadratic extension of $F$ can be written
$E_d=F(j_d)$, $j_d^2=d=u\varpi$ with $u\in\mathcal O_F^\times$.
There are two isomorphism classes, distinguished by the square
class of $u$.  For $t=a+bj_d$,
\begin{equation}\label{W:eq:ramified-R}
 R_{E_d}(a+bj_d)
 =\frac{|a^2-db^2|_F}{|3a^2+db^2|_F}.
\end{equation}
Again normalize $a,b$ to be integral and not both divisible
by $\varpi$.  The terms $3a^2$ and $db^2$ have different
valuations whenever both $a$ and $b$ are units.  A direct
valuation calculation therefore gives
\begin{equation}\label{W:eq:ramified-values}
 R_{E_d}(a+bj_d)=
 \begin{cases}
 9,& a\in\mathcal O_F^\times,
       \quad b\in\varpi\mathcal O_F,\\
 3,& a,b\in\mathcal O_F^\times,\\
 1,& a\in\varpi\mathcal O_F,
       \quad b\in\mathcal O_F^\times.
 \end{cases}
\end{equation}
For example, in the original cyclic algebra $d=\varpi$,
\[
 \begin{array}{c|ccc}
 t&j&1+j&1+\varpi j\\\hline
 R(t)&1&3&9\\
 (L_D1)(t^*)&3&9&27
 \end{array}
\]
for the central-character normalization of Example~\ref{W:ex:ranktwo-data}.
The three values illustrate why the unramified torus alone is not
an adequate test.

There is a natural finite quotient
\[
 B_d=E_d^\times/(F^\times U_1)\simeq C_2\times C_3,
\]
where $U_1\subset\mathcal O_{E_d}^\times$ consists of
$a+bj_d$ with $a\in\mathcal O_F^\times$ and
$b\in\varpi\mathcal O_F$; equivalently
$U_1=\mathcal O_F^\times(1+\mathfrak p_{E_d}^2)$.  The $C_2$ coordinate
records the parity of the $E_d$-valuation, and on even elements
the $C_3$ coordinate is $\bar c\in\mathbb F_3$ for
$1+cj_d$ with $c\in\mathcal O_F$ read modulo $\varpi$.  The multiplication law
\[
 (1+cj_d)(1+c'j_d)
 =(1+dcc')\left(1+\frac{c+c'}{1+dcc'}j_d\right)
\]
shows that $\bar c$ is additive, and $j_d^2=d\in F^\times$
provides the order-two factor.

Write $\epsilon$ for the nontrivial character of $C_2$
and $\xi$ for a character of order $3$ of $C_3$.
The identity \eqref{W:eq:ramified-values} says that $R_{E_d}$
has values $9$ at the identity of $B_d$, $3$ at its two
nonidentity elements in $C_3$, and $1$ on its three odd elements.
Consequently
\begin{equation}\label{W:eq:ram-fourier}
 \boxed{\quad
 \begin{aligned}
 R_{E_d}
 &=1+\mathrm{reg}_{C_2}+\mathrm{reg}_{C_2\times C_3}\\
 &=3\cdot1+2\epsilon
   +\xi+\xi^2+\epsilon\xi+\epsilon\xi^2.
 \end{aligned}\quad}
\end{equation}
Here the regular representations are inflated from the indicated
quotients.  All coefficients are nonnegative integers and sum
to $3+2+1+1+1+1=9$.  This applies to \emph{both} ramified
quadratic isomorphism classes, independently of $u$.

\begin{proposition}[All quadratic torus restrictions pass
an integral-character test]\label{W:prop:all-three-tori}
In the wild example $F=\Q_3(\zeta_3)$, $r=2$, $n=3$,
let $D/F$ be the quaternion division algebra.  The function
$R_E$ is the character of a nine-dimensional permutation
representation of a finite quotient of $E^\times/F^\times$,
for every quadratic maximal subfield $E\subset D$.
More precisely, for unramified $E$ it is
$1\oplus 2\,\mathbb C[C_4]$, and for either ramified $E$
it is $1\oplus\mathbb C[C_2]\oplus\mathbb C[C_2\times C_3]$.
For a suitable genuine central character inducing the trivial
ordinary central character, the restriction of $L_D\mathbf1$
to $\widetilde{E^\times}$ is therefore an \emph{actual genuine
representation character} of total dimension $27$, namely
\[
 \begin{array}{c|c}
 E/F&\text{multiplicities of three-dimensional genuine torus irreducibles}\\\hline
 \text{unramified}&(3,2,2,2)\\
 \text{ramified, either type}&(3,2,1,1,1,1).
 \end{array}
\]
These multiplicities are invariant under the nontrivial
$F$-automorphism of each $E$.
\end{proposition}
\begin{proof}
A $p$-adic field of odd residue characteristic has four
square classes; hence $F$ has exactly three quadratic field
extensions, one unramified and two ramified.  Every quadratic
field extension embeds in the quaternion division algebra.
Equations \eqref{W:eq:unram-fourier} and \eqref{W:eq:ram-fourier}
prove the first assertion.  Lemma~\ref{W:lem:toral-Fourier},
together with the degree-three torus Stone--von Neumann theorem,
converts those Fourier coefficients into the displayed
multiplicities; their weighted dimensions are $3\cdot9=27$.
For the trivial ordinary character one can choose
$\tth$ with $\tth(a^*)=1$:
the odd-$n$ scalar star kills $\mu_3\subset F^\times$
(Section~\ref{I:sec:tori}), and a faithful character of
$\mu_3$ extends across the abelian scalar central group
modulo the star image.  The nontrivial Galois automorphism
acts by inversion on the respective finite quotients' character
parameters (with the $C_2$ factor fixed), so the displayed
multiplicities are invariant.
\end{proof}

\begin{remark}[Standing conventions for the rest of the chapter]\label{W:rem:conventions}
From here on $F=\Q_3(\zeta_3)$, $n=3$, $D$ is the quaternion division
algebra, $E_0=F(i)$, and $j$, $\varpi$ are as in
Example~\ref{W:ex:ranktwo-data}.  We write $\Ocal_D$ for the maximal
order, $\Pfr=\Pfr_D=j\Ocal_D$ for its maximal ideal, $v_D$ for the
valuation with $v_D(j)=1$ (so $v_D(\varpi)=2$ and $v_D(3)=4$), and
\[
 K^m=1+\Pfr^m,\qquad U_m=F^\times K^m,\qquad \tU_m=p^{-1}(U_m)\subset\tH
 \qquad(m\ge1).
\]
Chapter~\ref{ch:III}'s log records that the standard scalar cocycle is trivial;
accordingly $\tZ\simeq\mu_3\times F^\times$ and a genuine central
character has the form $\tth(\zeta,a)=\zeta\theta_0(a)$.  For
$\rho=\mathbf1$ the compatibility $\tth(a^*)=1$ forces
$\theta_0^3=1$; we fix one such $\theta_0$ (there are $81$; by
Remark~\ref{W:rem:theta-independence} the choice is immaterial) and
note that $\tth(\tZ)=\mu_3$.  The central character $\theta_\rho$
of any $\rho$ to which $L_D$ applies is trivial on $\mu_3$, since
$\theta_\rho(\zeta)=\tth(\zeta^*)=\tth(1)=1$.
\end{remark}

\section{The pro-3 structure of $K^1$ and the cube map}\label{W:sec:depth}
This section describes the pro-$3$ group on which the rest of the argument takes place. It determines the good elements of that group by studying the cube map, and it computes the transferred trivial character there.

The whole wild phenomenon lives on the pro-$3$ group
\[
 Q:=K^1F^\times/F^\times\;\simeq\;K^1/U_F^{(1)},\qquad
 Q_m:=\text{image of }K^m\quad(m\ge1),
\]
where $U_F^{(k)}=1+\pfr_F^k$ and $K^1\cap F^\times=U_F^{(1)}$.
Since $U_F^{(k)}=F\cap K^{2k}$, the successive quotients are
\begin{equation}\label{W:eq:Qgraded}
 Q_m/Q_{m+1}\simeq\begin{cases}\Fq_9,&m\text{ odd},\\ \Fq_9/\Fq_3,&m\text{ even},\end{cases}
 \qquad [Q:Q_6]=3^8,\quad [Q:Q_3]=27,\quad [Q_3:Q_6]=3^5 .
\end{equation}
Every element of $D^\times\setminus F^\times$ is regular and lies in a
unique maximal subfield; in particular, $Q$ is torsion-free, because
$E^\times/F^\times$ has no $3$-torsion for quadratic $E$ (\cref{I:lem:FcapEn}) and no $2$-torsion inside a pro-$3$ group.
The group $D^\times/F^\times K^1$ has order $8$; the subgroup
\begin{equation}\label{W:eq:P0}
 P_0:=\langle 1+i,\ j\rangle F^\times/F^\times\simeq D_8
\end{equation}
is a complement to $Q$ in $D^\times/F^\times$ (indeed $(1+i)^4=-4$,
$j^2=\varpi$ and $j(1+i)j^{-1}=2(1+i)^{-1}$ in $D^\times$), and every
finite $2$-subgroup of $D^\times/F^\times$ is conjugate into $P_0$
(Schur--Zassenhaus).

\subsection{The cube map}
This subsection collects the properties of the cube map on $Q$ in one lemma. Its main output is the description of the good elements. They have depth exactly $3$ or depth at least $6$.

\begin{lemma}\label{W:lem:cube}
\begin{enumerate}[label=(\alph*)]
\item For $k\ge3$ the cube map $K^k\to K^{k+4}$ is bijective.
\item The cube map $Q\to Q$ is injective.
\item Let $t\in Q_1\setminus Q_2$ have leading coefficient
$\bar w\in\Fq_9^\times$ (that is, $t=1+w\pi+\cdots$ with $\pi=j$).
Then $t^3\in Q_3\setminus Q_4$ with leading coefficient $\bar w^{5}$;
in particular every nonzero leading coefficient occurs.
\item $\Cubes(Q_2)=Q_6$.
\item The set $\Good(Q)$ of images of good elements is the set of cubes,
and
\[
 \Good(Q)=G_3\sqcup Q_6,\qquad G_3:=\Cubes(Q_1\setminus Q_2)\subset Q_3\setminus Q_4 .
\]
Thus a good element of $Q$ has depth exactly $3$ or depth at least $6$.
\item $G_3$ is stable under multiplication by $Q_6$.
\item $V:=Q_3/Q_6$ is an elementary abelian group of order $3^5$, and
the image $\bar G_3$ of $G_3$ in $V$ has exactly $72$ elements, with
$\bar G_3=-\bar G_3$; it spans $V$.
\end{enumerate}
\end{lemma}
\begin{proof}
(a) For $u\in\Pfr^k$, $(1+u)^3=1+3u+3u^2+u^3$ with $3u\in\Pfr^{k+4}$,
$3u^2\in\Pfr^{2k+4}$, $u^3\in\Pfr^{3k}$; for $k\ge3$ both error terms
lie in $\Pfr^{k+5}$, and $u\mapsto3u$ is a bijection
$\Pfr^k/\Pfr^{k+1}\to\Pfr^{k+4}/\Pfr^{k+5}$.  Successive approximation
gives surjectivity; injectivity follows from (b) and
$\mu_3\cap K^3=1$.

(b) If $t^3=t'^3$ in $Q$ with $t\ne1$, both $t,t'$ lie in the torus
of the regular element $t^3$ and $(t'/t)^3\in F^\times$, hence
$t'/t\in F^\times$ by \cref{I:lem:FcapEn}.

(c) With $t=1+u$, $u=w\pi+\cdots$: $3u\in\Pfr^5$, $3u^2\in\Pfr^6$ and
$u^3\equiv w\sigma(w)w\pi^3=\Nm(\bar w)\bar w\,\pi^3 \pmod{\Pfr^4}$,
since $\pi a=\sigma(a)\pi$ for $a\in E_0$; and
$\Nm(\bar w)\bar w=\bar w^{5}$, with $5$ prime to $8$.

(d) Cubes of $K^2$ lie in $K^6$.  Conversely every element of $Q_6$ is
regular or trivial and lies in some torus $A=E^\times\cap F^\times K^1/F^\times$.
For $E=E_0$ (so $E_0\cap K^6=U^{(3)}_{E_0}$): $U^{(4)}_{E_0}$ consists of
cubes of $U^{(2)}_{E_0}$ (cubing $U^{(k)}_{E_0}\to U^{(k+2)}_{E_0}$ is bijective
for $k\ge2$, as $e(E_0/\Q_3)=2$), and $U^{(3)}_{E_0}/U^{(4)}_{E_0}\simeq\Fq_9$ is
covered by $U_F^{(3)}$ (the $\Fq_3$-part) together with the cubes
$(1+\varpi w)^3\equiv1+\varpi^3(w^3-w)$, whose leading coefficients
fill the Artin--Schreier image $\Fq_3 i$.  For ramified $E$
(so $E\cap K^6=U^{(6)}_E$, $e(E/\Q_3)=4$): $U^{(7)}_E$ consists of cubes of
$U^{(3)}_E$, and $U^{(6)}_E/U^{(7)}_E\simeq\Fq_3$ is covered by $U_F^{(3)}$.
In both cases the cube root lies in $E\cap K^2$.

(e) The good elements of $H$ are $F^\times\cdot\{\text{cubes}\}$
(\cref{I:prop:good}(a), applied in each torus), so
$\Good(Q)=\Cubes(Q)$.  By (c), cubes of depth-one elements have depth
exactly $3$; by (d), cubes of $Q_2$ exhaust $Q_6$.

(f) Fix $t\in Q_1\setminus Q_2$ with leading coefficient $a$ and put
$\tau=t-1$.  For $u\in\Pfr^k$ one has
$(t(1+u))^3=t^3\bigl(1+L(u)+O(\Pfr^{2k})\bigr)$ with
$L=1+\Ad(t^{-1})+\Ad(t^{-2})$, and
$L(u)\equiv(\Ad(t^{-1})-1)^2u\equiv[[u,\tau],\tau]\pmod{\Pfr^{k+3}}$
(ring commutators), because $3u\in\Pfr^{k+4}$.  For $u=b\pi^3$ the
double commutator is $-aD\pi^5$ with $D=b\sigma(a)-a\sigma(b)$, which
vanishes iff $b\in\Fq_3a$.  Hence cubing maps
$W:=t\cdot(1+\Fq_3a\pi^3+\Pfr^4)$ into $t^3Q_6$.  Cubing is a local
homeomorphism near regular elements (its differential $L$ is
invertible, its determinant on $\mathfrak d/\mathfrak f$ having absolute
value $|3|_F\,R(t)^{-2}=3^{-4}$ for depth-one $t$), so the image of
$W$ is open in $t^3Q_6$, and its measure is
$\mu(W)\cdot3^{-4}=3^{-4}\cdot3^{-4}=3^{-8}=\mu(Q_6)$.  Therefore
cubing maps $W$ onto $t^3Q_6$, i.e.\ $t^3Q_6\subset G_3$.
(The finite verification in Appendix~\ref{W:app:scripts}, script~2,
confirms this statement in $Q/Q_7$.)

(g) $[K^3,K^3]\subset K^6$ and $(K^3)^3\subset K^7$, so $V$ is
elementary abelian of order $[Q_3:Q_6]=3^5$.  By (f) and the
Jacobian computed in (f), $\mu(G_3)=\tfrac89\cdot3^{-4}$, which is
one third of $\mu(Q_3\setminus Q_4)$; hence $|\bar G_3|=\tfrac13\cdot216=72$.
The explicit description is in Proposition~\ref{W:prop:G3explicit};
spanning and the symmetry $\bar G_3=-\bar G_3$ are read off from it.
\end{proof}

\subsection{Coordinates on $V$ and the set $\bar G_3$}
This subsection introduces coordinates on the finite group $V=Q_3/Q_6$ and describes the good classes of depth three explicitly in them. The later finite computations are carried out in these coordinates.

Write $\Pfr^3=\varpi\Ocal_{E_0}j\oplus\varpi^2\Ocal_{E_0}$ and
$\Pfr^6=\varpi^3\Ocal_D$.  Since $(1+X)(1+Y)\equiv1+X+Y$ modulo
$\Pfr^6$ for $X,Y\in\Pfr^3$, the map $1+X\mapsto X$ identifies
$K^3/K^6$ with the additive group $\Pfr^3/\Pfr^6$, and writing
\[
 X\equiv\varpi(A+B\varpi)j+\varpi^2C,\qquad A,B,C\in\Z_9=W(\Fq_9),
\]
gives coordinates $(\bar A,\bar B,\bar C)\in\Fq_9^3$ (levels $3,5,4$
respectively).  The scalars $U_F^{(2)}$ contribute exactly the
$\Fq_3$-line in the $C$-coordinate.  Hence
\begin{equation}\label{W:eq:Vcoords}
 V=Q_3/Q_6\simeq\Fq_9\oplus(\Fq_9/\Fq_3)\oplus\Fq_9\ni(\bar A,\ \bar C\bmod\Fq_3,\ \bar B),
\end{equation}
and we use the $\Fq_3$-basis $(A_{\mathrm{re}},A_{\mathrm{im}},C_{\mathrm{im}},B_{\mathrm{re}},B_{\mathrm{im}})$
in the tables.

\begin{proposition}[\Pctag\ The good depth-three classes]\label{W:prop:G3explicit}
With $\Nm(a)=a\bar a$ and $\Tr(a)=a+\bar a$ on $\Fq_9$,
\begin{equation}\label{W:eq:G3formula}
 \bar G_3=\Bigl\{\bigl(a,\ x_4,\ y+\Nm(a)\bigl(\Tr(a\bar y)+x_4^2-1\bigr)a\bigr)
 \ :\ a\in\Fq_9^\times,\ x_4\in\Fq_3 i,\ y\in\Fq_9\Bigr\},
\end{equation}
a set of $72$ elements forming $36$ lines through $0$ (each line is
$\bar G_3\cap\bar A_E$ for the ramified tori $E$ whose image in $V$ is
a fixed plane).  Explicitly, for $a\in\Fq_9^\times$ and $z\in\Fq_3\,i\bar a^{-1}$,
\[
 \bar G_3=\{(a,0,z-\Nm(a)a)\}\cup\{(a,\pm i,z+\Nm(a)a)\}.
\]
\end{proposition}
\begin{proof}
For a ramified torus $E=F(j')$ with $j'=c\pi+c_2\pi^2+c_3\pi^3+\cdots$
pure (so $\Tr(c_2)=0$, $j'^2=u'\varpi$ with $\bar u'=\Nm(c)$), the cubes
of depth-one elements $1+j'w$, $w\in\Ocal_E^\times$, are
$\equiv1+(d'w_0^3-\varpi^2w_0)j'$ modulo $\Pfr^6$ and scalars, where
$w\equiv w_0\in\Ocal_F^\times$.  Reading off coordinates gives the
vectors $\pm(\Nm(c)c,\Nm(c)c_2,\Nm(c)c_3+(\bar u_1'-1)c)$ with
$\bar u_1'=\Tr(c\bar c_3)+c_2^2$, which after the substitution
$a=\Nm(c)c$, $y=\Nm(c)c_3$ is \eqref{W:eq:G3formula}.  The map
$y\mapsto y+\Nm(a)\Tr(a\bar y)a$ has kernel $\Fq_3a$ and image
$\Fq_3\,i\bar a^{-1}$, giving the second description and the count
$8\cdot3\cdot3=72$.  The formula was verified by exact enumeration
of the $648$ depth-one classes modulo $K^4U_F^{(1)}$
(Appendix~\ref{W:app:scripts}, script~1).
\end{proof}

\subsection{The transferred trivial character on $\tU_1$}
This subsection computes the values of $L_D\mathbf1$ on the preimage $\tU_1$ of $U_1$. In particular it shows that the restriction is supported on $\tU_3$.

\begin{lemma}\label{W:lem:Phi}
Put $\Phi:=L_D\mathbf1|_{\tU_1}$ and, for a good $g\in U_1$, write
$g=zt^3$ with $z\in F^\times$ and $t\in K^1$ ($t=1$ if $g$ is central);
its \emph{star lifts} are the three elements $\zeta\tilde z\,t^*$,
$\zeta\in\mu_3$.  Then
\begin{equation}\label{W:eq:Phivalues}
 \Phi(\zeta\tilde z\,t^*)=\begin{cases}
 27\,\zeta\,\tth(\tilde z),& t\in F^\times K^2\ (\text{i.e. } g\in U_6),\\
 9\,\zeta\,\tth(\tilde z),& t\in F^\times K^1\setminus F^\times K^2\ (\text{i.e. } g\in F^\times G_3),
 \end{cases}
\end{equation}
and $\Phi$ vanishes at all lifts of non-good elements.  In particular
$\Phi$ is supported on $\tU_3$.
\end{lemma}
\begin{proof}
By \eqref{W:eq:LD}, $\Phi(\tilde z\,t^*)=3\tth(\tilde z)R_E(t)$ with
$E=F(t)$.  If $t\in F^\times K^2$ and $E=E_0$, then $t\in F^\times U^{(1)}_{E_0}$
and $R=9$ by \eqref{W:eq:unram-values}; if $E$ is ramified, $t\in F^\times U^{(2)}_E$,
which is the case $a\in\Ocal_F^\times$, $b\in\varpi\Ocal_F$ of
\eqref{W:eq:ramified-values}, so $R=9$.  If $t\in F^\times K^1\setminus F^\times K^2$,
then $E$ is ramified ($E_0\cap K^1\subset K^2$) and $t\in F^\times(U^{(1)}_E\setminus U^{(2)}_E)$,
the case $a,b$ units, so $R=3$.  Good elements of $U_1$ are
$F^\times\cdot\Cubes(K^1)$, whose images in $Q$ are $G_3\sqcup Q_6$ by
Lemma~\ref{W:lem:cube}(e).
\end{proof}

\section{Reduction to the pro-3 Sylow subgroup}\label{W:sec:reduction}
This section reduces Target~A to a statement about the group $\tU_1$. It passes to a finite quotient of $\tH$, treats the elements of $2$-power order directly, and computes the norm of the restricted function. Theorem~\ref{W:thm:reduction} combines these steps.

\subsection{A finite quotient}
This subsection replaces $\tH$ by a finite quotient $\Gamma$ on which $L_D\mathbf1$ is defined, so that the character theory of finite groups applies. It also shows that the toral restrictions of Section~\ref{W:sec:examplewild} descend to $\Gamma$.

Let $\phi:=L_D\mathbf1$.  By Chapter~\ref{ch:III} the row of $\phi$ in
\eqref{W:eq:matrixchars} is finite: $\phi=\sum_{\tau\in S}a_\tau\Theta_\tau$
with $S$ finite.  Each $\tau\in S$ is a finite-dimensional smooth
representation with central character $\tth$, so $\ker\tau$ is an open
normal subgroup of $\tH$ with $\ker\tau\cap\tZ=\ker\tth$.  Put
\begin{equation}\label{W:eq:N}
 N:=\tU_1\cap\bigcap_{\tau\in S}\ker\tau\cap N_6 ,
\end{equation}
where $N_6\subset\tU_6$ is the open normal subgroup introduced in
Definition~\ref{W:def:psi6} below (for the present section one may take
$N_6=\tU_m$ for any large $m$).  Then $N$ is open and normal in $\tH$,
$N\cap\tZ=\ker\tth$, and
\[
 \Gamma:=\tH/N\quad\text{is finite},\qquad
 T:=\tZ N/N\simeq\tZ/\ker\tth=\mu_3\ \text{is central},\qquad
 \Gamma_0:=\tU_1/N .
\]
$\Gamma_0$ is a normal $3$-subgroup of index $8$ with $\Gamma/\Gamma_0\simeq D_8$,
hence the unique Sylow $3$-subgroup; $\phi$ and $\Phi=\phi|_{\tU_1}$ are
class functions on $\Gamma$ and $\Gamma_0$.  For a maximal subfield
$E$ write $T_E:=\widetilde{E^\times}N/N\subset\Gamma$.

\begin{lemma}[Toral restrictions descend]\label{W:lem:toral-descent}
For every maximal subfield $E$, $\phi|_{T_E}$ is the character of a
representation of $T_E$.
\end{lemma}
\begin{proof}
By Proposition~\ref{W:prop:all-three-tori}, $\phi|_{\widetilde{E^\times}}$ is
the character of a (finite-dimensional, genuine) representation $W$.
Since $\phi(gm)=\phi(g)$ for $m\in N\cap\widetilde{E^\times}$, in
particular $\Theta_W(m)=\Theta_W(1)=\dim W$, so $m$ acts trivially on $W$.
\end{proof}

\subsection{Two-groups}
This subsection treats the elements of $2$-power order in $\Gamma$. It computes the restriction of $L_D\mathbf1$ to the relevant $2$-subgroup and describes the centralizers of $2$-elements. Both facts are used in the proof of the reduction theorem.

\begin{lemma}[The star section on $2$-groups]\label{W:lem:2groups}
Let $\hat P_0\subset\tH$ be the preimage of $P_0F^\times$ and
$\Gamma_{P_0}:=\hat P_0N/N$.  Then $\Gamma_{P_0}\simeq\mu_3\times D_8$,
the elements of $2$-power order form the unique complement
$\tilde P_0\simeq D_8$ to $\mu_3$, and the map
$u\mapsto\sigma(u):=(u^{1/3})^*$ ($u^{1/3}$ the cube root in the
group $P_0$ of order $8$) is the inverse of the projection
$\tilde P_0\to P_0$.  Consequently
\begin{equation}\label{W:eq:phiP0}
 \phi|_{\Gamma_{P_0}}=3\,\psi_{P}+3\sum_{\pi}(\dim\pi)\,\Theta_\pi ,
\end{equation}
where $\psi_P$ is the genuine linear character trivial on $\tilde P_0$ and
$\pi$ runs over the genuine irreducible representations of $\Gamma_{P_0}$;
in particular $\phi|_{\Gamma_{P_0}}$ is a character.
\end{lemma}
\begin{proof}
$\hat P_0\cap\tU_1=\tZ$ because $P_0\cap Q=1$, so $\hat P_0\cap N=\ker\tth$
and $\Gamma_{P_0}$ is a central extension of $D_8$ by $T=\mu_3$.
Since $|D_8|$ is prime to $3$, $H^2(D_8,\mu_3)=H^1(D_8,\mu_3)=0$: the
extension splits and the complement is unique; it consists of the
elements of $2$-power order.  On each cyclic subgroup of $P_0$ the star
map is a homomorphism (it is a homomorphism on every covered torus, as
$(\tilde x\tilde y)^3=\tilde x^3\tilde y^3[\tilde y,\tilde x]^3$ with
$[\tilde y,\tilde x]\in\mu_3$), so $\sigma(u)^8=1$ and $\sigma(u)\in\tilde P_0$
is the unique element of the complement over $u$.  Every nontrivial
$u\in P_0$ is of order $2$ or $4$ in a quadratic torus and has
$R_E(u)=1$ by \eqref{W:eq:unram-values}--\eqref{W:eq:ramified-values}
(the cases $b$ a unit, resp.\ $a\in\varpi\Ocal_F$), so by \eqref{W:eq:LD}
$\phi(\zeta\tilde z\sigma(u))=3\zeta\tth(\tilde z)$ for $u\ne1$ and
$=27\zeta\tth(\tilde z)$ for $u=1$.  That is $3\psi_P+24\delta$, with
$\delta$ the genuine delta function at $\mu_3$, and
$8\delta=\sum_\pi(\dim\pi)\Theta_\pi$.
\end{proof}

\begin{lemma}[Centralizers of $2$-elements]\label{W:lem:centralizers}
Let $y\in\Gamma$ be an element of $2$-power order whose image in
$\Gamma/\Gamma_0\simeq D_8$ is nontrivial, $\eta\in D^\times$ a
representative of that image in $D^\times/F^\times$ and $E_y=F(\eta)$.
(For $\rho=\mathbf1$ every nontrivial $2$-element qualifies, $T$ being
$\mu_3$; in Theorem~\ref{W:thm:reduction-all} the central $\mu_2\subset T$
is excluded, as its centralizer is all of $\Gamma_0$.)
\begin{enumerate}[label=(\alph*)]
\item $C_{\Gamma_0}(y)\subset T_{E_y}$.
\item If $B_2\subset\Gamma$ is a $2$-subgroup whose image in
$\Gamma/\Gamma_0\simeq D_8$ is noncyclic, then $C_{\Gamma_0}(B_2)=T$.
\end{enumerate}
\end{lemma}
\begin{proof}
(a) Replacing $y$ by $y^2$ if necessary (which does not change $E_y$ and
enlarges the centralizer) we may assume $\eta$ has order $2$ modulo
$F^\times$, so $s:=\Ad(\eta)$ is an involutive automorphism of the
pro-$3$ group $Q$.  For $x\in Q$ put $b:=(x^{-1}s(x))^{-1/2}$ (square
roots are unique in a pro-$3$ group) and $a:=xb^{-1}$; one checks
$s(a)=a$, $s(b)=b^{-1}$ and $x^{-1}s(x)=b^{-2}$.  If
$[y,x]\in N$ then $b^{-2}\in\bar N$ (the image of $N$ in $Q$, an
$s$-stable closed normal subgroup), hence $b\in\bar N$ and
$x\in Q^s\bar N$.  Finally $Q^s=\{x:\eta x\eta^{-1}\in\pm xF^\times\}$;
the sign $-$ would give a nontrivial $2$-torsion element of $Q$, so
$Q^s=(E_y^\times\cap F^\times K^1)/F^\times$, and $x$ lifts into
$\widetilde{E_y^\times}N\tZ$.

(b) The (noncyclic) image of $B_2$ in $D_8$ contains the central involution
$\bar a^2$ (the class of $i$) and a reflection.  By (a) applied to the
class of $i$, an element $b\in C_{\Gamma_0}(B_2)$ is the image of
$\tilde zu$ with $u$ a lift of $u_0\in U^{(1)}_{E_0}$.  A reflection $\eta_2$
acts on $E_0$ by $\sigma$ and commutes with $b$ in $\Gamma$, so
$\sigma(u_0)u_0^{-1}\in\bar N$, i.e.\ $u_0^{-2}\in\bar N F^\times$,
hence $u_0\in\bar N$ and $b\in T$.
\end{proof}

\subsection{The norm of $\Phi$}
This subsection computes the norm of $L_D\mathbf1$ and of its restrictions to the subgroups $\tU_m$. These norms are used in the reduction theorem and in the later sections.

\begin{lemma}\label{W:lem:norm}
With respect to the probability Haar measures,
$\langle\phi,\phi\rangle_{\tH}=1$ and, for $m\ge1$,
\[
 \langle\Phi|_{\tU_m},\Phi|_{\tU_m}\rangle_{\tU_m}=[Q_{m'}:Q_m],\qquad
 m'=\begin{cases}1,&m\le3,\\2,&4\le m\le6,\\m-4,&m\ge7.\end{cases}
\]
In particular the norms for $m=1,3,4,6$ are $1,\,27,\,27,\,729$.
\end{lemma}
\begin{proof}
Weyl integration on the compact group $H/Z$ reads
$\int_{H/Z}f\,d\mu=\sum_E\tfrac{c_E}{2}\int_{E^\times/F^\times}\Delta(t)^2f(t)\,dt$
for class functions $f$, with constants $c_E$ fixed by $f\equiv1$.
On each torus the cube map $s\mapsto s^3$ is an injective endomorphism
of $E^\times/F^\times$ with image the good elements, of index $9$; it
pushes Haar measure forward to $\tfrac19$ of Haar measure on the image.
Substituting $t=s^3$ and using $|\phi|^2(s^3)=9R(s)^2$ and
$R(s)=\Delta(s)/\Delta(s^3)$ gives
$\int_{E^\times/F^\times}\Delta(t)^2|\phi|^2(t)\mathbf 1_{U}(t)\,dt
 =\int\Delta(s)^2\mathbf 1_{U}(s^3)\,ds$.
For $U=H$ this is unitarity of the row.  For $U=U_m$, note that
$s^3\in U_m$ iff $s\in U_{m'}$ by Lemma~\ref{W:lem:cube}(c)--(e) and (a),
so $\int_{U_m/Z}|\phi|^2d\mu=\mu(U_{m'}/Z)$, and the normalized inner
product on $\tU_m$ is $\mu(U_{m'})/\mu(U_m)=[Q_{m'}:Q_m]$.
\end{proof}

\subsection{The reduction theorem}
This subsection states and proves the reduction. Target~A holds if and only if the restriction of $L_D\mathbf1$ to $\tU_1$ is a virtual character.

\begin{theorem}[Reduction to the Sylow $3$-subgroup]\label{W:thm:reduction}
The following are equivalent.
\begin{enumerate}[label=(\roman*)]
\item (Target A) $L_D\mathbf1$ is the character of a genuine irreducible
representation of $\tH$ with central character $\tth$.
\item $L_D\mathbf1$ is a virtual character of $\tH$.
\item $\Phi=L_D\mathbf1|_{\tU_1}$ is a virtual character of $\tU_1$.
\item $\Phi$ is the character of an irreducible representation of
$\Gamma_0=\tU_1/N$, necessarily of dimension $27$.
\end{enumerate}
Moreover, in that case the representation in (i) has dimension $27$,
its restriction to $\tU_1$ is irreducible, and it is unique.
\end{theorem}
\begin{proof}
(i)$\Rightarrow$(ii)$\Rightarrow$(iii) are clear, and (iii)$\Leftrightarrow$(iv)
by Lemma~\ref{W:lem:norm} ($\langle\Phi,\Phi\rangle=1$, $\Phi(1)=27>0$).
Assume (iii).  By Brauer's characterization of characters it suffices
to show that $\phi|_B$ is a virtual character for every elementary
subgroup $B=B_2\times B_3\subset\Gamma$ ($B_\ell$ an $\ell$-group, one
of them cyclic).  If $B_2=1$, then $B\subset\Gamma_0$ and $\phi|_B=\Phi|_B$.
If $B_2=\langle y\rangle$ is cyclic and nontrivial, then $y\in T_{E_y}$
and $B_3\subset C_{\Gamma_0}(y)\subset T_{E_y}$ by Lemma~\ref{W:lem:centralizers}(a),
so $B\subset T_{E_y}$ and Lemma~\ref{W:lem:toral-descent} applies.  If
$B_2$ is noncyclic, then $B_3\subset T$ by Lemma~\ref{W:lem:centralizers}(b),
so after conjugation $B\subset\Gamma_{P_0}$ and Lemma~\ref{W:lem:2groups}
applies.  Hence $\phi$ is a virtual character of $\Gamma$; by unitarity
(Lemma~\ref{W:lem:norm}) and $\phi(1)=27>0$ it is an irreducible
character, which is (i).  The restriction of that representation to
$\tU_1$ has character $\Phi$, of norm one, hence is irreducible; a
character determines the representation.
\end{proof}

\begin{remark}[Independence of the choice of $\tth$]\label{W:rem:theta-independence}
Two compatible choices $\tth,\tth'$ for $\rho=\mathbf1$ differ by a
character $\theta_0$ of $F^\times/F^{\times3}$.  Twisting a
representation of $\tH$ by $(\theta_0\circ\Nrd\circ p)^{-1}$ changes
its central character by $\theta_0(z)^{-2}=\theta_0(z)$ and multiplies
its character at $\tilde zt^*$ by $\theta_0(\Nrd(zt^3))^{-1}=\theta_0(z)$;
this is exactly the relation between $L_D\mathbf1$ computed with
$\tth$ and with $\tth\theta_0$.  So Target~A does not depend on the
choice.
\end{remark}

\section{The finite model and the defect function}\label{W:sec:finite}
This section studies the restriction of $L_D\mathbf1$ to $\tU_3$ in a finite model. It introduces Hypothesis (H1) and a defect function, and it shows that Target~A is equivalent to (H1).

Throughout this section we assume the following hypothesis, which is
necessary for Target~A (Lemma~\ref{W:lem:H1necessary}) and will be
proved in Section~\ref{W:sec:model}.

\begin{definition}[Hypothesis (H1)]\label{W:def:psi6}\index{hypothesis h1@Hypothesis (H1)}
Let $\psi_6:\tU_6\to\mu_3$ be the function
$\psi_6(\zeta\tilde z\,t^*):=\zeta\tth(\tilde z)$ on star lifts (all
elements of $\tU_6$ are star lifts, Lemma~\ref{W:lem:cube}(d)).
\textbf{Hypothesis (H1)}: $\psi_6$ is a (continuous) character of
$\tU_6$.  Under (H1) put $N_6:=\ker\psi_6$, an open subgroup of
$\tU_6$, normal in $\tH$ because the star section is conjugation
equivariant and $\tth$ is central.
\end{definition}

\begin{example}[What (H1) says over one maximal subfield]\label{W:ex:H1torus}
Let $E\subset D$ be a maximal subfield, and consider the elements of $\tU_6$ that lie over $E^\times$. The star map is a homomorphism on each covered torus (\S\ref{W:subsec:coverbackground}), and $\tZ$ is central. So for two such elements $\zeta\tilde z\,t^*$ and $\zeta'\tilde z'\,t'^*$ with $t,t'\in E^\times$, the product is $\zeta\zeta'\,\tilde z\tilde z'\,(tt')^*$, and
\[
 \psi_6\bigl(\zeta\tilde z\,t^*\cdot\zeta'\tilde z'\,t'^*\bigr)
 =\zeta\zeta'\,\tth(\tilde z\tilde z')
 =\psi_6(\zeta\tilde z\,t^*)\,\psi_6(\zeta'\tilde z'\,t'^*).
\]
Thus $\psi_6$ is multiplicative on the part of $\tU_6$ over any single maximal subfield, and this needs no hypothesis. The content of (H1) is the same identity for star lifts $t^*$ and $t'^*$ with $t$ and $t'$ in different maximal subfields.
\end{example}


\begin{lemma}\label{W:lem:H1necessary}
If $L_D\mathbf1=\Theta_\tau$ with $\tau$ irreducible, then $\tau|_{\tU_6}=27\,\psi_6$;
in particular (H1) holds.
\end{lemma}
\begin{proof}
By Lemma~\ref{W:lem:Phi}, $|\Theta_\tau|=27=\dim\tau$ on $\tU_6$, so $\tU_6$
acts by scalars, by $\psi_6$.
\end{proof}

Under (H1) define
\[
 \bar\Gamma:=\tU_1/N_6\ (\text{order }3^9),\qquad
 \tV:=\tU_3/N_6,\qquad
 \mu_3\simeq\tU_6/N_6\subset\tV\to V=Q_3/Q_6\to0 .
\]
$\tV$ is a central extension of the elementary abelian group $V$ by
$\mu_3$; let $\omega:V\times V\to\mu_3$ be its commutator form.  By
Lemma~\ref{W:lem:Phi}, $\Phi$ factors through $\bar\Gamma$, with
$\bar\Gamma_6:=\tU_6/N_6=\mu_3$ carrying $\Phi=27\,\mathrm{id}$.  The
elements $\bar\sigma(g)$ ($g\in G_3$) are the images of the star lifts;
by Lemma~\ref{W:lem:cube}(f) and the next lemma they depend only on the
class of $g$ in $\bar G_3\subset V$.

\begin{lemma}[Three consequences of (H1)]\label{W:lem:H1consequences}
Assume (H1).
\begin{enumerate}[label=(\alph*)]
\item For $a\in H$ and $b\in U_6$, $\psi_6([\tilde a,\tilde b])=1$.  Hence for
$i+j\ge6$ the map $\kappa_{i,j}(a,b):=\psi_6([\tilde a,\tilde b])$ on
$K^i\times K^j$ is bimultiplicative, alternating, conjugation
invariant, and trivial if $a$ or $b$ lies in $U_6$; $\kappa_{1,5}$
factors through $(Q_1/Q_2)\times(Q_5/Q_6)\simeq\Fq_9\times\Fq_9$.
\item $\kappa_{1,5}(a,b)=\zeta^{\beta\Tr(a\bar b)}$ for a fixed sign
$\beta\in\{\pm1\}$.
\item (H3) $\omega=0$: $\tV$ is abelian, hence $\tV\simeq\mu_3\times V$.
\item (H2) For $g\in G_3$ and $u\in Q_6$, $\bar\sigma(gu)=\bar\sigma(g)$.
\end{enumerate}
\end{lemma}
\begin{proof}
(a) $\ker\psi_6$ is normal in $\tH$ and $\tU_6/\ker\psi_6$ is central
in $\tH/\ker\psi_6$ (as $\psi_6$ is conjugation invariant), so
$[\tilde a,\tilde b]\in\ker\psi_6$ for $b\in U_6$.  Bimultiplicativity
follows from $[xy,z]=x[y,z]x^{-1}[x,z]$.  A character of $Q$ trivial on
cubes and commutators factors through $Q/Q_2$, since
$[Q,Q]Q^3\supseteq Q_2$ (level by level: at even levels the commutators
$[1+a\pi,1+b\pi^{k-1}]$ produce the trace-zero line and $U_F$ the rest;
at odd levels they produce all of $\Fq_9$).

(b) Write $\kappa_{1,5}(a,b)=\zeta^{\Tr(\alpha ab)+\Tr(\beta a\bar b)}$ with
$\alpha,\beta\in\Fq_9$ (these span the bilinear forms).  Conjugation
by $\xi\in\Ocal_{E_0}^\times$ multiplies both $a$ and $b$ by
$\lambda=\xi/\bar\xi\in\mu_4$, forcing $\alpha=0$; conjugation by $j$
applies Frobenius to both, forcing $\beta\in\Fq_3$.  It remains to see
$\beta\ne0$, and for this we use the ramified torus $E=F(j)$,
$j^2=\varpi$, explicitly.  Its unit filtration is $U_E^{(k)}$ with
$U_E^{(2l)}/U_E^{(2l+1)}$ represented by $1+a\varpi^l$ ($a\in\Fq_3$, scalars)
and $U_E^{(2l+1)}/U_E^{(2l+2)}$ by $1+a\varpi^lj$ ($a\in\Fq_3$), of
$D$-depth $2l+1$.  For $a\in\Ocal_F^\times$,
$(1+aj)^3=1+a^3\varpi j+3a^2\varpi+3aj$, and $v_E(3)=4$, so the cube
has leading term $a^3\varpi j$ at $E$-depth $3$ and its next
$j$-component $3aj$ at $E$-depth $5$ is determined by $a$; cubes of
$U_E^{(3)}$ lie in $U_E^{(7)}$.  Hence modulo $F^\times E^{\times3}$ the
classes of $U_E^{(3)}$ and $U_E^{(k)}$, $k\ge6$, are trivial, and
$E^\times/F^\times E^{\times3}$, which has order $9$ (\cref{I:prop:good} and \eqref{W:eq:ranktwo-comm}), is the $\Fq_3$-vector space
with basis the classes of $1+j$ ($D$-depth one, leading coefficient
$1$) and $1+\varpi^2j$ ($D$-depth five, leading coefficient $1$).  The
commutator pairing \eqref{W:eq:ranktwo-comm} is alternating and
nondegenerate on this two-dimensional space (\cref{I:prop:good}),
so $[\,\widetilde{1+j},\widetilde{1+\varpi^2j}\,]\ne1$; by (a) this
commutator is $\kappa_{1,5}(1,1)=\zeta^{\beta\Tr(1)}=\zeta^{2\beta}$,
whence $\beta\ne0$.  (Consistently, $N_{E/F}(1+aj)=1-a^2\varpi\in U_F^{(1)}$
and $N_{E/F}(1+a\varpi^2j)\in U_F^{(5)}$ pair trivially in $F$, so the
commutator is the Hilbert pairing $(1+j,1+\varpi^2j)_E^{-1}$ of $E$ at the
top level $1+5=pe_E/(p-1)=6$, which is nondegenerate \cite[Ch.~IV]{FV02}.)

(c) Let $t$ have depth one and $y\in K^4$, $c:=[\tilde t,\tilde y]\in\tilde K^5$,
$d:=[c^{-1},\tilde t]\in\tilde K^6$, $e:=[d^{-1},\tilde t]$.  Then
$[\tilde t^3,\tilde y]=c^{t^2}c^{t}c$ and, expanding with $c^t=cd$,
$c^{t^2}=cd^2e$, one gets $[\tilde t^3,\tilde y]=c^3\cdot Y$ with $Y$ a
product of $d^3$, $e$ and commutators each involving $d$ or $e$.  As
$c^3$ is a star lift, $d,e\in\tU_6$ and (a) holds, $\psi_6([\tilde t^3,\tilde y])=1$.
Thus $\omega(\bar G_3,V_4)=0$ with $V_4=Q_4/Q_6$, hence $\omega(V,V_4)=0$
since $\bar G_3$ spans $V$ (Lemma~\ref{W:lem:cube}(g)).  So $\omega$ is
induced from an alternating form on $V/V_4=Q_3/Q_4\simeq\Fq_9$, a
two-dimensional $\Fq_3$-space, on which conjugation by $j$ acts by
Frobenius, of determinant $-1$; invariance forces $\omega=0$.  So $\tV$
is abelian.  It also has exponent $3$: for $v\in K^3$ the cube of any
lift is the star lift of $v^3\in K^9\subset U_6$, on which $\psi_6=1$,
and for $z\in F^\times$ the cube of $\tilde z$ is the star lift of
$z^3$, on which $\psi_6=\theta_0(z)^3=1$; so every element of
$\tV=\tU_3/N_6$ cubes into $N_6$.  An abelian group of exponent $3$ is
an $\Fq_3$-vector space, so the subgroup $\mu_3$ has a complement.
(Abelianness alone would not suffice: $C_9\to C_3$ is a nonsplit abelian
extension with genuine characters.)

(d) Let $g=t^3$, $t$ of depth one with leading coefficient $a$, and
$gu=t'^3$ with $t'=t(1+w)$, $w\in\Fq_3a\pi^3+\Pfr^4$
(proof of Lemma~\ref{W:lem:cube}(f)).  In $\bar\Gamma$ put $x=\tilde t$,
$y=\widetilde{1+w}$, $c=[y,x]_L=yxy^{-1}x^{-1}$.  The ring computation
in Lemma~\ref{W:lem:cube}(f) shows $[1+w,t]\in K^5$ with level-five
coefficient in $\Fq_3\,ia$, so $c\in\bar\Gamma_5$.  Since $\bar\Gamma_2$
has exponent $3$ (cubes of $\tU_2$ are star lifts), $c^3=1$;
$[c,y]\in\bar\Gamma_8=\mu_3$ equals $\omega$-type pairing $\kappa_{5,3}$,
trivial by (c); and $d:=c^{-1}xcx^{-1}\in\bar\Gamma_6=\mu_3$ is central.
A direct expansion then gives $(xy)^3=d^{-1}x^3y^3$, and $y^3=1$ in
$\bar\Gamma$ because $(1+w)^3\in Q_9$.  Hence
$\bar\sigma(gu)\bar\sigma(g)^{-1}=d^{-1}=\kappa_{1,5}(a,\,s\,ia)^{\pm1}$
for some $s\in\Fq_3$, and $\Tr(a\,\overline{s\,i\,a})=-s\Nm(a)\Tr(i)=0$.
\end{proof}

\subsection{The $27$ characters and the defect function}
This subsection describes the genuine characters of the finite group $\tV$ and defines the defect function $e$. It then records what Target~A forces on $e$.

Assume (H1).  Fix a complement $V_0$ to $\mu_3$ in $\tV$; genuine
characters of $\tV$ are then $\mu_h(\zeta,v)=\zeta\cdot\zeta^{h(v)}$,
$h\in V^\vee$.  The star section gives a canonical lift
$\bar\sigma(v)\in\tV$ for each $v\in\bar G_3$, and we define the
\emph{defect function} $e:\bar G_3\to\Fq_3$ by
$\bar\sigma(v)=(\zeta^{e(v)},v)$; it satisfies $e(-v)=-e(v)$
(star is a homomorphism on each torus line).  The group
$\Heis:=Q_1/Q_3$ (extraspecial of order $27$ and exponent $3$) acts on
$\tV$ by conjugation, fixing $\mu_3$; relative to $V_0$ the action is
$(\zeta,v)\mapsto(\zeta\zeta^{\beta_s(v)},sv)$ with $\beta_s\in V^\vee$,
and equivariance of the star section gives
\begin{equation}\label{W:eq:cocycle}
 e(sv)-e(v)=\beta_s(v)\qquad(v\in\bar G_3,\ s\in\Heis).
\end{equation}

\begin{proposition}[What Target A forces]\label{W:prop:forces}
Suppose $L_D\mathbf1=\Theta_\tau$.  Then, with $\Lambda\subset V^\vee$ the
set of $h$ with $\mu_h$ a constituent of $\tau|_{\tV}$:
\begin{enumerate}[label=(\alph*)]
\item $\tau|_{\tV}$ is multiplicity free, $|\Lambda|=27$, and $\Heis$ acts
simply transitively on $\Lambda$ (by $h\mapsto(h-\beta_s)\circ s^{-1}$);
\item for $v\in V\setminus(\bar G_3\cup\{0\})$ the values $h(v)$, $h\in\Lambda$,
are equidistributed in $\Fq_3$;
\item for $v\in\bar G_3$ the value $-e(v)$ is taken $15$ times and the
other two values $6$ times each;
\item equivalently, for every $h\in V^\vee$ the number
$N_e(h):=\#\{\text{lines }\ell\subset\bar G_3\cup\{0\}:\ h|_\ell=-e|_\ell\}$
is $20$ if $h\in\Lambda$ and $11$ otherwise.
\end{enumerate}
Conversely, if (H1) holds and $e$ satisfies (d) for some $27$-set
$\Lambda$ on which $\Heis$ acts simply transitively, then
$\tau_1:=\Ind_{\tU_3}^{\tU_1}\lambda$, $\lambda=\mu_h$ for any $h\in\Lambda$,
is irreducible with character $\Phi$, and Target~A holds by
Theorem~\ref{W:thm:reduction}.
\end{proposition}
\begin{proof}
By Lemma~\ref{W:lem:norm}, $\langle\Phi|_{\tU_3},\Phi|_{\tU_3}\rangle=27=\Phi(1)$,
which forces multiplicity one with $27$ constituents; Clifford theory
for the normal subgroup $\tU_3$ of $\tU_1$, with $[\tU_1:\tU_3]=27$,
gives (a).  Evaluating $\sum_{h\in\Lambda}\mu_h=\Phi|_{\tV}$ at
$(1,v)$ and at $\bar\sigma(v)$ gives (b) and (c), using
\eqref{W:eq:Phivalues} and (H2); (d) is the Fourier dual statement
(each line $\{v,-v\}$ contributes $2$ or $-1$ to
$\sum_{v\in\bar G_3}\zeta^{h(v)+e(v)}=3N_e(h)-36\in\{24,-3\}$).  The
converse is Mackey irreducibility and the Frobenius character formula,
whose Fourier transform is (b)--(c), together with the vanishing of
$\Phi$ and of the induced character off $\tU_3$.
\end{proof}

\subsection{Symmetry pins down the defect function}
This subsection uses the conjugation action of $D^\times$ to cut the possible defect functions down to a two-parameter family, and then determines the defect function. The conclusion is Theorem~\ref{W:thm:H1equiv}.

The group $D^\times$ acts on everything by conjugation through
$G_{\mathrm{sym}}:=D^\times/F^\times K^3$, of order $216$, an extension of
$D_8=\langle\Ad(1+i),\Ad(j)\rangle$ by $\Heis$.  Its action on $V$ is
recorded in Appendix~\ref{W:app:tables}; one checks $V^{D_8}=0$ and
$(V^\vee)^{D_8}=0$.  Since $|D_8|$ is prime to $3$, there is a
$D_8$-stable complement $V_0$, and it is unique.  With respect to it,
$e$ is $D_8$-invariant.  The $D_8$-orbits on $\bar G_3$ have sizes
$4,4,8,8,8,8,8,8,8,8$; six of them contain $-v$ together with $v$,
forcing $e=0$ there, and the other four form two pairs $\{O,-O\}$.
Hence $e$ is determined by two parameters in $\Fq_3$.

\begin{proposition}[\Pctag\ Uniqueness of the design]\label{W:prop:design}
Among the $3^{10}$ $D_8$-invariant functions $e$ on $\bar G_3$,
exactly two are odd and satisfy condition (d) of
Proposition~\ref{W:prop:forces}, namely $e_*$ and $-e_*$, where $e_*$ is
tabulated in Appendix~\ref{W:app:tables}; it vanishes on $20$ of the $36$
lines and takes the values $1$ and $2$ on $8$ lines each.  Among the
nine odd $D_8$-invariant $e$, exactly three admit linear
$\beta_s$ as in \eqref{W:eq:cocycle} for the generators
$s=\Ad(1+j),\Ad(1+ij)$ of $\Heis$: $e=0$ and $e=\pm e_*$.  For
$e=\pm e_*$ the resulting $\Heis$-action on $V^\vee$ is simply
transitive on $\Lambda_{\pm}:=\{h:N_{\pm e_*}(h)=20\}$, $0\in\Lambda_\pm$,
and (b)--(c) hold.  For $e=0$ one has $\beta_s=0$ for all $s$.
\end{proposition}
\begin{proof}
Exact enumeration over $\Fq_3$ (Appendix~\ref{W:app:scripts}, scripts
3, 4, 6): the $D_8$-orbits, the counts $N_e(h)$ for all $243$ covectors
and all candidates, the linear extension problem \eqref{W:eq:cocycle},
the orbit of $0$ under the affine action, and the Fourier identities.
The distribution of $N_0(h)$ is $\{9^{80},12^{90},15^{72},36^{1}\}$,
that of $N_{\pm e_*}(h)$ is $\{11^{216},20^{27}\}$.
\end{proof}

\begin{theorem}[\Pctag\ Target A is equivalent to (H1)]\label{W:thm:H1equiv}
$L_D\mathbf1$ is a genuine irreducible character if and only if (H1)
holds.  In that case the defect function is $\pm e_*$, and
$\tau|_{\tU_1}\simeq\Ind_{\tU_3}^{\tU_1}\lambda_0$ where $\lambda_0$ is the
genuine character of $\tU_3$ trivial on $N_6$ and on the unique
$D_8$-stable complement $V_0$; explicitly
\begin{equation}\label{W:eq:lambda0}
 \lambda_0(y)=\Bigl(\prod_{d\in D_8}\tilde d\,y\,\tilde d^{-1}\Bigr)^{-1}
 \in\mu_3\qquad(y\in\tU_3),
\end{equation}
the product taken in the abelian group $\tV$ (it lies in $\tV^{D_8}=\mu_3$).
\end{theorem}
\begin{proof}
Necessity is Lemma~\ref{W:lem:H1necessary}.  Assume (H1).  By
Lemma~\ref{W:lem:H1consequences}, (H2) and (H3) hold, so the defect
function $e$ is defined, is odd and $D_8$-invariant for the
$D_8$-stable complement, and satisfies \eqref{W:eq:cocycle}; by
Proposition~\ref{W:prop:design}, $e\in\{0,e_*,-e_*\}$.  If $e=0$, then
$\beta_s=0$ for $s=\Ad(1+j)$; but for the depth-five vector
$v=e_{B_{\mathrm{re}}}$ of the torus $F(j)$, fixed by $s$,
$\zeta^{\beta_s(v)}$ is the commutator of lifts, i.e.\
$\kappa_{1,5}(1,1)=\zeta^{2\beta}\ne1$ by Lemma~\ref{W:lem:H1consequences}(b).
Hence $e=\pm e_*$, and the converse part of Proposition~\ref{W:prop:forces}
applies with $\Lambda=\Lambda_\pm\ni0$; the character $\mu_0$ is
$\lambda_0$, and \eqref{W:eq:lambda0} is the formula for the projection
$\tV=\mu_3\times V_0\to\mu_3$, since $\prod_{d}d\cdot x=\lambda_0(x)^8=\lambda_0(x)^{-1}$
for $x\in\tV$.  Uniqueness of $V_0$: complements form a torsor under
$V^\vee$, and $(V^\vee)^{D_8}=0$.
\end{proof}

\begin{remark}
The two signs correspond to the two identifications of $\mu_3$ with
$\Z/3$ (equivalently, to replacing the cover by its opposite); the
statement of Theorem~\ref{W:thm:H1equiv} is symmetric.  The
$\Heis$-cocycle $\beta$ attached to $e_*$ restricted to
$V_5=Q_5/Q_6$ reproduces $\kappa_{1,5}$ with $\beta=-1$, an independent
consistency check between the design computation and the torus data.
\end{remark}

\section{The cover in the Kazhdan--Patterson model}\label{W:sec:model}
It remains to prove (H1).  We do this in an explicit model of the
cover.  By \cref{I:lem:center} and \cref{I:prop:restriction}(a), the cover of
$D^\times$ is the restriction, along
\begin{equation}\label{W:eq:iota}
 \iota:D^\times\hookrightarrow\mathrm{GL}_2(E_0),\qquad
 x+yj\longmapsto\begin{pmatrix}x&\varpi y\\\sigma(y)&\sigma(x)\end{pmatrix},
\end{equation}
of the parameter-$c_v$ Kazhdan--Patterson cover of $\mathrm{GL}_2(E_0)$,
whose cocycle is $\sigma_{E_0}(g,h)\,(\det g,\det h)_{E_0}^{c_v}$ with
$\sigma_{E_0}$ the parameter-$0$ cocycle.  We use the Kubota-type
representative
\begin{equation}\label{W:eq:kubota}
 \sigma_{E_0}(g,h)=\Bigl(\frac{x(gh)}{x(g)},\ \frac{x(gh)}{x(h)\det g}\Bigr)_{E_0},
 \qquad x\begin{pmatrix}a&b\\c&d\end{pmatrix}=\begin{cases}c,&c\ne0,\\ d,&c=0,\end{cases}
\end{equation}
and write $\beta(g,h)\in\Z/3$ for the exponent of
$\sigma_{E_0}(g,h)(\det g,\det h)_{E_0}^{2}$ (before the relabeling
$\zeta\mapsto\zeta^{r'}$, which multiplies all exponents by $2$ and is
irrelevant for everything below).  Elements of the cover are pairs
$(g,e)$ with $(g,e)(h,f)=(gh,e+f+\beta(g,h))$.

\begin{lemma}[\Etag\ The model is the cover]\label{W:lem:model}
\begin{enumerate}[label=(\alph*)]
\item \eqref{W:eq:kubota} is a $2$-cocycle on $\mathrm{GL}_2(E_0)$ whose
restriction to the diagonal torus is $(a_1,b_2)_{E_0}$.
\item On scalars the twisted cocycle is trivial, so $z\mapsto(z,0)$ is a
homomorphic section over $F^\times$; and for $z\in F^\times$, $g\in\mathrm{GL}_2(E_0)$,
$\beta(z,g)=(z,\,x(g)\det g)_{E_0}$.
\item For $x,y\in E_0^\times$, the commutator of $(\iota x,0)$ and
$(\iota y,0)$ has exponent $(\Nm x,\Nm y)_F+(x,y)_{E_0}$, in accordance
with \cref{I:prop:comm} (before relabeling, $c=1$).
\item The extension so defined is isomorphic, as a central extension of
$D^\times$ by $\mu_3$, to the cover of Chapter~\ref{ch:I}; any two such isomorphisms
differ by $\chi\circ\Nrd$ for a character $\chi$ of $F^\times$ with
values in $\mu_3$, which does not affect star lifts, good elements, or
$L_D\mathbf1$.
\end{enumerate}
\end{lemma}
\begin{proof}
(a) is Kubota's cocycle for $\mathrm{GL}_2$: for $\mathrm{SL}_2$ it is
\cite[Thm.~2]{Kub67} and \cite[Ch.~1]{Kub69}, and the extension to $\mathrm{GL}_2$ by
the factor $\det g$ in the second entry is the case $r=2$, $c=0$ of the
Kazhdan--Patterson cocycle \cite[\S0.1]{KP84}, with the $n$-th power Hilbert
symbol for the $n$-fold cover (the same formula serves in
Section~\ref{W:sec:targetC}); it is the one external input of this section
and is labeled \Etag.  (The twist $(\det g,\det h)^{c}$ is
bimultiplicative, hence a cocycle, so the twisted formula is a cocycle
as well.)  As a safeguard against transcription errors, the cocycle
identity was also checked on $250$ random triples and the diagonal
restriction on $100$ random pairs by exact computation (script~7); this
check is corroboration, not part of the proof.  (b) With $x(zI)=z$, $x(zg)=zx(g)$, $\det(zI)=z^2$ one gets
$\sigma(zI,g)=(x(g),z^{-1})=(z,x(g))$ and the twist contributes
$(z^2,\det g)^2=(z,\det g)$; for $g=wI$ the two terms cancel, since
$(z,w)_{E_0}=(z,w)_F^2$.  (c) is a direct evaluation; it was also
checked on $200$ random pairs.  (d) Steps~1--3 of the proof of \cref{I:prop:restriction}(b), apply verbatim to two cocycles on $\mathrm{GL}_2(E_0)$
with the same diagonal restriction; the last assertion is \cref{I:rem:D2tame} (the star map $x^*=\tilde x^{\,3}$ absorbs $\chi\circ\det$).
\end{proof}

\paragraph{The Hilbert symbols.}
The cubic Hilbert symbols of $F$ and of $E_0$ enter only through their
classes.  By Moore's theorem (cited in \cref{I:prop:restriction}), a
continuous Steinberg symbol $F^\times\times F^\times\to\mu_3$ is a power of
the Hilbert symbol.  We therefore computed the alternating bilinear
forms on $F^\times/F^{\times3}\simeq\Fq_3^4$ and $E_0^\times/E_0^{\times3}\simeq\Fq_3^6$
annihilated by all Steinberg relations $(a,1-a)=1$ (exact linear
algebra over $\Fq_3$ on several thousand relations, script~7); in both
cases the solution space is one-dimensional and nondegenerate, so each
symbol is determined up to the common sign, and the relative sign is
fixed by norm compatibility $(a,b)_{E_0}=(a,\Nm b)_F$, which the
computed forms satisfy.  Classes modulo cubes are computed from the
filtrations $U^{(k)}$ by explicit discrete logarithms.  The hand
arguments below use only the following standard facts about the cubic
symbol of $E_0$, for which $e(E_0/\Q_3)=2$ and $pe/(p-1)=3$:
bimultiplicativity and antisymmetry; the Steinberg relation
$(u,1-u)=1$; Galois invariance $(\sigma u,\sigma v)=(u,v)$; triviality
on cubes, in particular on $-1$, on $\mu_8$ and on
$U^{(4)}_{E_0}\subset E_0^{\times3}$; and the depth bound
$(U^{(i)}_{E_0},U^{(j)}_{E_0})=1$ for $i+j\ge4$ (see \cite[Ch.~IV]{FV02} or
\cite[Ch.~XIV]{Ser79}).  A consequence used repeatedly: for $m\in U^{(3)}_{E_0}$
and any $u\in E_0^\times$,
\begin{equation}\label{W:eq:level0}
 (u,m)_{E_0}=(\varpi,m)_{E_0}^{\,v_{E_0}(u)},
\end{equation}
because $u=\varpi^{v(u)}\xi u_1$ with $\xi\in\mu_8$ and $u_1\in U^{(1)}_{E_0}$.
The computed tables are consistent with these facts (for instance
$(1+\varpi,1+\varpi^3)=(1+\varpi^2,1+\varpi^2i)=1$ while
$(1+\varpi,1+\varpi^2)\ne1$).  With
$\varpi,\ 1+\varpi,\ 1+\varpi^2,\ 1+\varpi^3$ as basis of $F^\times/F^{\times3}$ the
form is $(\varpi,1+\varpi^3)=(1+\varpi,1+\varpi^2)=\zeta$ and all other basis
pairings trivial; the $6\times6$ matrix for $E_0$ is in
Appendix~\ref{W:app:tables}.

\subsection{The cube identity and the star section on $K^6$}
This subsection proves (H1). It first shows that the standard section of the model is compatible with cubes. It then identifies the star section on $K^6$ and verifies the identity that makes $\psi_6$ a character.

\begin{proposition}[Cubes of the standard section]\label{W:prop:cubeidentity}
Let $w=x+yj\in D^\times$ with $y\ne0$, $T:=\Tr_{E_0/F}(x)\ne0$ and
$A:=x^2+x\sigma(x)+\sigma(x)^2+\varpi\Nm(y)\ne0$.  Then
$(w,0)^3=(w^3,0)$ in the cover.  The hypotheses hold for every
$w\in K^1\setminus E_0$; for $w\in E_0^\times$ the conclusion holds as well.
\end{proposition}
\begin{proof}
Put $n=\Nrd(w)=\Nm(x)-\varpi\Nm(y)$.  From $w^2=(x^2+\varpi\Nm y)+yTj$ and
$w^3=\cdots+yA\,j$ one reads off $x(\iota w)=\sigma(y)$,
$x(\iota w^2)=\sigma(y)T$, $x(\iota w^3)=\sigma(y)A$, with $T,A\in F$.
Writing $h_E,h_F$ for the exponents of the symbols and using
$h_E(a,b)=2h_F(a,b)$ for $a,b\in F$:
\[
 \beta(w,w)=h_E(T,T/n)=h_F(T,n),\qquad
 \beta(w^2,w)=h_E(A/T,A/n^2)=h_F(nT,A)+h_F(T,n),
\]
the twists vanishing because $h_E(n,n)=h_E(n^2,n)=0$.  Hence the exponent
of $(w,0)^3$ is $h_F(n,T)+h_F(nT,A)$.  Now $A=T^2-n$; with $u:=n/T^2$,
\[
 h_F(nT,T^2-n)=h_F(nT,T^2)+h_F(nT,1-u)=2h_F(n,T)+h_F(uT^3,1-u)=2h_F(n,T),
\]
by the Steinberg relation $h_F(u,1-u)=0$.  So the exponent is
$3h_F(n,T)=0$.  For $w\in K^1$, $x\equiv1\bmod\varpi$ gives $T\equiv2$ and
$v_F(x^2+x\sigma x+\sigma x^2)=2$ exactly, so $A\ne0$.  For $w=x\in E_0^\times$,
$x(\iota w^k)=\sigma(x)^k$ and $\det=\Nm(x)^k$, and the same bookkeeping
gives $\sigma(w,w)+\sigma(w^2,w)=-h_E(\sigma x,x)+h_E(\sigma x,x)=0$.
\end{proof}

\begin{corollary}[The star section on $K^6$]\label{W:cor:starK6}
\begin{enumerate}[label=(\alph*)]
\item For every cube $g=t^3$ with $t\in K^1$, the star lift is
$\sigma(g)=(g,0)$.
\item Every $X\in K^6$ can be written $X=zv^3$ with $z\in U_F^{(1)}$ and
$v\in K^2$ (Lemma~\ref{W:lem:cube}(d)); $z$ is determined modulo $F^{\times3}$,
and $z\equiv\Nrd(X)^{-1}\pmod{F^{\times3}}$.
\item The star lift of $X=x_X+y_Xj\in K^6$, that is, the element
$\ell(X):=\tilde z\,\tilde v^{\,3}$ with $\tilde z=(z,0)$, is
\begin{equation}\label{W:eq:ell}
 \ell(X)=(X,\kappa(X)),\qquad
 \kappa(X)=\begin{cases}h_E(y_X,\Nrd X),& y_X\ne0,\\ 0,& y_X=0.\end{cases}
\end{equation}
Equivalently, $\kappa(X)=c(X)\,v_{E_0}(y_X)$, where $c(X)\in\Fq_3$ is the
real part of the level-six digit of $X$ (the digit of $(x_X-1)/\varpi^3$).
\item (H1) is equivalent to the identity
\begin{equation}\label{W:eq:H1identity}
 \kappa(XX')=\kappa(X)+\kappa(X')+\beta(X,X')\qquad(X,X'\in K^6),
\end{equation}
and when it holds, $\psi_6$ is given on $\tilde K^6$ by
\begin{equation}\label{W:eq:psi6formula}
 \psi_6(X,e)=\zeta^{\,e-\kappa(X)}\,\theta_0(\Nrd X)^{-1},
\end{equation}
where $\tth(\zeta,z)=\zeta\theta_0(z)$ on $\tZ=\mu_3\times F^\times$,
the splitting being $z\mapsto(z,0)$ of Lemma~\ref{W:lem:model}(b), and
$\theta_0$ is any character with $\theta_0^3=1$.
\end{enumerate}
\end{corollary}
\begin{proof}
(a) is Proposition~\ref{W:prop:cubeidentity}.  (b) Existence is
Lemma~\ref{W:lem:cube}(d) ($\Cubes(Q_2)=Q_6$).  Taking reduced norms,
$\Nrd(X)=z^2\Nrd(v)^3$, so $z^2\equiv\Nrd X$ modulo $F^{\times3}$; if
also $X=z'v'^3$ then $(z/z')^2\in F^{\times3}$, hence $z/z'\in F^{\times3}$
because $F^\times/F^{\times3}$ has exponent $3$.  Finally
$z\equiv z^4\equiv\Nrd(X)^2\equiv\Nrd(X)^{-1}$ modulo cubes.  (c) $\ell(X)=(z,0)(v,0)^3=(z,0)(v^3,0)=(X,\beta(z,v^3))$
by (a), and $\beta(z,v^3)=h_E(z,x(\iota v^3)\det\iota v^3)$ by
Lemma~\ref{W:lem:model}(b).  If $y_X\ne0$ then $x(\iota v^3)=\sigma(y_X/z)$
and $\det\iota v^3=\Nrd(v)^3$, so
$\beta(z,v^3)=h_E(z,\sigma(y_X))-h_E(z,z)=h_E(z,y_X)$ by Galois
invariance and $h_E(z,z)=h_E(z,-1)=0$; by (b) this is
$-h_E(\Nrd X,y_X)=h_E(y_X,\Nrd X)$, which does not depend on the choice
of $z$ because $h_E(F^{\times3},\cdot)=0$.  If $y_X=0$ then
$x(\iota v^3)=\sigma(x_X/z)$, $\beta(z,v^3)=h_E(z,x_X)=h_E(x_X,\Nrd X)=h_E(x_X,x_X\sigma x_X)=h_E(x_X,\sigma x_X)=0$,
since $x_X\in U^{(3)}_{E_0}$ and $3+3\ge4$.  The digit formula follows
from \eqref{W:eq:level0}, from $h_E(\varpi,\Nrd X)=h_E(\varpi,\Nm x_X)=2h_E(\varpi,x_X)=-h_E(\varpi,x_X)$
(as $\Nrd X\equiv\Nm x_X$ modulo $U^{(7)}_F\subset F^{\times3}$ and
$h_E(\varpi,\sigma x_X)=h_E(\varpi,x_X)$), and from $h_E(1+\varpi^3,\varpi)=1$,
$h_E(1+\varpi^3i,\varpi)=0$ (Appendix~\ref{W:app:tables}); it is recorded
because the tables and scripts use it.  (d)
$\ell(X)\ell(X')=(XX',\kappa(X)+\kappa(X')+\beta(X,X'))$, so
\eqref{W:eq:H1identity} says exactly that $\ell$ is a homomorphism on
$K^6$.  If it is, then $\ell$ is a homomorphism on $U_6=F^\times K^6$,
because $\ell(zX)=(z,0)\ell(X)$ with $(z,0)$ central and $z\mapsto(z,0)$
a homomorphism (Lemma~\ref{W:lem:model}(b)); and then
$\tU_6=\mu_3\times\ell(U_6)$ and $\psi_6$ is the character which is
$\zeta\mapsto\zeta$ on $\mu_3$ and $zv^3\mapsto\theta_0(z)$ on
$\ell(U_6)$.  The latter is well defined and multiplicative on $U_6$:
if $zv^3\cdot z'v'^3=z''v''^3$ then $z''\in zz'F^{\times3}$ by the norm
argument of (b), and $\theta_0^3=1$.  Since $(X,e)=\zeta^{e-\kappa(X)}\ell(X)$
and $\theta_0(z)=\theta_0(\Nrd X)^{-1}$ by (b), this is
\eqref{W:eq:psi6formula}; the factor $\theta_0(\Nrd X)^{-1}$ is itself a
character of $K^6$, so the choice of $\theta_0$ affects neither (H1) nor
$\ell$.  Conversely (H1) implies \eqref{W:eq:H1identity}, since $\psi_6$
is then a character whose kernel meets each fiber of $\tU_6\to U_6$ in
exactly one point, namely $\ell(u)$.
\end{proof}

\begin{proposition}[(H1) holds]\label{W:prop:H1}
The identity \eqref{W:eq:H1identity} holds for all $X,X'\in K^6$.
Consequently $\psi_6$ is a character of $\tU_6$, given on
$\tilde K^6$ by \eqref{W:eq:psi6formula}.
\end{proposition}
\begin{proof}
Write $X=a+bj$, $X'=a'+b'j$ with $a,a'\in U^{(3)}_{E_0}$, $b,b'\in\varpi^3\Ocal_{E_0}$,
and $n=\Nrd X$, $n'=\Nrd X'$, so that $n,n'\in U_F^{(3)}$ and
$XX'=(aa'+\varpi b\sigma(b'))+(ab'+b\sigma(a'))j$.  Three preliminary
remarks.  First, $n=\Nm(a)-\varpi\Nm(b)\equiv\Nm(a)$ modulo
$U^{(7)}_F\subset E_0^{\times3}$, and likewise $n'\equiv\Nm(a')$, so in
every symbol $n$ and $n'$ may be replaced by $\Nm a$ and $\Nm a'$.
Second, by \eqref{W:eq:level0} and Galois invariance, for every
$u\in E_0^\times$,
\begin{equation}\label{W:eq:lambda}
 h_E(u,\Nm a)=v(u)\,\lambda(a),\qquad \lambda(a):=h_E(\varpi,\Nm a)=2h_E(\varpi,a)=h_E(a,\varpi),
\end{equation}
and similarly for $a'$.  Third, all symbols between two of
$a,\sigma a,a',\sigma a',n,n'$ vanish ($3+3\ge4$), and the twist
$2h_E(n,n')$ vanishes.  We write $v$ for $v_{E_0}$.

\emph{Main case: $b\ne0$, $b'\ne0$, $R\ne0$}, where
\[
 \tau:=\sigma(b'/b),\qquad R:=a'+\sigma(a)\tau,
\]
so that $x(\iota X)=\sigma(b)$, $x(\iota X')=\sigma(b')$ and
$x(\iota XX')=\sigma(ab'+b\sigma(a'))=\sigma(b)R$.  Inserting into
\eqref{W:eq:kubota} and using $h_E(R,R)=h_E(R,-1)=0$,
\begin{equation}\label{W:eq:betaK6}
 \beta(X,X')=h_E\Bigl(R,\frac{R}{\tau n}\Bigr)+2h_E(n,n')=-h_E(R,\tau)-h_E(R,n).
\end{equation}
On the other side, $y_{XX'}=ab'+b\sigma(a')=b\,\sigma(R)$, so by
\eqref{W:eq:ell}, Galois invariance and the first remark,
\[
 \kappa(XX')=h_E(b\sigma(R),nn')=h_E(b,n)+h_E(b,n')+h_E(R,n)+h_E(R,n').
\]
Hence \eqref{W:eq:H1identity} is equivalent to
$h_E(b,n')+h_E(R,n)+h_E(R,n')=h_E(b',n')-h_E(R,\tau)-h_E(R,n)$, i.e.\ (using $h_E(b'/b,n')=h_E(\sigma\tau,n')=h_E(\tau,n')$) to
\begin{equation}\label{W:eq:H1reduced}
 h_E(R,\tau)=h_E(\tau,n')-h_E(R,n')+h_E(R,n)=v(\tau)\lambda(a')-v(R)\bigl(\lambda(a')-\lambda(a)\bigr),
\end{equation}
the second equality by \eqref{W:eq:lambda}.  To prove
\eqref{W:eq:H1reduced}, factor $R=a'(1+\tau_2)$ with
$\tau_2:=\sigma(a)\tau/a'$; then $\tau_2\ne0,-1$ and
$\tau=\tau_2\cdot a'/\sigma(a)$, so
\[
 h_E(R,\tau)=h_E(a',\tau)+h_E(1+\tau_2,\tau_2)+h_E\bigl(1+\tau_2,\,a'/\sigma(a)\bigr).
\]
The first term is $v(\tau)\lambda(a')$ by \eqref{W:eq:level0} and
\eqref{W:eq:lambda}.  The second vanishes by the Steinberg relation
($h_E(1+\tau_2,\tau_2)=h_E(1-(-\tau_2),-\tau_2)+h_E(1+\tau_2,-1)=0$).
For the third, $a'/\sigma(a)\in U^{(3)}_{E_0}$ and $v(1+\tau_2)=v(R)$, so
by \eqref{W:eq:level0} it equals
$v(R)\bigl(h_E(\varpi,a')-h_E(\varpi,\sigma a)\bigr)=v(R)\bigl(h_E(\varpi,a')-h_E(\varpi,a)\bigr)=-v(R)\bigl(\lambda(a')-\lambda(a)\bigr)$.
This is \eqref{W:eq:H1reduced}.  Note that no case distinction on
$v(\tau)$ or on $v(1+\tau)$ was needed.

\emph{Case $b=0$, $b'\ne0$.}  Then $x(\iota X)=\sigma(a)$,
$x(\iota XX')=\sigma(a)\sigma(b')$, and \eqref{W:eq:kubota} gives
$\beta=h_E(\sigma b',\sigma a/n)=h_E(\sigma b',1/a)=-h_E(b',\sigma a)$.  Also
$\kappa(X)=0$, $\kappa(X')=h_E(b',n')$ and
$\kappa(XX')=h_E(ab',nn')=h_E(a,n)+h_E(a,n')+h_E(b',n)+h_E(b',n')$, with
$h_E(a,n)=h_E(a,\sigma a)=0$ and $h_E(a,n')=0$ by the third remark.  So
\eqref{W:eq:H1identity} reduces to $h_E(b',n)+h_E(b',\sigma a)=0$; by the
first remark $n\sigma(a)\equiv a\sigma(a)^2$, so the left side is
$h_E(b',a/\sigma a)+3h_E(b',\sigma a)=v(b')\bigl(h_E(\varpi,a)-h_E(\varpi,\sigma a)\bigr)=0$.

\emph{Case $b\ne0$, $b'=0$.}  Then $x(\iota X')=\sigma(a')$,
$x(\iota XX')=\sigma(b)a'$, and
\[\beta=h_E\bigl(a',\sigma(b)a'/(\sigma(a')n)\bigr)=h_E(a',\sigma b)=v(b)\lambda(a'),\]
the other three symbols vanishing by the third remark.  Also
$\kappa(X')=0$ and
\[\kappa(XX')=h_E(b\sigma(a'),nn')=h_E(b,n)+h_E(b,n')=\kappa(X)+v(b)\lambda(a').\]
Both sides of \eqref{W:eq:H1identity} equal $\kappa(X)+v(b)\lambda(a')$.

\emph{Case $b=b'=0$.}  All three $\kappa$ vanish and
$\beta=h_E(a,\sigma a')=0$.

\emph{Case $b,b'\ne0$, $R=0$.}  Then $XX'\in E_0$, so $\kappa(XX')=0$,
and $\tau=-a'/\sigma(a)$, $b'=\sigma(\tau)b=-b\,\sigma(a')/a$.  Now
$x(\iota XX')=\sigma(aa'+\varpi b\sigma b')$ and
$aa'+\varpi b\sigma(b')=aa'+\varpi\Nm(b)\tau=(a'/\sigma a)\,n$, so
$x(\iota XX')=\sigma(a')n/a$ and
\[
 \beta=h_E\Bigl(\frac{\sigma(a')n}{a\,\sigma(b)},\ \frac{\sigma(a')}{a\,\sigma(b')}\Bigr)+2h_E(n,n').
\]
Expanding bimultiplicatively, discarding the symbols that vanish by the
third remark, and using Galois invariance, \eqref{W:eq:level0} and
$v(b')=v(b)$:
\begin{align*}
 \beta&=-h_E(\sigma a',\sigma b')-h_E(n,\sigma b')+h_E(a,\sigma b')-h_E(\sigma b,\sigma a')+h_E(\sigma b,a)+h_E(\sigma b,\sigma b')\\
 &=v(b)\bigl[h_E(\varpi,a')+2h_E(\varpi,a)-h_E(\varpi,a)-h_E(\varpi,a')+h_E(\varpi,a)\bigr]+h_E(b,b'),
\end{align*}
and $h_E(b,b')=h_E(b,-b\sigma(a')/a)=h_E(b,\sigma a')-h_E(b,a)=v(b)\bigl[h_E(\varpi,a')-h_E(\varpi,a)\bigr]$.
Altogether $\beta=v(b)\bigl[h_E(\varpi,a)+h_E(\varpi,a')\bigr]$, while
$\kappa(X)+\kappa(X')=h_E(b,n)+h_E(b',n')=v(b)\bigl[2h_E(\varpi,a)+2h_E(\varpi,a')\bigr]$.
The sum $\kappa(X)+\kappa(X')+\beta=3v(b)[\cdots]=0=\kappa(XX')$.

This exhausts all cases, since $x(\iota XX')$ is $\sigma(b)R$ when
$b,b'\ne0$ and the entries used in each case are the ones prescribed by
\eqref{W:eq:kubota}.  The last assertion is Corollary~\ref{W:cor:starK6}(d).
\end{proof}

\begin{remark}[Checks and a second route]\label{W:rem:H1-checks}
(a) Script~12 evaluates \eqref{W:eq:H1reduced} and \eqref{W:eq:H1identity}
in the explicit model on $12{,}000$ structured instances covering all
five cases of the proof (including $v(1+\tau)\in\{1,\dots,5\}$ and
$X'\in K^8,\dots,K^{12}$); all pass, and $\beta(X,X')\ne0$ in about
$44\%$ of them, so the identity is not vacuous.  Scripts~8 and~9 are
earlier checks of the same identity through cube roots computed by
Hensel's lemma and on a stratified sample.  None of these is part of
the proof.
(b) A second, independent argument reduces (H1) to a finite problem.
The subgroup $K^4$ is a uniform pro-$3$ group of dimension $8$, so
central $\mu_3$-extensions of $K^4$ are classified by alternating forms
on $K^4/K^8$ \cite[Ch.~V]{Laz65}; \cite[\S5]{SW00}; the form of the cover is
$\Ad(D^\times)$-invariant and vanishes on pairs of commuting elements by
\cref{I:prop:comm} and the depth bounds for the symbols of
$E_0$ and of the ramified quadratic extensions of $F$; script~5a finds
that the space of such forms is zero, so the cover splits over $K^4$
and the star section is a homomorphism on $K^8$.  Script~5b then shows
that every odd function on $Q_6/Q_8$ whose non-additivity is
$D^\times$-equivariant is linear.  Together these reduce (H1) to
multiplicativity of $\ell$ on $K^6\times K^8$, which is a special case
of Proposition~\ref{W:prop:H1}.  We record this route because the
Lazard-form computation is of independent interest for Target~B.
\end{remark}

\section{The main theorem}\label{W:sec:maintheorem}
This section assembles the earlier sections into the main theorem, Target~A, and lists what its proof depends on. A corollary extends the result to the representations trivial on $K^1$ and their twists.

\begin{theorem}[\Pctag\ Target A]\label{W:thm:main-targetA}
Let $F=\Q_3(\zeta_3)$, $n=3$, $D$ the quaternion division algebra over
$F$, $\tH$ its scalar-central $3$-fold cover, and $\tth$ a genuine
central character compatible with the trivial character of $F^\times$.
Then $L_D\mathbf1=\Theta_\tau$ for a unique genuine irreducible
representation $\tau$ of $\tH$, of dimension $27$ and central character
$\tth$.  Moreover:
\begin{enumerate}[label=(\alph*)]
\item $\tau|_{\tU_1}$ is irreducible and equals
$\Ind_{\tU_3}^{\tU_1}\lambda_0$, with $\lambda_0$ as in \eqref{W:eq:lambda0};
$\tau|_{\tU_6}=27\,\psi_6$, and $\tau|_{\tU_3}$ is the sum of the $27$
distinct genuine characters of $\tU_3/N_6$ whose restriction to
$\tU_4/N_6$ exhausts the genuine characters of that group.
\item On the covered quadratic tori, $\tau$ restricts as in
Proposition~\ref{W:prop:all-three-tori}: $(3,2,2,2)$ on the unramified
torus and $(3,2,1,1,1,1)$ on each ramified torus.
\item $\Theta_\tau$ is supported on the star lifts of good elements and
of the center; its values are $27\zeta\tth(\tilde z)$ at
$\zeta\tilde zt^*$ with $t\in F^\times K^2$, $9\zeta\tth(\tilde z)$ for
$t\in F^\times K^1\setminus F^\times K^2$, and $3\zeta\tth(\tilde z)$ for
$t\notin F^\times K^1$.
\item $\tau$ is the unique extension to $\tH$ of $\tau|_{\tU_1}$ with
character values $3\zeta\tth(\tilde z)$ at the star lifts of the elements
of $2$-power order modulo the center.
\end{enumerate}
\end{theorem}
\begin{proof}
(H1) holds by Proposition~\ref{W:prop:H1}; Theorem~\ref{W:thm:H1equiv}
gives Target~A together with (a), and Theorem~\ref{W:thm:reduction} gives
uniqueness and irreducibility of the restriction.  The dependencies of
this chain are itemized in Remark~\ref{W:rem:dependencies}.  (b) is
Proposition~\ref{W:prop:all-three-tori}; (c) collects
\eqref{W:eq:Phivalues} and the proof of Lemma~\ref{W:lem:2groups}; (d)
follows from (c) since the four extensions of an $\tH$-stable
irreducible of $\tU_1$ differ by the linear characters of $D_8$, which
are distinguished by their values on $P_0$.
\end{proof}

\begin{remark}[Dependencies of Theorem~\ref{W:thm:main-targetA}]\label{W:rem:dependencies}
The proof uses the following inputs and nothing else.
\begin{enumerate}[label=(D\arabic*),leftmargin=*]
\item \Etag\ From Chapter~\ref{ch:I}: the cover of $D^\times$ is the restriction
along $\iota$ of the parameter-$c_v$ Kazhdan--Patterson cover of
$\mathrm{GL}_2(E_0)$ (\cref{I:lem:center}, \cref{I:prop:restriction}); the commutator
formula on covered tori (\cref{I:prop:comm}); good elements are
$F^\times\cdot\{\text{cubes}\}$ in each torus (\cref{I:prop:good}(a));
$E^\times/F^\times$ has no $3$-torsion (\cref{I:lem:FcapEn}); the star map
(\cref{I:thm:ident}, \cref{I:rem:D2tame}).  From Chapter~\ref{ch:III}: the normalization
\eqref{W:eq:LD} of $L_D$ with $b_D=3$, $c_D=27$, and the value $R=1$ at
elements of $2$-power order.
\item \Etag\ Kubota's cocycle on $\mathrm{GL}_2(E_0)$
(Lemma~\ref{W:lem:model}(a)), and the standard properties of the cubic
Hilbert symbols of $F$ and $E_0$ listed after Lemma~\ref{W:lem:model}.
\item \Pctag\ The explicit values of the two symbols
(Appendix~\ref{W:app:tables}), obtained as the unique Steinberg symbols
by exact linear algebra (script~7).  They enter the hand arguments only
through $h_E(1+\varpi^3,\varpi)\ne0$ (used to exclude $e=0$ in
Proposition~\ref{W:prop:design} via $\kappa_{1,5}(1,1)\ne1$) and through
the digit form of $\kappa$ used by the scripts; the proof of
Proposition~\ref{W:prop:H1} does not use them.
\item \Pctag\ Four exhaustive finite computations: the $72$ elements of
$\bar G_3$ from the $648$ depth-one classes (script~1, also proved by
hand in Proposition~\ref{W:prop:G3explicit}); the cube map on $Q/Q_7$
(script~2, corroborating Lemma~\ref{W:lem:cube}); the actions on $V$ and
their orbits (script~3); and the design search over all odd
$D_8$-invariant functions on the $36$ lines together with the
Heisenberg cocycle extension and the Fourier identities (scripts~4
and~6).  Each is an enumeration of an explicitly described finite set.
\item Everything else is proved by hand: the toral theorem and
the three torus calculations, the positivity criterion, the depth
lemma, the Brauer reduction (Theorem~\ref{W:thm:reduction}), the
consequences of (H1) and the uniqueness of the design given the finite
data (Section~\ref{W:sec:finite}), the cube identity, and (H1) itself.
\end{enumerate}
No sampled computation is load bearing.  Scripts~5, 8--12 are
corroboration (Remarks~\ref{W:rem:H1-checks} and the next remark).
Theorem~\ref{W:thm:leveltwo} depends on no further computation; its
Steps~2--3 and Proposition~\ref{W:prop:free} are proved by hand from the
items above.
\end{remark}

\begin{remark}[End-to-end verification]
Independently of the chain of reductions, the character of
$\Ind_{\tU_3}^{\tU_1}\lambda_0$ was computed in the explicit model of
Section~\ref{W:sec:model} from formula~\eqref{W:eq:lambda0} (products of the
eight $D_8$-conjugates, $27$ coset representatives of $\tU_1/\tU_3$) and
compared with \eqref{W:eq:Phivalues} at star lifts of cubes of depth-one
elements, at elements of $\tU_6$ with all three values of $\psi_6$, at
products $\sigma(g)\ell(k)\tilde z\zeta^m$, and at non-good elements of
depths $3,4,5$; all values agree, including phases (scripts~11, 11b).
Script~10 checks, again in the explicit model, that all $124$ sampled
coplanar triples of good lines have trivial defect and that all $186$
sampled four-term relations $\sum v_i=0$ in $\bar G_3$ have defect
$\sum e_*(v_i)$; that $\omega=0$ on $300$ pairs; and that
$\kappa_{1,5}(a,b)=\zeta^{-\Tr(a\bar b)}$ on $300$ pairs.
\end{remark}

\begin{corollary}[\Pctag\ Level-zero rows and their norm twists]\label{W:cor:levelzero}
Let $\rho$ be one of the five irreducible representations of $D^\times$
trivial on $K^1$ (inflated from $D^\times/F^\times K^1\simeq D_8$), and let
$\chi$ be a character of $F^\times$ trivial on $\mu_3$.  Then
$L_D\Theta_{\rho\otimes(\chi\circ\Nrd)}$, computed with any compatible
genuine central character, is the character of a genuine irreducible
representation of $\tH$, of dimension $27\dim\rho$.
\end{corollary}
\begin{proof}
First let $\chi=1$.  On $\tU_1$, $\Theta_\rho\circ p=\dim\rho$, so
$L_D\Theta_\rho|_{\tU_1}=(\dim\rho)\Phi$ is a character by
Theorem~\ref{W:thm:main-targetA}; on $\Gamma_{P_0}$ it is
$3\,\psi_P\otimes\rho+3\dim\rho\sum_\pi(\dim\pi)\Theta_\pi$; and its
restrictions to covered tori are characters by
Lemma~\ref{W:lem:toral-Fourier}.  The proof of Theorem~\ref{W:thm:reduction}
applies verbatim (Lemma~\ref{W:lem:norm} holds for every row by
Chapter~\ref{ch:III}), giving an irreducible character.  For general $\chi$, write
$\chi=\chi_1^3$: this is possible because $\chi$ is trivial on $\mu_3$,
so that its restriction to $U_F^{(1)}\simeq\mu_3\times\Z_3^2$ factors
through $\Z_3^2$, whose character group is $3$-divisible, while on
$\varpi^{\Z}\times\mu_2$ cube roots exist trivially.  Twisting the
representation $\tau_\rho$ just obtained by $\chi_1\circ\Nrd\circ p$
multiplies its character at $\tilde zt^*$ by $\chi_1(z)^2\chi(\Nrd t)$
and changes its central character to
$\tth':=\tth\cdot(\chi_1^2\circ p)$, which is compatible with
$\theta_\rho\chi^2=\theta_{\rho\otimes\chi\circ\Nrd}$ because
$\tth'(a^*)=\theta_\rho(a)\chi_1(a^3)^2$.  The result is
$L_D\Theta_{\rho\otimes\chi\circ\Nrd}$ computed with $\tth'$; by
Remark~\ref{W:rem:theta-independence} any other compatible choice is
obtained by a further twist.
\end{proof}

\section{Target B: other rows}\label{W:sec:plan}
This section turns to the other rows of the character matrix, that is, to $L_D\Theta_\rho$ for a general irreducible $\rho$. It extends the reduction theorem to every row and compares each row with the row of the trivial character by a product formula. It then proves Target~B for the representations trivial on $K^3$.

Throughout this section $\rho$ is an irreducible smooth representation of
$D^\times$ (finite dimensional, as $D^\times/F^\times$ is compact) with
central character $\theta_\rho$ trivial on $\mu_3$, $\theta_0$ is a
character of $F^\times$ with $\theta_0^3=\theta_\rho$, and
$\tth(\zeta,z)=\zeta\theta_0(z)$ is the corresponding compatible genuine
central character, in the coordinates of Lemma~\ref{W:lem:model}(b).  We
write $\phi_\rho:=L_D\Theta_\rho$ and $\Phi_\rho:=\phi_\rho|_{\tU_1}$.
For $\rho=\mathbf1$ we take, once and for all, $\theta_1:=\mathbf1$,
i.e.\ $\tth_1(\zeta,z)=\zeta$ (a compatible choice, by
Remark~\ref{W:rem:theta-independence}; it is the normalization of the
scripts), and write $\phi_1$, $\tth_1$, $\lambda_0$ for the objects of
Sections~\ref{W:sec:finite}--\ref{W:sec:maintheorem} in this normalization.
Two reductions are free.

\paragraph{Twisting.}  For a character $\chi$ of $F^\times$ trivial on
$\mu_3$, $\rho\otimes(\chi\circ\Nrd)$ is irreducible with central
character $\theta_\rho\chi^2$, and by the argument of
Corollary~\ref{W:cor:levelzero}, $L_D\Theta_{\rho\otimes\chi\circ\Nrd}$ is
an irreducible character if and only if $L_D\Theta_\rho$ is.  Squaring
is surjective onto the characters of $F^\times=\varpi^{\Z}\times\mu_2\times U_F^{(1)}$
that are trivial at $-1$, so we may and do assume
\begin{equation}\label{W:eq:normalized-central}
 \theta_\rho\in\{\mathbf1,\ \mathrm{sgn}\},\qquad
 \mathrm{sgn}(\varpi^k\varepsilon u):=\varepsilon\quad(\varepsilon=\pm1,\ u\in U_F^{(1)}),
\end{equation}
according to whether $\rho(-1)=+1$ or $-1$, and accordingly
$\theta_0=\theta_1=\mathbf1$ or $\theta_0=\mathrm{sgn}$.

\paragraph{Unitarity.}  By Chapter~\ref{ch:III}, $\langle\phi_\rho,\phi_\rho\rangle_{\tH}=1$
and $\phi_\rho(1)=27\dim\rho>0$, so $\phi_\rho$ is an irreducible
character as soon as it is a virtual character.

\subsection{The reduction theorem for every row}
This subsection extends the reduction theorem of Section~\ref{W:sec:reduction} from the trivial character to every $\rho$.

\begin{theorem}[Reduction, all rows]\label{W:thm:reduction-all}
$L_D\Theta_\rho$ is the character of a genuine irreducible representation
of $\tH$ if and only if $\Phi_\rho$ is a virtual character of $\tU_1$.
\end{theorem}
\begin{proof}
Only ``if'' needs proof.  Assume \eqref{W:eq:normalized-central}.  The
finite quotient $\Gamma=\tH/N$ of Section~\ref{W:sec:reduction} is built
from the finite row of $\phi_\rho$ exactly as before, with
$N\cap\tZ=\ker\tth\cap\ker\tth_1$ (so that both $\phi_\rho$ and $\phi_1$
are functions on $\Gamma$); then $T=\tZ N/N$ is $\mu_3$ if
$\theta_\rho=\mathbf1$ and $\mu_6=\mu_3\times\mu_2$ if
$\theta_\rho=\mathrm{sgn}$; $\Gamma_0=\tU_1/N$ is the Sylow
$3$-subgroup in the first case and is $\Gamma_0^{(3)}\times\mu_2$ with
$\Gamma_0^{(3)}$ the Sylow $3$-subgroup in the second; in both cases
$\Gamma/\Gamma_0$ is $D_8$.  We check the three inputs of the proof
of Theorem~\ref{W:thm:reduction}.

\emph{Tori.}  $\phi_\rho|_{\widetilde{E^\times}}$ is a genuine
representation character for every maximal subfield $E$
(Lemma~\ref{W:lem:toral-Fourier}: $\Theta_\rho|_{E^\times}$ is a sum of
characters $\alpha$ with $\alpha|_{F^\times}=\theta_\rho$), so
Lemma~\ref{W:lem:toral-descent} holds verbatim.

\emph{Centralizers.}  Lemma~\ref{W:lem:centralizers} concerns $\Gamma_0$
and the images of tori and is unchanged; the $\mu_2$ in $T$ is central.

\emph{Two-groups.}  Let $\hat P_0$ be the preimage of $P_0F^\times$ and
$\Gamma_{P_0}=\hat P_0N/N$, a central extension of $D_8$ by $T$; it
contains a Sylow $2$-subgroup of $\Gamma$.  Every $\eta\in\hat P_0\subset D^\times$
is good: $\eta^4\in F^\times$, so $\eta=\eta^4\cdot(\eta^{-1})^3$.  Put
$x_\eta:=\widetilde{\eta^4}\,(\tilde\eta^{-1})^3$, a lift of $\eta$ with
$x_{\eta c}=x_\eta\tilde c$ for $c\in F^\times$.  By \eqref{W:eq:LD} and $R=1$
off $F^\times K^1$,
\[
 \phi_\rho(x_\eta)=3\,\theta_0(\eta^4)\,\Theta_\rho(\eta^{-1})
 =3\,\theta_0(\eta^4)\,\Theta_{\rho^\vee}(\eta)\qquad(\eta\notin F^\times),
\]
and $\phi_\rho=27\dim\rho\cdot\tth$ on $\tZ$.  Let $\psi_P$ be the
genuine character of order $3$ of $\Gamma_{P_0}$ trivial on a Sylow
$2$-subgroup (unique, by Schur--Zassenhaus for the $\mu_3$-part of $T$);
on the image of $\tZ$ it is $\tth_1$, i.e.\ $\psi_P(z,0)=\theta_1(z)$, and
therefore $\psi_P(x_\eta)=\theta_1(\eta^4)\psi_P(\tilde\eta^{-1})^3=\theta_1(\eta^4)$.
(We do not twist $\rho^\vee$ by $\theta_1^2\circ\Nrd$: on scalars that
twist is $\theta_1(z^2)^2=\theta_1(z)$, not trivial.)

If $\theta_\rho=\mathbf1$, then $\theta_0=\theta_1$, so
$\phi_\rho(x_\eta)=3\psi_P(x_\eta)\Theta_{\rho^\vee}(\eta)$, and
$\rho^\vee$ is trivial on $F^\times$, so $\Theta_{\rho^\vee}$ is the
character of a representation of $D_8=\hat P_0/F^\times$.  With
$\Gamma_{P_0}=\mu_3\times\tilde P_0$ as in Lemma~\ref{W:lem:2groups},
\[
 \phi_\rho|_{\Gamma_{P_0}}=\psi_P\otimes\Bigl(3\,\rho^\vee|_{D_8}+3\dim\rho\sum_{\pi\in\Irr(D_8)}(\dim\pi)\,\Theta_\pi\Bigr),
\]
a character (the second term is $24\dim\rho$ times the delta function at
$1\in D_8$).

If $\theta_\rho=\mathrm{sgn}$, then $\theta_0=\mathrm{sgn}\cdot\theta_1$ and
$\phi_\rho(x_\eta)=3\psi_P(x_\eta)\,\mathrm{sgn}(\eta^4)\,\Theta_{\rho^\vee}(\eta)$,
where now $\rho^\vee$ has central character $\mathrm{sgn}$.
Here $\Gamma_{P_0}=\mu_3\times\Gamma''$ with $\Gamma'':=\Gamma_{P_0}/\mu_3$,
the quotient of $\hat P_0$ by $\mu_3$ and the lifts of
$\ker\mathrm{sgn}\subset F^\times$, a central extension of $D_8$ by
$\mu_2$; $\rho^\vee|_{\hat P_0}$ factors through $\Gamma''$.  With $a,b$
the images of $1+i$ and $j$ one has $a^8=1$ (as $(1+i)^4=-4$ and
$\mathrm{sgn}(-4)=-1$), $b^2=1$ and $bab^{-1}=a^3$ (as
$j(1+i)j^{-1}=2(1+i)^{-1}$ and $\mathrm{sgn}(2)=-1$), so
$\Gamma''\simeq SD_{16}$, the semidihedral group, with $a^4$ the image of
$-1$.  The function $f(\eta):=\mathrm{sgn}(\eta^4)$ is $-1$ on the classes
$\{a,a^3\},\{a^5,a^7\}$ and $+1$ elsewhere.  Having central character
$\mathrm{sgn}$, $\rho^\vee|_{\Gamma''}$ is a sum of copies of the two
faithful two-dimensional characters $\chi_1,\chi_2$ of $SD_{16}$
($a\mapsto\pm i\sqrt2$, $a^2\mapsto0$, $a^4\mapsto-2$, $0$ off
$\langle a\rangle$); a direct check on the seven classes gives
$f\chi_1=\chi_2$ and $f\chi_2=\chi_1$.  Hence $f\cdot\Theta_{\rho^\vee}$
is a character of $\Gamma''$, and with $\delta^{\mathrm{sgn}}$ the
function $\pm1$ at $a^{0},a^4$ and $0$ elsewhere,
$8\delta^{\mathrm{sgn}}=2\chi_1+2\chi_2$, so
\[
 \phi_\rho|_{\Gamma_{P_0}}=\psi_P\otimes\bigl(3f\,\Theta_{\rho^\vee|_{\Gamma''}}+3\dim\rho\,(2\chi_1+2\chi_2)\bigr)
\]
is a character.

\emph{Brauer's theorem.}  Let $B=B_2\times B_3\subset\Gamma$ be elementary,
and let $\bar B_2$ be the image of $B_2$ in $\Gamma/\Gamma_0=D_8$.  (When
$T\supset\mu_2$, $B_2$ may be noncyclic while $\bar B_2$ is cyclic, e.g.\
$B_2=\langle a^4,b\rangle$; the case distinction must therefore be made
on $\bar B_2$.)  If $\bar B_2=1$, then $B\subset\Gamma_0$ and
$\phi_\rho|_B=\Phi_\rho|_B$ is a virtual character by hypothesis.  If
$\bar B_2=\langle\bar y\rangle$ is cyclic and nontrivial, pick $y\in B_2$
over $\bar y$; then $B_2\subset\langle y\rangle\cdot(B_2\cap T)$, $B_3\subset C_{\Gamma_0}(y)\subset T_{E_y}$
by Lemma~\ref{W:lem:centralizers}(a), $y\in T_{E_y}$ and $T\subset T_{E_y}$,
so $B\subset T_{E_y}$ and Lemma~\ref{W:lem:toral-descent} applies.  If
$\bar B_2$ is noncyclic, then $B_3\subset C_{\Gamma_0}(B_2)\subset T$ by
Lemma~\ref{W:lem:centralizers}(b) (whose proof uses only the images of the
elements of $B_2$ in $D_8$), so $B\subset B_2T$, which is conjugate into
$\Gamma_{P_0}$ because $\Gamma_{P_0}$ contains a Sylow $2$-subgroup; the
two-group computation applies.

Brauer's theorem now gives that $\phi_\rho$ is a virtual character of
$\Gamma$, and unitarity finishes the proof.
\end{proof}

\subsection{The cube-root transform}
This subsection writes $L_D\Theta_\rho$ as the product of $L_D\mathbf1$ and a function built from $\Theta_\rho$ and the cube root. It then gives a sufficient criterion for a row in terms of this product.

\begin{lemma}[Product formula]\label{W:lem:product}
Assume \eqref{W:eq:normalized-central}.  There is a unique function
$F_\rho$ on the set of star lifts of good elements of $H$ such that
\begin{equation}\label{W:eq:product}
 \phi_\rho=\phi_1\cdot F_\rho\quad\text{on that set},\qquad
 F_\rho(\zeta\tilde z\,t^*)=\theta_0\theta_1^{-1}(z)\,\Theta_\rho(t),
\end{equation}
and both sides of \eqref{W:eq:product} vanish elsewhere.  $F_\rho$ is
invariant under conjugation by $\tH$ and under multiplication by
$\mu_3$.  On the star lifts of good elements of $U_3$,
\begin{equation}\label{W:eq:Frho-U3}
 F_\rho(\tilde z\,t^*)=\varepsilon_\rho(z)\,\Theta_\rho(t)\qquad(zt^3\in U_3,\ t\in K^1),
\end{equation}
where $\varepsilon_\rho=\mathbf1$ if $\theta_\rho=\mathbf1$ and
$\varepsilon_\rho=\mathrm{sgn}$ if $\theta_\rho=\mathrm{sgn}$; here
$t\in K^1$ is determined by $zt^3$ up to $U_F^{(1)}$, on which
$\Theta_\rho$ is invariant, and $z$ up to $(U_F^{(1)})^3\subset\ker\varepsilon_\rho$.
In particular $g\mapsto\varepsilon_\rho(z)$ is the character of $U_3=F^\times K^3$
trivial on $K^3$ and equal to $\varepsilon_\rho$ on $F^\times$ (well
defined because $F^\times\cap K^3=U_F^{(2)}$), and it is $\Ad(U_1)$-invariant.
\end{lemma}
\begin{proof}
Both $\phi_\rho$ and $\phi_1$ vanish off star lifts of good elements, and
there $\phi_\rho(\zeta\tilde zt^*)=3\zeta\theta_0(z)R(t)\Theta_\rho(t)$,
$\phi_1(\zeta\tilde zt^*)=3\zeta\theta_1(z)R(t)$ with $R(t)\ne0$.  If
$\tilde zt^*=\tilde z't'^*$ then $t'=ct$ with $c\in F^\times$ and
$z'=zc^{-3}$ (as in Lemma~\ref{W:lem:Phi}), and
$\theta_0\theta_1^{-1}(c^{-3})\Theta_\rho(ct)=\theta_\rho(c)^{-1}\theta_\rho(c)\Theta_\rho(t)$,
so $F_\rho$ is well defined.  Invariance is clear.  For $zt^3\in U_3$
with $t\in K^1$: $\theta_0\theta_1^{-1}=\varepsilon_\rho$ on $F^\times$
by the normalization, and $c\in F^\times\cap K^1=U_F^{(1)}$, on which
$\theta_\rho=1$.
\end{proof}

\begin{remark}[Adams operations]
$F_\rho$ restricted to $\tU_1$ is the Adams-type transform
$\psi^{1/3}\Theta_\rho$, ``$\Theta_\rho$ at the cube root,'' corrected by
$\varepsilon_\rho$.  At a tame place $\psi^{1/3}$ is a Galois conjugation
of characters (\cref{III:lem:tame-adams}), which is why Target~B is trivial
there; at the wild place it is not even a virtual character in general,
and the content of the next criterion is that it need not be: it must
only be compatible with the character $\lambda_0$ on the good set.
\end{remark}

\begin{proposition}[\Pctag\ A sufficient criterion]\label{W:prop:criterion}
Assume \eqref{W:eq:normalized-central} and Theorem~\ref{W:thm:main-targetA}
(the proposition uses $\lambda_0$ and $\Theta_{\Ind\lambda_0}=\Phi_1$, so
it inherits the label of Target~A).  Suppose there is a virtual
character $\Psi$ of $U_3$ (a $\Z$-combination of characters of
finite-dimensional smooth representations) which is $\Ad(U_1)$-invariant
and satisfies $\Psi(zt^3)=F_\rho(\tilde zt^*)$ for all good $zt^3\in U_3$.
Then $L_D\Theta_\rho$ is the character of a genuine irreducible
representation $\tau_\rho$ of $\tH$ of dimension $27\dim\rho$, and
\begin{equation}\label{W:eq:tau-rho}
 \Theta_{\tau_\rho}|_{\tU_1}=\sum_i n_i\,\Theta_{\Ind_{\tU_3}^{\tU_1}(\lambda_0\otimes\sigma_i\circ p)},
 \qquad \Psi=\sum_i n_i\Theta_{\sigma_i},\ n_i\in\Z .
\end{equation}
\end{proposition}
\begin{proof}
For a finite-dimensional representation $\sigma$ of $U_3$ with
$\Ad(U_1)$-invariant character, the induced character satisfies, for
$x\in\tU_3$ and $g$ running over $\tU_1/\tU_3$,
\[
 \Theta_{\Ind(\lambda_0\otimes\sigma\circ p)}(x)=\sum_g\lambda_0(gxg^{-1})\Theta_\sigma(p(gxg^{-1}))
 =\Theta_\sigma(p(x))\sum_g\lambda_0(gxg^{-1})=\Theta_\sigma(p(x))\,\Theta_{\Ind\lambda_0}(x),
\]
and both sides vanish off $\tU_3$.  By Theorem~\ref{W:thm:main-targetA}(a),
$\Theta_{\Ind\lambda_0}=\Phi_1$, which is supported on star lifts of good
elements of $U_3$ (Lemma~\ref{W:lem:Phi}).  Summing with the $n_i$ and
using \eqref{W:eq:product}, $\sum_in_i\Theta_{\Ind(\lambda_0\otimes\sigma_i\circ p)}=\Phi_1\cdot(\Psi\circ p)=\Phi_1F_\rho=\Phi_\rho$
on $\tU_1$.  So $\Phi_\rho$ is a virtual character of $\tU_1$, and
Theorem~\ref{W:thm:reduction-all} applies.
\end{proof}

\subsection{All representations trivial on $K^2$}
This subsection proves Target~B for the representations of $D^\times$ that are trivial on $K^2$.

Let $\pi_3:U_3\to U_3/U_4\simeq Q_3/Q_4\simeq\Fq_9$ be the depth-three
digit (the $A$-coordinate of \eqref{W:eq:Vcoords}), a homomorphism which
is $\Ad(U_1)$-invariant because $[U_1,U_3]\subset U_4$, and identify
$K^1/K^2\simeq\Fq_9$ by $1+wj\mapsto w$.

\begin{theorem}[\Pctag\ Target B for level at most one]\label{W:thm:levelone}
Assume Theorem~\ref{W:thm:main-targetA}.  Let $\rho$ be an irreducible
representation of $D^\times$ trivial on $K^2$, with
$\theta_\rho|_{\mu_3}=1$ (automatic, since $\mu_3\subset U_F^{(1)}\subset K^2$).
Then $L_D\Theta_\rho$, computed with any compatible genuine central
character, is the character of a genuine irreducible representation
$\tau_\rho$ of $\tH$ of dimension $27\dim\rho$.  After the normalization
\eqref{W:eq:normalized-central},
\[
 \tau_\rho|_{\tU_1}=\bigoplus_{\psi}m_\psi\,\Ind_{\tU_3}^{\tU_1}\bigl(\lambda_0\otimes(\varepsilon_\rho\cdot\psi\circ\pi_3)\circ p\bigr),
 \qquad \rho|_{K^1}=\bigoplus_\psi m_\psi\,\psi ,
\]
the sum over the characters $\psi$ of $K^1/K^2\simeq\Fq_9$; the
induced representations occurring are irreducible and pairwise
distinct.
The same conclusion holds for $\rho\otimes(\chi\circ\Nrd)$ for every
character $\chi$ of $F^\times$ trivial on $\mu_3$; these are all the
irreducible representations of $D^\times$ whose restriction to $K^2$
is a multiple of a character of the form $\chi\circ\Nrd$.
\end{theorem}
\begin{proof}
By twisting we may assume \eqref{W:eq:normalized-central}.  Let $g=zt^3$
be a good element of $U_3$ with $t\in K^1$.  If $t\in F^\times K^2$, then
$g\in U_6$, $\pi_3(g)=0$ and $\Theta_\rho(t)=\dim\rho=\Theta_{\rho|_{K^1}}(0)$.
Otherwise $t\equiv1+wj$ modulo $K^2$ with $w\ne0$, and by
Lemma~\ref{W:lem:cube}(c) the depth-three digit of $g$ is $w^5$; since
$w^4=\Nm_{\Fq_9/\Fq_3}(w)=\pm1$, $w^5=\pm w$, so $w=\pm\pi_3(g)$.  Because
$\Theta_\rho$ is a class function and $i(1+wj)i^{-1}=1-wj$, we get
$\Theta_\rho(t)=\Theta_{\rho|_{K^1}}(\pi_3(g))$ in both cases.  By
\eqref{W:eq:Frho-U3}, $F_\rho(\tilde zt^*)=\varepsilon_\rho(z)\Theta_{\rho|_{K^1}}(\pi_3(g))$,
which is the value at $g$ of the character of the representation
$\Psi:=\varepsilon_\rho\otimes(\rho|_{K^1})\circ\pi_3$ of $U_3$ (with
$\rho|_{K^1}$ regarded as a representation of $K^1/K^2$).  Its character
is $\Ad(U_1)$-invariant, so Proposition~\ref{W:prop:criterion} applies
with all $n_i=m_\psi\ge0$, giving the displayed decomposition.
The restriction $\tau_\rho|_{\tU_1}$ has norm
$\int_{U_1/Z}|\Theta_\rho|^2\,d\mu$ by the Weyl-integration computation
of Lemma~\ref{W:lem:norm} (substituting $t=s^3$, $|\phi_\rho|^2(s^3)=9R(s)^2|\Theta_\rho(s)|^2$),
and this is $\langle\Theta_\rho|_{K^1},\Theta_\rho|_{K^1}\rangle=\sum_\psi m_\psi^2$.
On the other hand $\langle\sum m_\psi\chi_\psi,\sum m_\psi\chi_\psi\rangle\ge\sum m_\psi^2\langle\chi_\psi,\chi_\psi\rangle\ge\sum m_\psi^2$
for the characters $\chi_\psi$ of the induced representations, with
equality only if each $\chi_\psi$ with $m_\psi>0$ is irreducible and
they are pairwise orthogonal.  The last sentence: if $\rho|_{K^2}$ is a multiple of
$\chi\circ\Nrd$, then $\rho\otimes(\chi^{-1}\circ\Nrd)$ is trivial on
$K^2$; and $\chi$ may be modified by any character trivial on
$\Nrd(K^2)=U_F^{(1)}$ without changing this.
\end{proof}

\begin{remark}
(a) The theorem contains Corollary~\ref{W:cor:levelzero} and extends it to
the level-zero representations with $\rho(-1)=-1$ (the two-dimensional
representations of $D^\times/K^1$ induced from the characters of
$\mu_8$ of odd order), which are not twists of representations of $D_8$.
(b) For the four four-dimensional irreducible representations of
$\Fq_9\rtimes D_8$ the theorem gives $\dim\tau_\rho=108$ and
$\tau_\rho|_{\tU_1}$ a sum of four irreducibles of dimension $27$,
permuted transitively by $D_8$.
(c) Nothing in the proof used irreducibility of $\rho$ beyond the last
sentence: for every finite-dimensional representation $\rho$ of
$D^\times$ trivial on $K^2$, $L_D\Theta_\rho$ is the character of a
genuine representation of $\tH$.
\end{remark}

\subsection{Level two}\label{W:subsec:targetB}
This subsection proves Target~B for the representations that are trivial on $K^3$ but not on $K^2$. It first describes these representations and the finite group through which the argument runs.

Let now $\rho$ be irreducible and trivial on $K^3$ but not on $K^2$.
Normalize $\theta_\rho\in\{\mathbf1,\mathrm{sgn}\}$ as in
\eqref{W:eq:normalized-central}; then $\rho|_{K^1}$ factors through
$\mathrm{Heis}:=Q_1/Q_3=K^1/K^3U_F^{(1)}$, the Heisenberg group of order
$27$, and is nontrivial on its center $Q_2/Q_3$.  The group
$D_8=D^\times/F^\times K^1$ permutes the two three-dimensional
irreducible characters $\chi_3,\chi_3'$ of $\mathrm{Heis}$, so
$\rho|_{K^1}=e(\chi_3+\chi_3')$, and $e=1$ for irreducible $\rho$: the
stabilizer of $\chi_3$ in $D_8$ is the cyclic group of order $4$
generated by the class of $1+i$, which acts trivially on the center
$Q_2/Q_3$ of $\mathrm{Heis}$ (while $j$ inverts it), and a character
stable under a cyclic group extends.  So $\dim\rho=6$; there are four such $\rho$ for each normalized central
character.  For these $\rho$ the criterion of
Proposition~\ref{W:prop:criterion} \emph{fails}: the restriction of
$F_\rho$ to the good classes of $Q_3/Q_7$ is a rational but not an
integral combination of the restrictions of the $\Ad(Q_1)$-orbit sums of
$\Irr(Q_3/Q_7)$ (script~13).  Nevertheless:

\begin{lemma}[The group $H=Q_3/Q_7$]\label{W:lem:Hstructure}
Let $H=Q_3/Q_7$, of order $3^6$.  Then $H'=[H,H]=Q_6/Q_7$, of order $3$;
$Z(H)=Q_4/Q_7$, of index $9$; $H$ has $243$ linear characters and $54$
nonlinear irreducible characters, all of degree $3$, namely the
irreducible characters nontrivial on $H'$, $27$ above each of the two
nontrivial characters $\omega$ of $H'$; and
\begin{equation}\label{W:eq:regH}
 3\sum_{\chi'(1)=3}\chi'=\Reg_H-\mathrm{Inf}_{H/H'}^{H}\Reg_{H/H'},\qquad
 \sum_{\chi'(1)=3,\ \chi'|_{H'}=3\omega}\chi'=81\,\omega\,\mathbf 1_{H'} .
\end{equation}
\end{lemma}
\begin{proof}
$[K^3,K^3]\subset K^6$ and $[K^3,K^4]\subset K^7$, so $H'\subset Q_6/Q_7$
and $Q_4/Q_7\subset Z(H)$.  For $x=1+\varpi aj+\cdots$, $y=1+\varpi bj+\cdots$
of depth three, $[x,y]\equiv1+\varpi^3(a\bar b-b\bar a)$ modulo $K^7$
(odd-odd bracket), and $a\bar b-b\bar a=2i\,\mathrm{Im}(a\bar b)$ is
nonzero modulo $\Fq_3$ for suitable $b$ whenever $a\ne0$; so
$H'=Q_6/Q_7\simeq\Fq_9/\Fq_3$, of order $3$, and no element of depth
exactly three is central, so $Z(H)=Q_4/Q_7$ and $[H:Z(H)]=|Q_3/Q_4|=9$.
Hence $|H/H'|=243$ linear characters; every nonlinear irreducible
character has degree $\le\sqrt{[H:Z(H)]}=3$, hence $3$; their number is
$(729-243)/9=54$; they are exactly the characters nontrivial on $H'$,
and by Schur's lemma each restricts to $H'$ as $3\omega$.  The first
identity in \eqref{W:eq:regH} is the decomposition of the regular
character, since the linear characters of $H$ are the inflations of
those of $H/H'$.  For the second, the $\omega$-isotypic part of $\C[H]$
under the central subgroup $H'$ is the $H$-module $\C[H]e_\omega$ with
character $243\,\omega\mathbf 1_{H'}$ (as $|H/H'|=243$), and it contains
each irreducible character above $\omega$ with multiplicity equal to its
degree $3$; dividing by $3$ gives the identity, and comparing degrees
($81\cdot3=243=27\cdot9$) gives $27$ characters above each $\omega$.
\end{proof}

\begin{theorem}[\Pctag\ Target B for level two]\label{W:thm:leveltwo}
Assume Theorem~\ref{W:thm:main-targetA}.  Let $\rho$ be an irreducible
representation of $D^\times$ trivial on $K^3$ with $\theta_\rho|_{\mu_3}=1$.  Then $L_D\Theta_\rho$, computed
with any compatible genuine central character, is the character of a
genuine irreducible representation $\tau_\rho$ of $\tH$ of dimension
$27\dim\rho$.  If $\rho|_{K^2}$ is not scalar, then after normalizing
$\theta_\rho\in\{\mathbf1,\mathrm{sgn}\}$,
\[
 \tau_\rho|_{\tU_1}=(\varepsilon_\rho\circ p)\otimes\bigl(\Ind_{\tU_3}^{\tU_1}(\lambda_0\otimes\chi_1'\circ p)\ \oplus\ \Ind_{\tU_3}^{\tU_1}(\lambda_0\otimes\chi_2'\circ p)\bigr),
\]
where $\chi_1',\chi_2'$ are three-dimensional irreducible characters of
$Q_3/Q_7$ (one from each of the two orbits described in the proof); both
induced representations are irreducible of dimension $81$, and they are
distinct.
\end{theorem}
\begin{proof}
By Theorem~\ref{W:thm:levelone} we may assume $\rho|_{K^2}$ nonscalar, and
by twisting, \eqref{W:eq:normalized-central}.  By
Theorem~\ref{W:thm:reduction-all} it suffices to show that
$\Phi_\rho=\Phi_1F_\rho$ is a virtual character of $\tU_1$.  The proof
has three steps, all by hand; the finite computations of
scripts~13--15 corroborate Steps~2 and~3 and are not used.

\emph{Step 1: a finite quotient.}  For $g=zt^3$ good in $U_3$ with
$t\in K^1$, the class of $t$ in $\mathrm{Heis}$ depends only on the
class of $g$ in $Q_3/Q_7$: for $g\in Q_6$ this is the linearity of the
inverse Artin--Schreier map (Lemma~\ref{W:lem:cube}(d)), and for
$g\in G_3$ it follows from the proof of Lemma~\ref{W:lem:cube}(f)
($(t(1+u))^3\in t^3Q_6$ for $u\in\Pfr^3$); script~13 verifies it on
the $219$ good classes of $Q_3/Q_7$ ($216$ of depth three, $3$ in
$Q_6/Q_7$), and verifies that the cube root is $\Ad(Q_1)$-equivariant.
Hence $F_\rho(\tilde zt^*)=\varepsilon_\rho(z)\Theta_\rho(t)$
(Lemma~\ref{W:lem:product}) is a function of the class of $g$ in
$Q_3/Q_7$ (times the character $\varepsilon_\rho$ of $U_1$).  Moreover
$\Phi_\rho$ is invariant under right multiplication by $\ell(K^7)$: for
$n\in K^7$ and $g$ good in $U_3$, $\ell(gn)=\ell(g)\ell(n)$ (for
$g\in U_6$ by (H1); for $g\in G_3$ by Lemma~\ref{W:lem:H1consequences}(d),
since $\ell(n)\in N_6$ as $\theta_0(\Nrd n)=1$), and $gn$ is good
(Lemma~\ref{W:lem:cube}(f)).  So $\Phi_\rho$ is a class function on the
finite group $\hat\Gamma:=\tU_1/\ell(K^7)\ker\tth$, of order $3^{10}$,
supported on the normal subgroup $\hat H:=\tU_3/\ell(K^7)\ker\tth$ of
index $27$, which is a central extension of $H:=Q_3/Q_7$ (order $3^6$)
by $\mu_3$.  (In the normalization $\theta_1=\mathbf1$ fixed above,
$\ker\tth=\ker\tth_1=\{(1,z)\}$ when $\theta_\rho=\mathbf1$; when
$\theta_\rho=\mathrm{sgn}$ we use $\ker\tth_1$ in place of $\ker\tth$ and
carry the character $\varepsilon_\rho\circ p$ of $\tU_1$ separately, as
in Remark~(b) below.)

\emph{Step 2: the restriction to $\tU_3$.}  By
Theorem~\ref{W:thm:main-targetA}(a), $\Phi_1|_{\tU_3}=\sum_{t\in\tU_1/\tU_3}\lambda_0^t=\lambda_0\cdot(D\circ p)$
with $D(v)=\sum_{h\in\Lambda}\zeta^{h(v)}$, a function on $V=Q_3/Q_6$
which is $27$ at $v=0$, $9\zeta^{-e_*(v)}$ on $\bar G_3$ and $0$
elsewhere (Proposition~\ref{W:prop:forces}).  Hence
\[
 \Phi_\rho|_{\tU_3}=(\varepsilon_\rho\circ p)\,\lambda_0\cdot\bigl((\Theta_\rho\circ r)\,D\bigr)\circ p ,
\]
where $(\Theta_\rho\circ r)D$ is a class function on $H$ supported on
the good classes.  We compute it by hand, using
Lemma~\ref{W:lem:Hstructure}, whose second identity reads, for $h\in H$,
\begin{equation}\label{W:eq:deg3sum}
 \sum_{\chi'(1)=3,\ \chi'|_{H'}=3\omega}\chi'(h)=81\,\omega(h)\,\mathbf 1_{H'}(h),\qquad
 \sum_{\chi'(1)=3}\chi'(h)=\begin{cases}162,&h=1,\\-81,&h\in H'\setminus1,\\0,&h\notin H'.\end{cases}
\end{equation}
On the other side, let $\chi_3$ be one of the two three-dimensional
characters of $\mathrm{Heis}$, with central character $\omega_3$ on
$Q_2/Q_3$; it vanishes off the center.  For a good class $g$ of depth
exactly three, $r(g)\in\mathrm{Heis}$ has depth one, so
$(\chi_3\circ r)(g)=0$.  For $g\in Q_6/Q_7=H'$ one has $D(g)=27$, and
$r$ restricted to $H'$ is the inverse of the Artin--Schreier bijection
$Q_2/Q_3\to Q_6/Q_7$ of Lemma~\ref{W:lem:cube}(d), an isomorphism of
groups of order $3$; so $(\chi_3\circ r)(g)=3\omega_3(r(g))$ and
$(\chi_3\circ r)D=81\,(\omega_3\circ r)\mathbf 1_{H'}$.  Comparing with
\eqref{W:eq:deg3sum},
\begin{equation}\label{W:eq:chi3D}
 (\chi_3\circ r)\cdot D=\sum_{\chi'(1)=3,\ \chi'|_{H'}=3\,\omega_3\circ r}\chi',\qquad
 \bigl((\chi_3+\chi_3')\circ r\bigr)\cdot D=\sum_{\chi'(1)=3}\chi' ,
\end{equation}
each of the $54$ characters occurring exactly once.  (Script~14
confirms this from the character table of $H$; it also returns, for the
orbit sums of linear characters of $\mathrm{Heis}$, the sums of $27$,
resp.\ $108$, linear characters of $H$, in accordance with
Theorem~\ref{W:thm:levelone}.)  Thus
$\Phi_\rho|_{\tU_3}=(\varepsilon_\rho\circ p)\sum_{\mu\in S}\mu$, where
$S$ is the set of the $54$ genuine irreducible characters
$\mu=\lambda_0\cdot\chi'\circ p$ of $\hat H$ of degree $3$ with central
character $\zeta$ on $\mu_3$: a multiplicity-free character of
dimension $162$.

\emph{Step 3: Clifford theory.}  A class function on $\hat\Gamma$
supported on the normal subgroup $\hat H$ is determined by its
restriction to $\hat H$, and for $\mu\in\Irr(\hat H)$ with
$\hat\Gamma$-orbit $O$ and stabilizer $\hat\Gamma_\mu$, the character of
$\Ind_{\hat H}^{\hat\Gamma}\mu$ is supported on $\hat H$ with restriction
$[\hat\Gamma_\mu:\hat H]\sum_{\mu'\in O}\mu'$.  By
Proposition~\ref{W:prop:free} below, $\Ad(Q_1)$ acts freely on the genuine
irreducible characters of $\hat H$, so $S$ is the union of two orbits
$O_1,O_2$ of size $27$ (one for each $\omega$, by \eqref{W:eq:chi3D} and
the invariance of the central character), and
\[
 \Phi_\rho=(\varepsilon_\rho\circ p)\bigl(\Theta_{\Ind\mu_1}+\Theta_{\Ind\mu_2}\bigr),\qquad \mu_i\in O_i ,
\]
with $\Ind_{\hat H}^{\hat\Gamma}\mu_i$ irreducible of dimension $81$ and
distinct.  (Had some orbit $O\subset S$ had size $|O|<27$, then for an
irreducible constituent $\chi$ of $\Ind\mu_O$ with
$\chi|_{\hat H}=e\sum_O\mu'$ one would have
$\langle\Phi_\rho,\chi\rangle=e|O|/27\notin\Z$, since
$e^2\le[\hat\Gamma_\mu:\hat H]=27/|O|$.)  This proves the displayed
decomposition, and $\dim\tau_\rho=162=27\cdot6$.
\end{proof}

\begin{proposition}[\Pctag\ Freeness]\label{W:prop:free}
$\Ad(\tU_1)$ acts freely on the set of genuine irreducible characters of
$\hat H=\tU_3/\ell(K^7)\ker\tth_1$: for every genuine irreducible $\mu$
and every $t\in U_1\setminus U_3$, $\mu\circ\Ad(t)\ne\mu$.
Equivalently, $\Ind_{\tU_3}^{\tU_1}\mu$ is irreducible for every genuine
irreducible representation $\mu$ of $\tU_3$ trivial on
$\ell(K^7)\ker\tth_1$.
\end{proposition}
\begin{proof}
Let $A:=Q_4/Q_7$, central in $H$, and $\hat A$ its preimage in $\hat H$.
$\hat A$ is central in $\hat H$: for $x\in K^3$, $y\in K^4$ one has
$[x,y]\in K^7$, so $[\tilde x,\tilde y]=\zeta^{a}\ell([x,y])$ with
$\zeta^a=\psi_6([\tilde x,\tilde y])=\omega(x,y)=1$ by
Lemma~\ref{W:lem:H1consequences}(c), and $\ell([x,y])\in\ell(K^7)$.  Hence
$\mu|_{\hat A}=\mu(1)\,\hat\nu$ for a genuine character $\hat\nu$ of
$\hat A$, and $\mu\circ\Ad(t)=\mu$ implies $\hat\nu\circ\Ad(t)=\hat\nu$.
Write $\hat\nu=\lambda_0|_{\hat A}\cdot(\alpha\circ p)$ with $\alpha$ a
character of $A$, and let $\omega_\mu:=\alpha|_{H'}$ (possibly trivial).
Since $\lambda_0(\ell(u))=1$ for $u\in U_6$ ($\ell(u)\in N_6$, as
$\psi_6(\ell(u))=\theta_1(z_u)=1$), the restriction of $\hat\nu$ to
$\ell(Q_6/Q_7)$ is $\omega_\mu$.  For $y\in A$,
\[
 \frac{\hat\nu\circ\Ad(t)}{\hat\nu}(\tilde y)=\hat\nu([\tilde t,\tilde y])
 =\frac{\lambda_0^t}{\lambda_0}(\tilde y)\cdot\omega_\mu([t,y])
 =c_t(y)\,\omega_\mu([t,y]),
\]
where $c_t:=\lambda_0^t/\lambda_0=\zeta^{h_t}$ is a character of $V$
with $h_t\in\Lambda$, and $t\mapsto h_t$ is a bijection
$Q_1/Q_3\to\Lambda$ with $h_t=0$ iff $t\in Q_3$ (because
$\sum_t\lambda_0^t=\Phi_1|_{\tU_3}=\sum_{h\in\Lambda}\lambda_h$,
Theorem~\ref{W:thm:main-targetA}(a), and the $27$ characters $\lambda_h$
are distinct).  Let $t\notin Q_3$.

If $t\in Q_2$, then $[t,y]\in K^7$ for all $y\in K^4$ (the leading
brackets $[\varpi a,\varpi^2c]=0$ and $[\varpi a,\varpi^2dj]=\varpi^3d(a-\bar a)j$
have depth $\ge7$), so the displayed ratio is $c_t|_A$, which is
nontrivial: no nonzero element of $\Lambda$ vanishes on the coordinates
$(C_{\mathrm{im}},B_{\mathrm{re}},B_{\mathrm{im}})$ of $Q_4/Q_6$
(Appendix~\ref{W:app:tables}).

If $t$ has depth exactly one with leading coefficient $a\in\Fq_9^\times$,
take $y\in Q_5$ with leading coefficient $b$ (so $y=1+\varpi^2bj+\cdots$).
Then $c_t(y)=\lambda_0([\tilde t,\tilde y])=\zeta^{\kappa_{1,5}(t,y)}=\zeta^{\beta\Tr(a\bar b)}$
with $\beta\ne0$ (Lemma~\ref{W:lem:H1consequences}(b)), and $[t,y]\in Q_6/Q_7\simeq\Fq_9/\Fq_3$
is the class of $\varpi^3(a\bar b-\bar ab)=2i\,\mathrm{Im}(a\bar b)\varpi^3$,
so $\omega_\mu([t,y])=\zeta^{2\gamma\,\mathrm{Im}(a\bar b)}$ for some
$\gamma\in\Fq_3$ depending on $\mu$.  As $b$ runs over $\Fq_9$ so does
$u=a\bar b$, and $u\mapsto2\beta\,\mathrm{Re}(u)+2\gamma\,\mathrm{Im}(u)$ is
a nonzero $\Fq_3$-linear functional.  So the ratio is nontrivial on
$Q_5$.

In both cases $\hat\nu\circ\Ad(t)\ne\hat\nu$, hence $\mu\circ\Ad(t)\ne\mu$.
\end{proof}

\begin{remark}[Certificate]
Script~15 computes $\lambda_0$ on $\hat H$ in the cocycle model, the
twisted action $\mu^t=\lambda_0c_t(\chi'\circ\Ad t)$ of the $27$ coset
representatives on the $54$ characters of $S$ and on the $243$ genuine
linear characters, verifies that no nontrivial $t$ fixes any of them,
and writes the orbit table to \texttt{level2\_orbit\_certificate.txt}:
two orbits of size $27$ on $S$, nine of size $27$ on the linear
characters (the nine induced representations of
Theorem~\ref{W:thm:levelone}).  For comparison, the untwisted action
$\chi'\mapsto\chi'\circ\Ad t$ on the $54$ characters of $H$ has six
orbits of size $9$: the $\lambda_0$-twist is what makes the action free.
These computations corroborate Proposition~\ref{W:prop:free} and are not
used in its proof.
\end{remark}

\begin{remark}
(a) The criterion of Proposition~\ref{W:prop:criterion} is sufficient but
not necessary: here $\tau_\rho|_{\tU_1}$ \emph{is} of the form
$\Ind(\lambda_0\otimes\sigma)$ with $\sigma=\chi_1'+\chi_2'$, but
$\Theta_\sigma$ is not $\Ad(U_1)$-invariant, and the product formula
$\Theta_{\Ind(\lambda_0\otimes\sigma)}=\Theta_\sigma\Theta_{\Ind\lambda_0}$
fails.
(b) The $\varepsilon_\rho$ bookkeeping: for $\theta_\rho=\mathrm{sgn}$,
$\varepsilon_\rho$ is a character of $U_1$ trivial on $K^1$, so twisting
by $\varepsilon_\rho\circ p$ preserves virtual characters of $\tU_1$ and
commutes with everything above.
(c) Label.  The theorem rests on Target~A through $\lambda_0$, $\Lambda$
and $e_*$ only (Theorem~\ref{W:thm:main-targetA}(a),
Proposition~\ref{W:prop:design}, and the table of $\Lambda$); Steps~2 and
3 are hand arguments.  Scripts~13--15 are corroboration.
\end{remark}

\subsection{Beyond level two}\label{W:subsec:cuberoot}
This subsection describes the pattern common to the two proofs above and isolates the statements that an extension to higher levels would need. It is an outline, as its label says.

\Otag\ The proofs of Theorems~\ref{W:thm:levelone} and~\ref{W:thm:leveltwo}
both run as follows.  For $m\ge1$ put $N_m:=\max(3m+1,6)$,
$\mathrm{Heis}_m:=Q_1/Q_{m+1}$, $H_m:=Q_3/Q_{N_m}$,
$\hat H_m:=\tU_3/\ell(K^{N_m})\ker\tth_1$ and
$\hat\Gamma_m:=\tU_1/\ell(K^{N_m})\ker\tth_1$ (for $m=1$, $\hat H_1=\tU_3/N_6=\tV$,
whose genuine characters are the $243$ characters $\lambda_h$,
$h\in V^\vee$; for $m=2$ these are the groups of
Theorem~\ref{W:thm:leveltwo}).  The cube root induces a well-defined
\emph{surjection} $r$ from the good classes of $H_m$ onto
$\mathrm{Heis}_m$, many-to-one ($219$ good classes of $H_2$ onto the
$27$ elements of $\mathrm{Heis}_2$), whose fibers are the sets
$\{t^3Q_{N_m}\}$ for $t$ in a fixed class of $Q_1/Q_{m+1}$.  For
$\rho$ trivial on $K^{m+1}$, $\Phi_\rho|_{\tU_3}=\lambda_0\cdot((\Theta_\rho\circ r)D)$,
and the two statements
\begin{enumerate}[label=(B\arabic*)]
\item $(\Theta_\rho\circ r)D$ is a character of $H_m$ (for $m\le2$: for
the $D_8$-invariant character $\Theta_\rho|_{K^1}=\sum m_\psi\psi$ of
$\mathrm{Heis}_1$ it is $\sum_\psi m_\psi\sum_{h\in\Lambda}\zeta^h\cdot\psi\circ\pi_3$,
and for $\chi_3+\chi_3'$ it is \eqref{W:eq:chi3D});
\item every genuine irreducible constituent $\mu$ of $\Phi_\rho|_{\hat H_m}$
has trivial inertia in $\hat\Gamma_m/\hat H_m\simeq Q_1/Q_3$, or more
generally the inertia indices and multiplicities satisfy the
divisibility conditions of Step~3 (for $m\le2$: freeness,
Proposition~\ref{W:prop:free} and the free orbits of
Proposition~\ref{W:prop:design}),
\end{enumerate}
give $\tau_\rho|_{\tU_1}=(\varepsilon_\rho\circ p)\otimes\bigoplus_i\Ind_{\tU_3}^{\tU_1}\mu_i$,
the $\mu_i$ running over one representative of each $Q_1$-orbit of
genuine constituents.  For $m=1$ the $243$ genuine characters of
$\hat H_1$ fall into nine free orbits of size $27$, of which $\Lambda$
is one; the nine induced representations are those of
Theorem~\ref{W:thm:levelone}.

A warning about the shape of (B1): it holds for $D_8$-invariant
characters of $\mathrm{Heis}_m$, which are the only ones Target~B needs,
but not constituent-wise.  For a single additive character $\chi$ of
$\Fq_9$ the function $\chi\circ r$ on the depth-three good classes is
$\chi(\pm\pi_3(g))$ with a sign depending on $g$ ($w^5=\Nm(w)w$), which
is not $\chi'\circ\pi_3$ for any character $\chi'$, and
$\Phi_1\cdot(\chi\circ r)$ is not a virtual character of $\tU_1$ for
$\chi\ne\mathbf1$: with $\chi(a+bi)=\zeta^a$ and the character $\zeta^{h}$
of $V$ with $h=(0,0,0,1,1)$, the table of $e_*$ gives
$\langle D\cdot(\chi\circ r),\zeta^h\rangle_V=\tfrac13$ (the $72$ good
vectors contribute exponents with frequencies $28,22,22$).  At level two
the constituent-wise identity \eqref{W:eq:chi3D} holds for $\chi_3$
because $\chi_3$ vanishes off the center of $\mathrm{Heis}_2$ and
$1+i$ acts trivially on that center.

For $m\ge3$ both (B1) and (B2) are open; Section~\ref{W:sec:future}
describes what a proof would need and why the mechanism of level two
does not continue unchanged.

\section{Target C: uniform statements}\label{W:sec:targetC}
This section records which parts of the argument are uniform in the
odd prime $p$ and the field $F$.  Throughout: $p$ is an odd prime,
$F/\Q_p$ is finite with $\mu_p\subset F$ (so $(p-1)\mid e:=e(F/\Q_p)$),
$q=|k_F|$, $\varpi$ is a uniformizer, $n=p$, $E_0/F$ is the unramified
quadratic extension, $D=E_0\oplus E_0j$ with $j^2=\varpi$ and
$jx=\sigma(x)j$, $K^m=1+\Pfr^m$, and
\begin{equation}\label{W:eq:dgen}
 d:=\frac{pe}{p-1},\qquad k_0:=\frac{2e}{p-1}=\frac{2d}{p},
\end{equation}
so that $U^{(k)}_{E_0}\subset E_0^{\times p}$ for $k>d$ and
$(U^{(i)}_{E_0},U^{(j)}_{E_0})_p=1$ for $i+j>d$ \cite[Ch.~IV]{FV02}; in the
running example $p=3$, $e=2$, $d=3$, $k_0=2$.  The cover is the
restriction along $\iota$ of the Kazhdan--Patterson cover of
$\mathrm{GL}_2(E_0)$ with cocycle $\sigma_{E_0}(g,h)(\det g,\det h)^{c}_{E_0}$,
$\sigma_{E_0}$ as in \eqref{W:eq:kubota} with the $p$-th power symbol,
and $c\equiv-1/4\pmod p$, which is the condition for scalars to be
central (the computation of Lemma~\ref{W:lem:model}(b) gives
$\beta(z,g)-\beta(g,z)=(1+4c)(z,\det g)$); for $p=3$ this is $c=2$.  We
write $h_E$ for exponents of the $p$-th power symbol of $E_0$ and
$(a,b)_E=h_E(a,b)$ additively.

\begin{lemma}[Depth facts]\label{W:lem:depth-gen}
\begin{enumerate}[label=(\alph*)]
\item For $k>k_0$ the $p$-th power map $K^k\to K^{k+2e}$ is bijective.
\item If $w\in D^\times$ and $w^p\in F^\times$, then $w\in F^\times$.
\item Every element of $K^{2d}$ is of the form $zv^p$ with
$z\in U_F^{(1)}$ and $v\in K^{k_0}$; $z$ is unique modulo $F^{\times p}$
and $z\equiv\Nrd(X)^{(p+1)/2}\pmod{F^{\times p}}$.
\end{enumerate}
\end{lemma}
\begin{proof}
(a) As in Lemma~\ref{W:lem:cube}(a): $v_D(p)=2e$, and for $u\in\Pfr^k$ the
terms $\binom pi u^i$ ($2\le i\le p-1$) and $u^p$ lie in $\Pfr^{k+2e+1}$
when $k>k_0$.  (b) If $w\notin F$ then $F(w)$ is a quadratic field in
which $X^p-c$ ($c=w^p$) has a root; since $\mu_p\subset F$, $X^p-c$ is
either split or irreducible over $F$, and $p>2$ excludes the latter, so
$w\in\mu_pw_0\subset F$ for a root $w_0\in F$.  (c) Every $X\in K^{2d}\setminus F^\times$
lies in a unique maximal subfield $E$ (and $X\in F^\times\cap K^{2d}=U_F^{(d)}$
is trivially of the required form).  Recall the standard fact that for a
local field $L\supset\mu_p$ with $e_L=e(L/\Q_p)$, the $p$-th power map
$U^{(k)}_L\to U^{(k+e_L)}_L$ is bijective for $k>e_L/(p-1)$ \cite[Ch.~I]{FV02}.

\emph{$E$ ramified.}  Then $e_E=2e$, $k_E=k_F$, $E\cap K^{2d}=U_E^{(2d)}$
(as $v_D=v_E$), and $U_F^{(d)}\subset U_E^{(2d)}$ induces an isomorphism
$U_F^{(d)}/U_F^{(d+1)}\to U_E^{(2d)}/U_E^{(2d+1)}$ of residue groups
($\varpi_E^2=\varpi\cdot$unit).  Choose $z\in U_F^{(d)}$ with
$z^{-1}X\in U_E^{(2d+1)}$.  Since $k_0+1>2e/(p-1)=e_E/(p-1)$ and
$k_0+1+2e=2d+1$, the $p$-th power map $U_E^{(k_0+1)}\to U_E^{(2d+1)}$ is
bijective, so $z^{-1}X=v^p$ with $v\in U_E^{(k_0+1)}\subset K^{k_0+1}$.

\emph{$E$ unramified.}  Then $e_E=e$, $k_E=\Fq_{q^2}$, $E\cap K^{2d}=U_E^{(d)}$
(as $v_D=2v_E$), and $p$-th powers of $U_E^{(k)}$ fill $U_E^{(k+e)}$ for
$k>e/(p-1)$; the boundary level $k=e/(p-1)=k_0/2$ needs an argument.
Write $\zeta_p=1+\varpi^{e/(p-1)}a_0+\cdots$ with $a_0\in\Ocal_F^\times$
(possible since $v_F(\zeta_p-1)=e/(p-1)$), so that
$p/\varpi^{e}\equiv-\bar a_0^{\,p-1}$ in $k_F$ (from $\zeta_p^p=1$).  For
$u=1+\varpi^{e/(p-1)}a$, $u^p\equiv1+\varpi^d(a^p+(p/\varpi^e)a)$ modulo
$U_E^{(d+1)}$, so the induced map $k_E\to k_E$ on the graded pieces
$U_E^{(e/(p-1))}/U_E^{(e/(p-1)+1)}\to U_E^{(d)}/U_E^{(d+1)}$ is
$\bar a\mapsto\bar a^p-\bar a_0^{\,p-1}\bar a=\bar a_0^{\,p}\,\wp(\bar a/\bar a_0)$
with $\wp(x)=x^p-x$.  Its image is $\bar a_0^{\,p}\ker(\Tr_{k_E/\Fq_p})$, of
index $p$ in $k_E$.  The image of $U_F^{(d)}$ in $U_E^{(d)}/U_E^{(d+1)}$
is $k_F$, and $k_F\not\subset\bar a_0^{\,p}\ker(\Tr_{k_E/\Fq_p})$: for
$u\in k_F$, $\Tr_{k_E/\Fq_p}(u/\bar a_0^{\,p})=2\Tr_{k_F/\Fq_p}(u/\bar a_0^{\,p})$,
which is nonzero for suitable $u$ because $p$ is odd and the trace
$k_F\to\Fq_p$ is onto.  Hence $k_F$ and the image together span $k_E$:
choose $z\in U_F^{(d)}$ and $v_0\in U_E^{(e/(p-1))}$ with
$z^{-1}v_0^{-p}X\in U_E^{(d+1)}=(U_E^{(e/(p-1)+1)})^p$, say $=v_1^p$, and
put $v=v_0v_1\in U_E^{(e/(p-1))}\subset K^{k_0}$ (as $v_D=2v_E$).

In both cases $z\in U_F^{(d)}\subset U_F^{(1)}$ and $v\in E\cap K^{k_0}$.
Uniqueness and the formula for $z$: $\Nrd X=z^2\Nrd(v)^p$, and squaring
is bijective on $F^\times/F^{\times p}$, with inverse $y\mapsto y^{(p+1)/2}$.
\end{proof}

\begin{theorem}[The $p$-th power identity]\label{W:thm:pthpower}
Let $w=x+yj\in K^1$.  Then $(w,0)^p=(w^p,0)$ in the cover.
\end{theorem}
\begin{proof}
If $y=0$ the diagonal formula gives exponent
$\sum_{k=1}^{p-1}\bigl(h_E(x^k,\sigma x)+2ck\,h_E(\Nm x,\Nm x)\bigr)=\tfrac{p(p-1)}2h_E(x,\sigma x)=0$.
Let $y\ne0$, $T=\Trd w$, $n=\Nrd w$, and define $S_k\in F$ by
$S_0=0$, $S_1=1$, $S_{k+1}=TS_k-nS_{k-1}$, so that
$w^k=S_kw-nS_{k-1}$ by Cayley--Hamilton, i.e.\ $w^k=x_k+y_kj$ with
$y_k=S_ky$ and $x_k=S_kx-nS_{k-1}$.  By Lemma~\ref{W:lem:depth-gen}(b),
$S_k\ne0$ for $1\le k\le p$ ($S_k=0$ would put $w^k\in E_0$ hence,
as $w\notin E_0$, $w^k\in F$; for $k<p$ this forces $w\in F$ because
$K^1F^\times/F^\times$ is pro-$p$, and for $k=p$ by (b)).  In the cocycle
$\beta(w^k,w)$ one has $x(\iota w^k)=\sigma(y)S_k$, $x(\iota w)=\sigma(y)$,
$x(\iota w^{k+1})=\sigma(y)S_{k+1}$, $\det\iota w^k=n^k$, so with
$\tau=1/S_k$ and $R=S_{k+1}/S_k$ (all in $F$), by the computation
leading to \eqref{W:eq:betaK6},
\[
 \beta(w^k,w)=-(R,\tau)_E-(R,n^k)_E+c\,(n^k,n)_E
 =(S_{k+1},S_k)_E-k(S_{k+1},n)_E+k(S_k,n)_E ,
\]
using $(S_k,S_k)_E=(n,n)_E=0$.  Put $a_k:=(S_{k+1},S_k)_E$ and
$b_k:=(S_k,n)_E$ (so $a_1=b_1=0$).  Summing over $k=1,\dots,p-1$ and
collecting the $b$'s,
\begin{equation}\label{W:eq:pthsum}
 \sum_{k=1}^{p-1}\beta(w^k,w)=\sum_{k=1}^{p-1}a_k+\sum_{k=2}^{p}b_k\pmod p .
\end{equation}
The Cassini identity $S_{k+1}S_{k-1}-S_k^2=-n^{k-1}$ (induction on
$k$) and the Steinberg relation for $u=S_{k+1}S_{k-1}/S_k^2$,
$1-u=n^{k-1}/S_k^2$ (both nonzero for $2\le k\le p-1$) give, after
expanding bimultiplicatively,
\[
 2(a_k-a_{k-1})=(k-1)\,(b_{k+1}-2b_k+b_{k-1})\qquad(2\le k\le p-1).
\]
Hence $a_k=\tfrac12\sum_{j=2}^k(j-1)\Delta^2b_j$, where
$\Delta^2b_j:=b_{j+1}-2b_j+b_{j-1}$, and interchanging the sums,
\[
 \sum_{k=1}^{p-1}a_k=\sum_{j=2}^{p-1}c_j\,\Delta^2b_j,\qquad
 c_j:=\tfrac12(j-1)(p-j)\quad(\text{computations in }\Z[\tfrac12]/p).
\]
Extend $c_j$ to all integers $j$ by the same quadratic formula; then
$c_{j-1}-2c_j+c_{j+1}=-1$ for every $j$, $c_1=c_p=0$ and
$c_0=c_{p+1}=-p/2\equiv0$.  Expanding $\Delta^2b_j$ and reindexing,
\[
 \sum_{j=2}^{p-1}c_j\Delta^2b_j
 =\sum_{j=1}^{p}b_j\,(c_{j-1}-2c_j+c_{j+1})
  -\bigl(c_0b_1+c_1b_2+c_pb_{p-1}+c_{p+1}b_p\bigr)
 \equiv-\sum_{j=1}^{p}b_j ,
\]
the four correction terms being those of the full double sum
$\sum_{j=1}^p b_j(c_{j-1}-2c_j+c_{j+1})$ that do not occur in
$\sum_{j=2}^{p-1}c_j\Delta^2b_j$, each of which vanishes modulo $p$.
Since $b_1=0$, \eqref{W:eq:pthsum} vanishes.  For $p=3$ this is the computation of
Proposition~\ref{W:prop:cubeidentity}.
\end{proof}

\begin{theorem}[(H1) for every odd $p$]\label{W:thm:H1-gen}
Define the star section on $U_{2d}=F^\times K^{2d}$ by
$\ell(zv^p):=(z,0)(v,0)^p$ with $v\in K^{k_0}$.  It is well defined,
$\Ad(D^\times)$-equivariant, and a group homomorphism; explicitly
$\ell(X)=(X,\kappa(X))$ with
\[
 \kappa(X)=-\tfrac{p+1}2\,(y_X,\Nrd X)_E\quad(y_X\ne0),\qquad \kappa(X)=0\quad(y_X=0).
\]
Consequently, for $\theta_0$ with $\theta_0^p=1$ and
$\tth(\zeta,z)=\zeta\theta_0(z)$, the analogue $\psi_{2d}$ of $\psi_6$
($\zeta\tth(\tilde z)$ at $\zeta\tilde zv^*$) is a character of
$\tU_{2d}$, given on $\tilde K^{2d}$ by
\[
 (X,e)\longmapsto\zeta^{\,e-\kappa(X)}\,\theta_0(\Nrd X)^{(p+1)/2},
\]
since the value at the star lift $\ell(X)$ is $\theta_0(z)$ with
$z\equiv\Nrd(X)^{(p+1)/2}$ (Lemma~\ref{W:lem:depth-gen}(c)).  For $p=3$
the exponent $2\equiv-1$ recovers \eqref{W:eq:psi6formula}.
\end{theorem}
\begin{proof}
Theorem~\ref{W:thm:pthpower} gives $\ell(zv^p)=(zv^p,\beta(z,v^p))$, and
$\beta(z,g)=(z,x(g))_E+2c(z,\det g)_E$ with $\det g=\Nrd(v)^p$, so
$\kappa(X)=(z,y_X/z)_E=(z,y_X)_E=\tfrac{p+1}2(\Nrd X,y_X)_E$ by
Lemma~\ref{W:lem:depth-gen}(c) (and $0$ if $y_X=0$, as
$(z,x_X)_E=\tfrac{p+1}2(x_X\sigma x_X,x_X)_E=0$ for $x_X\in U^{(d)}$).
Well-definedness is the uniqueness in Lemma~\ref{W:lem:depth-gen}(c)
together with $h_E(F^{\times p},\cdot)=0$.  For multiplicativity we
repeat the proof of Proposition~\ref{W:prop:H1} with $m:=-\tfrac{p+1}2$,
so $2m\equiv-1$, for $X=a+bj$, $X'=a'+b'j\in K^{2d}$: now
$a,a'\in U^{(d)}$, $b,b'\in\varpi^d\Ocal$, $n,n'\in U_F^{(d)}$, all mutual
symbols among $a,\sigma a,a',\sigma a',n,n'$ vanish ($2d>d$), the twist
$c(n,n')_E$ vanishes, $n\equiv\Nm a$ modulo $U^{(2d+1)}\subset E_0^{\times p}$,
and \eqref{W:eq:level0}--\eqref{W:eq:lambda} hold with $U^{(3)}$ replaced by
$U^{(d)}$ (as $1+d>d$ and $\mu_{q^2-1}\subset E_0^{\times p}$), with
$\lambda(a)=(\varpi,\Nm a)_E=2(\varpi,a)_E$.  In the main case
\eqref{W:eq:H1identity} becomes
$(R,\tau)_E=m(\tau,n')_E-(m+1)(R,n)_E-m(R,n')_E=-v(\tau)(\varpi,a')_E-v(R)(\varpi,a)_E+v(R)(\varpi,a')_E$,
which is exactly the value of $(R,\tau)_E$ computed from
$R=a'(1+\tau_2)$; in the cases $b=0$, $b'=0$, $b=b'=0$ and $R=0$
the required identities are $(2m+1)v(b')(\varpi,a)_E=0$,
$2mv(b)(\varpi,a')_E=-v(b)(\varpi,a')_E$, $0=0$ and
$(2m+1)v(b)[(\varpi,a)_E+(\varpi,a')_E]=0$ respectively, all true since
$2m+1\equiv0$.  The passage from $K^{2d}$ to $U_{2d}$ and the formula
for $\psi_{2d}$ are as in Corollary~\ref{W:cor:starK6}(d).
\end{proof}

\begin{remark}[What is and is not uniform]\label{W:rem:uniform}
\begin{enumerate}[label=(\roman*)]
\item Uniform in $p$ and $F$ (for $n=p$): the toral theorem
(Section~\ref{W:sec:general-toral}); the positivity criterion; the cocycle
model and the $p$-th power identity; the depth facts; hypothesis (H1)
in the form of Theorem~\ref{W:thm:H1-gen}; the product formula and the
criterion of Section~\ref{W:sec:plan} (with $3$ replaced by $p$, $K^3$ by
$K^{p}$ and $\lambda_0$ by the character produced by the analogue of
Theorem~\ref{W:thm:H1equiv}, once available).
\item \Otag\ Expected to be uniform but not written: the Brauer
reduction.  $D^\times/F^\times K^1$ is now the group of order $2(q+1)$
generated by $k_{E_0}^\times/k_F^\times$ and $j$; the centralizer lemma
uses only that $p$ is odd; the $\ell$-elementary subgroups for
$\ell\ne p$ live in tori or in the preimage of a $2$-subgroup, and
the $2$-group lemma needs the same analysis as in
Theorem~\ref{W:thm:reduction-all}, with $SD_{16}$ replaced by the
appropriate extension.
\item \Ctag\ Not uniform: the finite design.  The good elements of
depth exactly $p$ (the analogue of $G_3$, with leading coefficients
$w^{1+(p-1)(q+1)/2}$, which need not exhaust $k_{E_0}^\times$ when
$\gcd(1+(p-1)(q+1)/2,\,q^2-1)>1$) and the quotient $V=Q_p/Q_{2d}$ depend on
$(p,e,q)$, and the uniqueness of the defect function was a computation.
A conceptual replacement for Section~\ref{W:sec:finite} is the main open
problem of Target~C; the next test cases are $F=\Q_3(\zeta_3,\sqrt3)$
($e=4$, $d=6$) and $F=\Q_5(\zeta_5)$ ($e=4$, $d=5$).
\item For composite odd $n$ and for $p\nmid n$ nothing new is needed
at tame places (Chapter~\ref{ch:III}); for $n=p^a m$ with $a\ge2$ the depth
bookkeeping changes ($n$-th powers of $K^k$ land in $K^{k+2ae}$) and we
have not examined it.
\end{enumerate}
\end{remark}

\subsection{Proof checkpoints}
This subsection lists what a positive solution for a general row must supply, and notes which items are now available.

A positive solution for a row must still supply (i) an explicit
cocycle/central-character normalization (now available, uniformly in
$p$), (ii) a finite-level factorization or an induced construction
(now: $\Ind$ from $\tU_3$, Theorems~\ref{W:thm:levelone}, \ref{W:thm:leveltwo}),
(iii) exact full-group character coefficients (now: the Fourier
identities of Proposition~\ref{W:prop:forces}), and (iv) a uniform
all-$\rho$ argument (open beyond level two).

\section{Future directions}\label{W:sec:future}
We stop at level two.  This section records, as precisely as we can,
what a proof of Target~B for all levels would need, what the first open
case looks like, and which parts of Target~C are within reach.  Nothing
here is claimed.

\subsection{Why level two is a natural stopping point}
Levels are unbounded (Section~\ref{W:subsec:divalg}), so no finite list of
levels proves Target~B; what is needed is an inductive mechanism, and the
proofs of Sections~\ref{W:sec:plan} do not contain one.  They rely on
features special to $m\le2$: the group $H_2=Q_3/Q_7$ has class two with
derived subgroup of order $3$, so its nonlinear characters are
controlled by the regular character (Lemma~\ref{W:lem:Hstructure});
$\Theta_\rho\circ r$ vanishes off a central subgroup, so (B1) is a
one-line identity; and the central subgroup $\hat A=\widehat{Q_4/Q_7}$
is large enough for central characters to detect every nontrivial
$t\in Q_1/Q_3$ (Proposition~\ref{W:prop:free}).  At $m=3$ all three fail
at once: $H_3=Q_3/Q_{10}$ has order $3^{11}$, class $3$, and its
largest obvious abelian normal subgroup is $Q_5/Q_{10}$ (note that
$Q_4/Q_{10}$ is \emph{not} abelian: $[1+\varpi^2i,1+\varpi^2j]\equiv1+2\varpi^4ij$
modulo $K^{10}$, of depth $9$); $\Theta_\rho\circ r$ for $\rho$ of level
three is supported on the preimages of the noncentral elements of
$\mathrm{Heis}_3=Q_1/Q_4$ as well; and the natural central subgroup
$\hat A_3=\widehat{Q_7/Q_{10}}$ lies inside $\tU_6$, where
$\lambda_0$ is trivial and only the brackets $[t,\cdot]$ remain to
separate characters.  The level-two theorem is therefore a complete
result about a finite family of rows, and an honest statement of the
present state is: Target~B is proved for every row of level $\le2$,
and open beyond.

\subsection{The two statements to prove at each level}
For $\rho$ of level $m$, with the notation of
Section~\ref{W:subsec:cuberoot}, Target~B for $\rho$ follows from (B1)
and (B2); conversely, if $\Phi_\rho$ is a virtual character of $\tU_1$
then its restriction to $\hat H_m$ is one, so (B1) is necessary.  The
two should be tested separately: a failure of (B1) (a non-integral
multiplicity in $(\Theta_\rho\circ r)D$) would be a genuine obstruction
to Target~B, exhibiting a row with a non-integral coefficient
$a_{\rho\tau}$; a failure of universal freeness in (B2) would not, since
the divisibility conditions of Step~3 of Theorem~\ref{W:thm:leveltwo}
($e|O|\in27\Z$ for each constituent) can hold with nontrivial inertia.

\subsection{Level three}
The first open case is $m=3$, $\rho$ trivial on $K^4$.  A finite
verification of (B1) and (B2) is feasible but needs a different method
from scripts~13--15: $H_3$ is too large for Burnside's algorithm, and
one should instead describe $\Irr(H_3)$ by Clifford theory relative to
the abelian normal subgroup $Q_5/Q_{10}$ (abelian because
$[K^5,K^5]\subset K^{10}$), determine the action of $H_3$ on its dual,
and compute the inner products of $(\Theta_\rho\circ r)D$ with the
resulting characters, where now $\Theta_\rho$ is a $D_8$-invariant
character of $\mathrm{Heis}_3$ of order $3^5$ and the good classes of
$H_3$ number $216\cdot3^4+\cdots$.  For (B2) one needs the twisted
action of $Q_1/Q_3$ on the genuine constituents, i.e.\ $\lambda_0$ on
$\hat H_3$ and the brackets; the certificate format of script~15 would
carry over.  Whatever the outcome, level three decides whether the
mechanism of level two (multiplicity-free restriction, free action,
induced irreducibles of dimension $27\dim\mu$) is the general shape of
$\tau_\rho|_{\tU_1}$ or a low-level accident.

\subsection{Toward an inductive mechanism}
Two ideas seem to us the most promising.  First, an orbit-theoretic
description.  Howe's Kirillov theory \cite{How77a} describes the characters
of the congruence subgroups $K^m$ (for $m$ large relative to $p$) by
coadjoint orbits on a Lie lattice; the genuine constituents of
$\Phi_\rho|_{\tU_3}$ should correspond to orbits shifted by the fixed
``metaplectic'' datum that $\lambda_0$ represents, and the twisted
action of $Q_1/Q_3$ on them to an affine coadjoint action.  The question
is whether the nonvanishing of $\kappa_{1,5}$, which forced regularity
at level two, continues to force it at all depths or whether deeper
stabilizers can absorb it.  Second, an inductive form of freeness: for
$m\ge3$ the subgroup $\hat A_m=\widehat{Q_{3m-2}/Q_{3m+1}}$ is central in
$\hat H_m$ with no hypothesis (for $x\in K^3$, $y\in K^{3m-2}\subset U_6$,
$[\tilde x,\tilde y]=\ell([x,y])\in\ell(K^{3m+1})$ because $\ell$ is an
$\Ad$-equivariant homomorphism on $U_6$), so a genuine irreducible
$\mu$ of $\hat H_m$ has a central character $\lambda_0\alpha_\mu$ on
$\hat A_m$ and $\mu\circ\Ad(t)=\mu$ forces $\alpha_\mu([t,A_m])=1$;
freeness for the constituents of exact level $m$ would follow from
nondegeneracy of the brackets $Q_1/Q_2\times Q_{3m-2}/Q_{3m-1}\to Q_{3m-1}/Q_{3m}$
and $Q_1/Q_2\times Q_{3m-1}/Q_{3m}\to Q_{3m}/Q_{3m+1}$ on the relevant
characters, and for the others from level $m-1$; but the characters
$\alpha_\mu$ that actually occur must first be identified, which is
(B1) again.

\subsection{Target C}
Of the uniform statements, the Brauer reduction (Theorem~\ref{W:thm:reduction-all})
should go through for every odd $p$ and $F$ with the same proof, the
$2$-group step being replaced by an analysis of the preimage of the
group $D^\times/F^\times K^1$ of order $2(q+1)$; the finite design of
Section~\ref{W:sec:finite} is the part with no uniform argument, and the
two test cases named in Remark~\ref{W:rem:uniform} ($e=4$ with $p=3$, and
$p=5$) would show whether the uniqueness of the defect function is
special to $(p,e,q)=(3,2,3)$.

\subsection{The duality}
Remark~\ref{W:rem:howe-duality} reformulates Target~B as the existence of
a representation $\Omega_\theta$ of $\tH\times H$ with the
multiplicity-one property of a Howe duality, whose level-$m$
truncations have explicit characters.  A construction of
$\Omega_\theta$ that does not go through $L_D$, for instance on a space
of functions on which $\tH$ and $H$ act through the cube-root map and
the star section, would prove Target~B and explain it; the induced
models of Theorems~\ref{W:thm:levelone} and~\ref{W:thm:leveltwo} are what
such a construction must produce at levels $\le2$.

\section{Interface with the larger Jacquet--Langlands program}
The companion problem is local.  Theorem~\ref{W:thm:main-targetA} and
Corollary~\ref{W:cor:levelzero} remove, for the rows they cover, the
``$L^2$-unitarity-to-integrality'' gap recorded in Chapter~\ref{ch:III},
at the wild place $\pfr_3$ of the running example; in particular the
character relation (D2) of Chapter~\ref{ch:I} now has, at $\pfr_3$, a proved
instance for $\rho=\mathbf1$ and the level-zero rows.  Even full
success on Target~B would not on its own prove the additional global
comparison of ordinary $D^\times$ and the covered group: that comparison
remains conditional on matching functions at exceptional split places,
complex-place normalization, a global orbital-coefficient identity,
and the support/germ hypotheses stated in Chapter~\ref{ch:IV}.  Conversely, the
global comparison could be established in a restricted test range
without first settling every local irreducible-character row.  These
projects should be developed in parallel but not conflated.

\begin{remark}[Target B as a Howe duality]\label{W:rem:howe-duality}
Fix a character $\theta$ of $F^\times$ trivial on $\mu_3$ and a
compatible $\tth$, and for $m\ge0$ let $\Irr_\theta^{\le m}$ be the
finite set of irreducible representations of $H$ with central character
$\theta$ and level $\le m$, i.e.\ the irreducible representations of
the finite group $H/K^{m+1}$ with that central character.  By
Lemma~\ref{W:lem:product} (whose proof does not use
\eqref{W:eq:normalized-central}), $L_D\Theta_\rho=\Theta_{\tau_1}\cdot F_\rho$
with $F_\rho(\zeta\tilde zt^*)=\tth\tth_1^{-1}(\tilde z)\Theta_\rho(t)$ on
star lifts of good elements, and $L_D\Theta_\rho=0$ elsewhere.  Hence
the finite sum
\[
 K_\theta^{(m)}(x,g):=\sum_{\rho\in\Irr_\theta^{\le m}}L_D\Theta_\rho(x)\,\overline{\Theta_\rho(g)}
 =\Theta_{\tau_1}(x)\,\tth\tth_1^{-1}(\tilde z)\,\Delta^{(m)}_\theta(t,g),\qquad x=\zeta\tilde zt^*,\ g\in H,
\]
where $\Delta^{(m)}_\theta(t,g):=\sum_{\rho\in\Irr_\theta^{\le m}}\Theta_\rho(t)\overline{\Theta_\rho(g)}$
is the $\theta$-isotypic class delta of the finite group $H/K^{m+1}$
(it vanishes unless $g\equiv zt$ modulo $K^{m+1}$ for some $z\in F^\times$
up to conjugacy, and is then $\theta(z)$ times the order of the
centralizer in $H/F^\times K^{m+1}$), is a well-defined class function on
$\tH\times H/K^{m+1}$, supported where $g$ is, up to the center and
modulo $K^{m+1}$, the cube root of $p(x)$, and explicit because the
trivial row is.  (The sum over all $\rho$ does not converge pointwise,
$K_\theta(1,1)=27\sum_\rho(\dim\rho)^2$; only the truncations, or the
pairing of the full sum against test functions in $g$, make sense.)
Since the $\Theta_\rho$ are orthonormal, the $H$-Fourier coefficients of
$K^{(m)}_\theta$ are the $L_D\Theta_\rho$, $\rho\in\Irr_\theta^{\le m}$;
so $K_\theta^{(m)}$ is the character of a finite-dimensional
representation $\Omega^{(m)}_\theta$ of $\tH\times H/K^{m+1}$ if and only
if every $L_D\Theta_\rho$ with $\rho$ of level $\le m$ is a character,
and in that case $\Omega^{(m)}_\theta=\bigoplus_{\rho}\tau_\rho\otimes\rho^\vee$
with each $\tau_\rho$ irreducible and the $\tau_\rho$ pairwise distinct
(Proposition~\ref{W:prop:positive} and unitarity).  Target~B for the
central character $\theta$ is therefore equivalent to: for every $m$,
$K^{(m)}_\theta$ is a character.  When it holds, the formal direct sum
$\Omega_\theta=\varinjlim\Omega^{(m)}_\theta$ has the multiplicity-one
Howe-duality property of Section~\ref{W:subsec:howe}, with the
$K^{(m)}_\theta$ in the role of the character of the Weil representation
restricted to the pair.  Theorems~\ref{W:thm:levelone} and~\ref{W:thm:leveltwo}
say that $K^{(2)}_\theta$ is a character and describe
$\Omega^{(2)}_\theta$ on $\tU_1\times H/K^3$ as a sum of induced
representations.  This is a reformulation, not a result: we know no
construction of $\Omega^{(m)}_\theta$ that does not go through $L_D$.
\end{remark}
\endgroup
